% !TEX program = lualatex
% ---------------------------------------------------------------------------------------------
%  Readback note 1: From a one-sided bound for primes to a zero-free half-plane
%
%  Build:  python3 tools/rb.py note <slug>
%
%  One unit of a Lean artifact, rewritten as ordinary mathematics. The prose is hand-written. Lean names,
%  line ranges and counts come from data.tex; quoted Lean statements are copied from the artifact by the build.
%
%  Rules of the text: standard terminology; every symbol is explained where it first appears and listed again
%  in the notation at the end; no letter has two meanings; a proof is complete or says what it takes for granted;
%  Lean names appear only in the LEAN lines and in the concordance.
%
%  The signature rule (the build enforces it): a function, relation or predicate that the artifact defines is
%  written only through its notation macro, which the build makes from the spec (`notation`) with one slot for
%  each explicit argument of the Lean definition. No argument is fixed and dropped, also not "for brevity" and
%  not for one section. Every argument and every hypothesis of a stated declaration stands inside its
%  statement, or in a Setting that the statement cites by number; only implicit and instance arguments (the
%  ambient objects) may stand in the conventions.
% ---------------------------------------------------------------------------------------------
\documentclass[11pt,a4paper]{article}
\ifdefined\pdfvariable\pdfvariable suppressoptionalinfo \numexpr 1+2+4+8\relax\fi
\usepackage[a4paper,hmargin=35mm,top=22mm,bottom=27mm,footskip=12mm,nohead]{geometry}

\usepackage{fontspec}
\directlua{
  luaotfload.add_fallback("gfallback", {"STIXTwoMath-Regular.otf:mode=node;", "DejaVuSans.ttf:mode=node;"})
  luaotfload.add_fallback("gmonofallback", {"DejaVuSansMono.ttf:mode=node;", "STIXTwoMath-Regular.otf:mode=node;"})
}
\setmainfont{NewCM10}[Extension=.otf,UprightFont=*-Book,ItalicFont=*-BookItalic,
  BoldFont=*-Bold,BoldItalicFont=*-BoldItalic,SmallCapsFont=*-Book,SmallCapsFeatures={Letters=SmallCaps},
  RawFeature={fallback=gfallback}]
\setsansfont{NewCMSans10}[Extension=.otf,UprightFont=*-Book,ItalicFont=*-BookOblique,
  BoldFont=*-Bold,BoldItalicFont=*-BoldOblique,RawFeature={fallback=gfallback}]
\setmonofont{JetBrainsMono}[Extension=.otf,UprightFont=*-Regular,ItalicFont=*-Italic,
  BoldFont=*-SemiBold,BoldItalicFont=*-SemiBoldItalic,Scale=MatchLowercase,
  RawFeature={-calt;-liga;fallback=gmonofallback}]
\usepackage{amsmath}
\usepackage{amsthm}
\usepackage[math-style=TeX]{unicode-math}
\setmathfont{NewCMMath-Book.otf}
\emergencystretch=2.5em\relax

\usepackage{xcolor}
\usepackage{tikz}
\usetikzlibrary{calc,arrows.meta,positioning,backgrounds}
\usepackage[protrusion=true,expansion=true]{microtype}
\usepackage{enumitem}
\usepackage{needspace}
\usepackage{booktabs}
\usepackage{array}
\usepackage{longtable}
\usepackage{fancyhdr}
\usepackage{lastpage}
\usepackage{multicol}
\usepackage[hidelinks,pdfusetitle]{hyperref}

\definecolor{ink}{HTML}{0F172A}
\definecolor{inksoft}{HTML}{475569}
\definecolor{inkfaint}{HTML}{94A3B8}
\definecolor{rulegrey}{HTML}{CBD5E1}
\definecolor{arith}{HTML}{C2410C}
\definecolor{est}{HTML}{7C3AED}
\definecolor{ana}{HTML}{0891B2}
\definecolor{ok}{HTML}{047857}
\definecolor{leancol}{HTML}{3730A3}

\title{From a one-sided bound for primes to a zero-free half-plane}
\author{grh-compact}
\date{}

% ------------------------------------------------------------------------------------- macros
\newcommand{\leanfont}{\ttfamily\fontsize{8.6}{10.5}\selectfont\color{leancol}}
\newcommand{\lean}[1]{\csname LN@#1\endcsname}
\newcommand{\fillin}[1]{{\color{arith}\itshape[#1]}}  % a place of the template that is not written yet; the build reports every one that is left
\newcommand{\leanquote}[1]{\csname LQ@#1\endcsname}
% generated by rb.py note; do not edit
\newcommand{\ConcDecls}{44}
\newcommand{\Pin}{72fead10}
\newcommand{\PinText}{that of commit {\leanfont 72fead10}}
\newcommand{\ArtifactName}{grh-compact}
\expandafter\newcommand\csname LN@chiK\endcsname{\href{https://github.com/lsernau/grh-lean/blob/72fead102f663d5c7a407964a13dc0cf9d977207/grh-compact/src/Lib/General/NumberTheory/QuadField.lean\#L20-L20}{{\leanfont chiK}}}
\expandafter\newcommand\csname LN@quadTwist\endcsname{\href{https://github.com/lsernau/grh-lean/blob/72fead102f663d5c7a407964a13dc0cf9d977207/grh-compact/src/Lib/General/NumberTheory/QuadField.lean\#L26-L27}{{\leanfont quadTwist}}}
\expandafter\newcommand\csname LN@quadTwistapply\endcsname{\href{https://github.com/lsernau/grh-lean/blob/72fead102f663d5c7a407964a13dc0cf9d977207/grh-compact/src/Lib/Analytic/Landau/PolarLandau.lean\#L20-L22}{{\leanfont quadTwist\_apply}}}
\expandafter\newcommand\csname LN@entireL\endcsname{\href{https://github.com/lsernau/grh-lean/blob/72fead102f663d5c7a407964a13dc0cf9d977207/grh-compact/src/Lib/General/NumberTheory/QuadField.lean\#L31-L31}{{\leanfont entireL}}}
\expandafter\newcommand\csname LN@chowlaSelbergConst\endcsname{\href{https://github.com/lsernau/grh-lean/blob/72fead102f663d5c7a407964a13dc0cf9d977207/grh-compact/src/Lib/General/NumberTheory/QuadField.lean\#L35-L35}{{\leanfont chowlaSelbergConst}}}
\expandafter\newcommand\csname LN@lambdaClass\endcsname{\href{https://github.com/lsernau/grh-lean/blob/72fead102f663d5c7a407964a13dc0cf9d977207/grh-compact/src/New/Endgame/PrimeError.lean\#L25-L25}{{\leanfont lambdaClass}}}
\expandafter\newcommand\csname LN@lambdaSum\endcsname{\href{https://github.com/lsernau/grh-lean/blob/72fead102f663d5c7a407964a13dc0cf9d977207/grh-compact/src/New/Endgame/PrimeError.lean\#L29-L29}{{\leanfont lambdaSum}}}
\expandafter\newcommand\csname LN@logDerivMean\endcsname{\href{https://github.com/lsernau/grh-lean/blob/72fead102f663d5c7a407964a13dc0cf9d977207/grh-compact/src/New/Endgame/PrimeError.lean\#L33-L34}{{\leanfont logDerivMean}}}
\expandafter\newcommand\csname LN@primeError\endcsname{\href{https://github.com/lsernau/grh-lean/blob/72fead102f663d5c7a407964a13dc0cf9d977207/grh-compact/src/New/Endgame/PrimeError.lean\#L38-L39}{{\leanfont primeError}}}
\expandafter\newcommand\csname LQ@primeError\endcsname{\leanquoteline{0}{def primeError [NeZero M]}\leanquoteline{4}{x := lambdaSum d M 0 x - lambdaSum d M 1 x - x.log / φ M - (logDerivMean d M 1).re}}
\expandafter\newcommand\csname LN@IsPolar\endcsname{\href{https://github.com/lsernau/grh-lean/blob/72fead102f663d5c7a407964a13dc0cf9d977207/grh-compact/src/Lib/Analytic/Landau/PolarLandau.lean\#L34-L36}{{\leanfont IsPolar}}}
\expandafter\newcommand\csname LQ@IsPolar\endcsname{\leanquoteline{0}{def IsPolar (E : ℝ → ℝ) (r : ℂ → ℂ) := Measurable E ∧ ∃ C Φ, Meromorphic Φ ∧ (∀ v ≥ 0, |E v| ≤ C *}\leanquoteline{2}{v.exp) ∧ (∀ s, 2 < re s → Φ s = ∫ v in Ioi 0, E v * ((1 - s) * v).exp) ∧ AnalyticOnNhd ℂ (fun s ↦}\leanquoteline{2}{(s - 1) * Φ s) \{s | 1 ≤ s.re\} ∧ ∀ ρ, 0 < ρ.re → (𝓝[≠] ρ).Tendsto (fun z ↦ (z - ρ) * Φ z) (𝓝 (r ρ))}}
\expandafter\newcommand\csname LN@entireLnezero\endcsname{\href{https://github.com/lsernau/grh-lean/blob/72fead102f663d5c7a407964a13dc0cf9d977207/grh-compact/src/Lib/Analytic/Landau/PolarLandau.lean\#L55-L56}{{\leanfont entireL\_ne\_zero}}}
\expandafter\newcommand\csname LN@LSeriesHasSummulvonMangoldt\endcsname{\href{https://github.com/lsernau/grh-lean/blob/72fead102f663d5c7a407964a13dc0cf9d977207/grh-compact/src/Lib/Analytic/Landau/PolarLandau.lean\#L58-L62}{{\leanfont LSeriesHasSum\_mul\_vonMangoldt}}}
\expandafter\newcommand\csname LN@tendstosubmullogDeriventireL\endcsname{\href{https://github.com/lsernau/grh-lean/blob/72fead102f663d5c7a407964a13dc0cf9d977207/grh-compact/src/Lib/Analytic/Landau/PolarLandau.lean\#L64-L73}{{\leanfont tendsto\_sub\_mul\_logDeriv\_entireL}}}
\expandafter\newcommand\csname LN@quadTwistneone\endcsname{\href{https://github.com/lsernau/grh-lean/blob/72fead102f663d5c7a407964a13dc0cf9d977207/grh-compact/src/Lib/Analytic/Landau/PolarLandau.lean\#L26-L29}{{\leanfont quadTwist\_ne\_one}}}
\expandafter\newcommand\csname LN@landau\endcsname{\href{https://github.com/lsernau/grh-lean/blob/72fead102f663d5c7a407964a13dc0cf9d977207/grh-compact/src/Lib/Analytic/Landau/PolarLandau.lean\#L88-L136}{{\leanfont IsPolar.\allowbreak{}landau}}}
\expandafter\newcommand\csname LQ@landau\endcsname{\leanquoteline{0}{lemma IsPolar.landau \{r\} (h : IsPolar E r) \{β : ℝ\} (hβ : 0 < β) (hr : ∀ x, β < x → 0 ≤ (r x).re)}\leanquoteline{4}{(hw : ∀ᶠ v in atTop, E v ≤ ((β - 1) * v).exp) \{ρ\} (hρ : β < ρ.re) : r ρ = 0}}
\expandafter\newcommand\csname LN@existsentireLquadTwistnezero\endcsname{\href{https://github.com/lsernau/grh-lean/blob/72fead102f663d5c7a407964a13dc0cf9d977207/grh-compact/src/Lib/Analytic/QuadraticTwist/ZeroFreeTwist.lean\#L193-L206}{{\leanfont exists\_entireL\_quadTwist\_ne\_zero}}}
\expandafter\newcommand\csname LN@laplacepow\endcsname{\href{https://github.com/lsernau/grh-lean/blob/72fead102f663d5c7a407964a13dc0cf9d977207/grh-compact/src/Lib/Analytic/Landau/PolarLandau.lean\#L47-L50}{{\leanfont laplace\_pow}}}
\expandafter\newcommand\csname LN@integrableOnlaplaceofabsle\endcsname{\href{https://github.com/lsernau/grh-lean/blob/72fead102f663d5c7a407964a13dc0cf9d977207/grh-compact/src/Lib/Analytic/Landau/PolarLandau.lean\#L40-L43}{{\leanfont integrableOn\_laplace\_of\_abs\_le}}}
\expandafter\newcommand\csname LN@summabletaylorofremovable\endcsname{\href{https://github.com/lsernau/grh-lean/blob/72fead102f663d5c7a407964a13dc0cf9d977207/grh-compact/src/Lib/Analytic/Landau/PolarLandau.lean\#L78-L82}{{\leanfont summable\_taylor\_of\_removable}}}
\expandafter\newcommand\csname LN@lambdaClassnonnegandabsle\endcsname{\href{https://github.com/lsernau/grh-lean/blob/72fead102f663d5c7a407964a13dc0cf9d977207/grh-compact/src/New/Endgame/PrimeError.lean\#L41-L43}{{\leanfont lambdaClass\_nonneg\_and\_abs\_le}}}
\expandafter\newcommand\csname LN@abslambdaSumexple\endcsname{\href{https://github.com/lsernau/grh-lean/blob/72fead102f663d5c7a407964a13dc0cf9d977207/grh-compact/src/New/Endgame/PrimeError.lean\#L45-L48}{{\leanfont abs\_lambdaSum\_exp\_le}}}
\expandafter\newcommand\csname LN@primeErroreqsum\endcsname{\href{https://github.com/lsernau/grh-lean/blob/72fead102f663d5c7a407964a13dc0cf9d977207/grh-compact/src/New/Endgame/PrimeError.lean\#L60-L64}{{\leanfont primeError\_eq\_sum}}}
\expandafter\newcommand\csname LN@relogDerivMeanone\endcsname{\href{https://github.com/lsernau/grh-lean/blob/72fead102f663d5c7a407964a13dc0cf9d977207/grh-compact/src/New/Endgame/PrimeError.lean\#L52-L54}{{\leanfont re\_logDerivMean\_one}}}
\expandafter\newcommand\csname LN@isPolarprimeError\endcsname{\href{https://github.com/lsernau/grh-lean/blob/72fead102f663d5c7a407964a13dc0cf9d977207/grh-compact/src/New/Endgame/PrimeError.lean\#L69-L111}{{\leanfont isPolar\_primeError}}}
\expandafter\newcommand\csname LQ@isPolarprimeError\endcsname{\leanquoteline{0}{lemma isPolar\_primeError (hd : d \% 8 = 7) (h : M.Coprime d) : IsPolar (primeError d M ·.exp) fun s ↦}\leanquoteline{4}{(∑ ξ, (analyticOrderNatAt (@entireL M \_ ξ) s - analyticOrderNatAt (entireL (quadTwist ξ d))}\leanquoteline{2}{s : ℂ)) / φ M / (s * (1 - s))}}
\expandafter\newcommand\csname LN@entireLnezeroofprimeErrorle\endcsname{\href{https://github.com/lsernau/grh-lean/blob/72fead102f663d5c7a407964a13dc0cf9d977207/grh-compact/src/New/Endgame/OneSidedCriterion.lean\#L28-L36}{{\leanfont entireL\_ne\_zero\_of\_primeError\_le}}}
\expandafter\newcommand\csname LQ@entireLnezeroofprimeErrorle\endcsname{\leanquoteline{0}{lemma entireL\_ne\_zero\_of\_primeError\_le \{M\} [NeZero M] \{a : ZMod M\} \{β\} (hβ : 1 / 2 < β)}\leanquoteline{2}{(ha : IsUnit a) (h : ∀ d : Primes, d = a → (d.1 \% 8 = 7 ∧ M.Coprime d) ∧ ∀ᶠ x in atTop, primeError}\leanquoteline{2}{d M x ≤ x \textasciicircum{} (β - 1)) ρ (hρ : β < ρ.re) (χ : DirichletCharacter ℂ M) : entireL χ ρ ≠ 0}}
\expandafter\newcommand\csname LN@eventuallymomentle\endcsname{\href{https://github.com/lsernau/grh-lean/blob/72fead102f663d5c7a407964a13dc0cf9d977207/grh-compact/src/New/Endgame/OneSidedCriterion.lean\#L41-L67}{{\leanfont IsPolar.\allowbreak{}eventually\_moment\_le}}}
\expandafter\newcommand\csname LN@existswindowsubmulle\endcsname{\href{https://github.com/lsernau/grh-lean/blob/72fead102f663d5c7a407964a13dc0cf9d977207/grh-compact/src/New/Endgame/OneSidedCriterion.lean\#L73-L99}{{\leanfont IsPolar.\allowbreak{}exists\_window\_sub\_mul\_le}}}
\expandafter\newcommand\csname LN@mainHeight\endcsname{\href{https://github.com/lsernau/grh-lean/blob/72fead102f663d5c7a407964a13dc0cf9d977207/grh-compact/src/New/Endgame/ShortAndEnough.lean\#L32-L32}{{\leanfont Data.\allowbreak{}mainHeight}}}
\expandafter\newcommand\csname LN@mainAsk\endcsname{\href{https://github.com/lsernau/grh-lean/blob/72fead102f663d5c7a407964a13dc0cf9d977207/grh-compact/src/New/Endgame/ShortAndEnough.lean\#L36-L37}{{\leanfont Data.\allowbreak{}mainAsk}}}
\expandafter\newcommand\csname LN@remainder\endcsname{\href{https://github.com/lsernau/grh-lean/blob/72fead102f663d5c7a407964a13dc0cf9d977207/grh-compact/src/New/Endgame/ShortAndEnough.lean\#L42-L42}{{\leanfont Data.\allowbreak{}remainder}}}
\expandafter\newcommand\csname LN@remaindereq\endcsname{\href{https://github.com/lsernau/grh-lean/blob/72fead102f663d5c7a407964a13dc0cf9d977207/grh-compact/src/New/Endgame/ShortAndEnough.lean\#L86-L112}{{\leanfont Data.\allowbreak{}remainder\_eq}}}
\expandafter\newcommand\csname LQ@remaindereq\endcsname{\leanquoteline{0}{lemma Data.remainder\_eq : D.remainder = D.a * (\#D.thick *}\leanquoteline{2}{(primeError P.d P.N₀ D.x - primeError P.d P.N₀ D.x₁) + \#D.tops * primeError P.d 1 D.x₁)}}
\expandafter\newcommand\csname LN@heightPerRowle\endcsname{\href{https://github.com/lsernau/grh-lean/blob/72fead102f663d5c7a407964a13dc0cf9d977207/grh-compact/src/New/Endgame/ShortAndEnough.lean\#L48-L70}{{\leanfont Data.\allowbreak{}heightPerRow\_le}}}
\expandafter\newcommand\csname LN@absmainHeightaddremainderle\endcsname{\href{https://github.com/lsernau/grh-lean/blob/72fead102f663d5c7a407964a13dc0cf9d977207/grh-compact/src/New/Endgame/ShortAndEnough.lean\#L74-L82}{{\leanfont Data.\allowbreak{}abs\_mainHeight\_add\_remainder\_le}}}
\expandafter\newcommand\csname LN@scaledInclusion\endcsname{\href{https://github.com/lsernau/grh-lean/blob/72fead102f663d5c7a407964a13dc0cf9d977207/grh-compact/src/New/Endgame/ShortAndEnough.lean\#L116-L116}{{\leanfont Data.\allowbreak{}scaledInclusion}}}
\expandafter\newcommand\csname LN@projectedPoincare\endcsname{\href{https://github.com/lsernau/grh-lean/blob/72fead102f663d5c7a407964a13dc0cf9d977207/grh-compact/src/New/Endgame/ShortAndEnough.lean\#L120-L121}{{\leanfont Data.\allowbreak{}projectedPoincare}}}
\expandafter\newcommand\csname LN@integralRamified\endcsname{\href{https://github.com/lsernau/grh-lean/blob/72fead102f663d5c7a407964a13dc0cf9d977207/grh-compact/src/New/Endgame/ShortAndEnough.lean\#L125-L126}{{\leanfont Data.\allowbreak{}integralRamified}}}
\expandafter\newcommand\csname LN@IndexLe\endcsname{\href{https://github.com/lsernau/grh-lean/blob/72fead102f663d5c7a407964a13dc0cf9d977207/grh-compact/src/New/Endgame/ShortAndEnough.lean\#L130-L132}{{\leanfont Data.\allowbreak{}IndexLe}}}
\expandafter\newcommand\csname LN@existsratlogeqandpadicnormeq\endcsname{\href{https://github.com/lsernau/grh-lean/blob/72fead102f663d5c7a407964a13dc0cf9d977207/grh-compact/src/New/Endgame/ShortAndEnough.lean\#L21-L26}{{\leanfont exists\_rat\_log\_eq\_and\_padic\_norm\_eq}}}
\expandafter\newcommand\csname LN@shortforms\endcsname{\href{https://github.com/lsernau/grh-lean/blob/72fead102f663d5c7a407964a13dc0cf9d977207/grh-compact/src/New/Endgame/ShortAndEnough.lean\#L138-L168}{{\leanfont Data.\allowbreak{}short\_forms}}}
\expandafter\newcommand\csname LN@enoughforms\endcsname{\href{https://github.com/lsernau/grh-lean/blob/72fead102f663d5c7a407964a13dc0cf9d977207/grh-compact/src/New/Endgame/ShortAndEnough.lean\#L173-L178}{{\leanfont Data.\allowbreak{}enough\_forms}}}
\expandafter\newcommand\csname LN@heightLevolumeratiosq\endcsname{\href{https://github.com/lsernau/grh-lean/blob/72fead102f663d5c7a407964a13dc0cf9d977207/grh-compact/src/New/Endgame/LongForms.lean\#L19-L32}{{\leanfont Data.\allowbreak{}heightLe\_volume\_ratio\_sq}}}
\expandafter\newcommand\csname LN@heightLevolumerows\endcsname{\href{https://github.com/lsernau/grh-lean/blob/72fead102f663d5c7a407964a13dc0cf9d977207/grh-compact/src/New/Endgame/LongForms.lean\#L39-L58}{{\leanfont Data.\allowbreak{}heightLe\_volume\_rows}}}
\expandafter\newcommand\csname LN@integralformslong\endcsname{\href{https://github.com/lsernau/grh-lean/blob/72fead102f663d5c7a407964a13dc0cf9d977207/grh-compact/src/New/Endgame/LongForms.lean\#L63-L71}{{\leanfont Data.\allowbreak{}integral\_forms\_long}}}
\expandafter\newcommand\csname LN@existsfamilycovolumeSqle\endcsname{\href{https://github.com/lsernau/grh-lean/blob/72fead102f663d5c7a407964a13dc0cf9d977207/grh-compact/src/New/Endgame/LongForms.lean\#L78-L92}{{\leanfont Data.\allowbreak{}exists\_family\_covolumeSq\_le}}}
\newcommand{\ConcordanceRows}{\crow{chiK}{\href{https://github.com/lsernau/grh-lean/blob/72fead102f663d5c7a407964a13dc0cf9d977207/grh-compact/src/Lib/General/NumberTheory/QuadField.lean\#L20-L20}{{\leanfont chiK}}}{QuadField.lean}{20--20}
\crow{quadTwist}{\href{https://github.com/lsernau/grh-lean/blob/72fead102f663d5c7a407964a13dc0cf9d977207/grh-compact/src/Lib/General/NumberTheory/QuadField.lean\#L26-L27}{{\leanfont quadTwist}}}{QuadField.lean}{26--27}
\crow{quadTwistapply}{\href{https://github.com/lsernau/grh-lean/blob/72fead102f663d5c7a407964a13dc0cf9d977207/grh-compact/src/Lib/Analytic/Landau/PolarLandau.lean\#L20-L22}{{\leanfont quadTwist\_apply}}}{PolarLandau.lean}{20--22}
\crow{entireL}{\href{https://github.com/lsernau/grh-lean/blob/72fead102f663d5c7a407964a13dc0cf9d977207/grh-compact/src/Lib/General/NumberTheory/QuadField.lean\#L31-L31}{{\leanfont entireL}}}{QuadField.lean}{31--31}
\crow{chowlaSelbergConst}{\href{https://github.com/lsernau/grh-lean/blob/72fead102f663d5c7a407964a13dc0cf9d977207/grh-compact/src/Lib/General/NumberTheory/QuadField.lean\#L35-L35}{{\leanfont chowlaSelbergConst}}}{QuadField.lean}{35--35}
\crow{lambdaClass}{\href{https://github.com/lsernau/grh-lean/blob/72fead102f663d5c7a407964a13dc0cf9d977207/grh-compact/src/New/Endgame/PrimeError.lean\#L25-L25}{{\leanfont lambdaClass}}}{PrimeError.lean}{25--25}
\crow{lambdaSum}{\href{https://github.com/lsernau/grh-lean/blob/72fead102f663d5c7a407964a13dc0cf9d977207/grh-compact/src/New/Endgame/PrimeError.lean\#L29-L29}{{\leanfont lambdaSum}}}{PrimeError.lean}{29--29}
\crow{logDerivMean}{\href{https://github.com/lsernau/grh-lean/blob/72fead102f663d5c7a407964a13dc0cf9d977207/grh-compact/src/New/Endgame/PrimeError.lean\#L33-L34}{{\leanfont logDerivMean}}}{PrimeError.lean}{33--34}
\crow{primeError}{\href{https://github.com/lsernau/grh-lean/blob/72fead102f663d5c7a407964a13dc0cf9d977207/grh-compact/src/New/Endgame/PrimeError.lean\#L38-L39}{{\leanfont primeError}}}{PrimeError.lean}{38--39}
\crow{IsPolar}{\href{https://github.com/lsernau/grh-lean/blob/72fead102f663d5c7a407964a13dc0cf9d977207/grh-compact/src/Lib/Analytic/Landau/PolarLandau.lean\#L34-L36}{{\leanfont IsPolar}}}{PolarLandau.lean}{34--36}
\crow{entireLnezero}{\href{https://github.com/lsernau/grh-lean/blob/72fead102f663d5c7a407964a13dc0cf9d977207/grh-compact/src/Lib/Analytic/Landau/PolarLandau.lean\#L55-L56}{{\leanfont entireL\_ne\_zero}}}{PolarLandau.lean}{55--56}
\crow{LSeriesHasSummulvonMangoldt}{\href{https://github.com/lsernau/grh-lean/blob/72fead102f663d5c7a407964a13dc0cf9d977207/grh-compact/src/Lib/Analytic/Landau/PolarLandau.lean\#L58-L62}{{\leanfont LSeriesHasSum\_mul\_vonMangoldt}}}{PolarLandau.lean}{58--62}
\crow{tendstosubmullogDeriventireL}{\href{https://github.com/lsernau/grh-lean/blob/72fead102f663d5c7a407964a13dc0cf9d977207/grh-compact/src/Lib/Analytic/Landau/PolarLandau.lean\#L64-L73}{{\leanfont tendsto\_sub\_mul\_logDeriv\_entireL}}}{PolarLandau.lean}{64--73}
\crow{quadTwistneone}{\href{https://github.com/lsernau/grh-lean/blob/72fead102f663d5c7a407964a13dc0cf9d977207/grh-compact/src/Lib/Analytic/Landau/PolarLandau.lean\#L26-L29}{{\leanfont quadTwist\_ne\_one}}}{PolarLandau.lean}{26--29}
\crow{landau}{\href{https://github.com/lsernau/grh-lean/blob/72fead102f663d5c7a407964a13dc0cf9d977207/grh-compact/src/Lib/Analytic/Landau/PolarLandau.lean\#L88-L136}{{\leanfont IsPolar.\allowbreak{}landau}}}{PolarLandau.lean}{88--136}
\crow{existsentireLquadTwistnezero}{\href{https://github.com/lsernau/grh-lean/blob/72fead102f663d5c7a407964a13dc0cf9d977207/grh-compact/src/Lib/Analytic/QuadraticTwist/ZeroFreeTwist.lean\#L193-L206}{{\leanfont exists\_entireL\_quadTwist\_ne\_zero}}}{ZeroFreeTwist.lean}{193--206}
\crow{laplacepow}{\href{https://github.com/lsernau/grh-lean/blob/72fead102f663d5c7a407964a13dc0cf9d977207/grh-compact/src/Lib/Analytic/Landau/PolarLandau.lean\#L47-L50}{{\leanfont laplace\_pow}}}{PolarLandau.lean}{47--50}
\crow{integrableOnlaplaceofabsle}{\href{https://github.com/lsernau/grh-lean/blob/72fead102f663d5c7a407964a13dc0cf9d977207/grh-compact/src/Lib/Analytic/Landau/PolarLandau.lean\#L40-L43}{{\leanfont integrableOn\_laplace\_of\_abs\_le}}}{PolarLandau.lean}{40--43}
\crow{summabletaylorofremovable}{\href{https://github.com/lsernau/grh-lean/blob/72fead102f663d5c7a407964a13dc0cf9d977207/grh-compact/src/Lib/Analytic/Landau/PolarLandau.lean\#L78-L82}{{\leanfont summable\_taylor\_of\_removable}}}{PolarLandau.lean}{78--82}
\crow{lambdaClassnonnegandabsle}{\href{https://github.com/lsernau/grh-lean/blob/72fead102f663d5c7a407964a13dc0cf9d977207/grh-compact/src/New/Endgame/PrimeError.lean\#L41-L43}{{\leanfont lambdaClass\_nonneg\_and\_abs\_le}}}{PrimeError.lean}{41--43}
\crow{abslambdaSumexple}{\href{https://github.com/lsernau/grh-lean/blob/72fead102f663d5c7a407964a13dc0cf9d977207/grh-compact/src/New/Endgame/PrimeError.lean\#L45-L48}{{\leanfont abs\_lambdaSum\_exp\_le}}}{PrimeError.lean}{45--48}
\crow{primeErroreqsum}{\href{https://github.com/lsernau/grh-lean/blob/72fead102f663d5c7a407964a13dc0cf9d977207/grh-compact/src/New/Endgame/PrimeError.lean\#L60-L64}{{\leanfont primeError\_eq\_sum}}}{PrimeError.lean}{60--64}
\crow{relogDerivMeanone}{\href{https://github.com/lsernau/grh-lean/blob/72fead102f663d5c7a407964a13dc0cf9d977207/grh-compact/src/New/Endgame/PrimeError.lean\#L52-L54}{{\leanfont re\_logDerivMean\_one}}}{PrimeError.lean}{52--54}
\crow{isPolarprimeError}{\href{https://github.com/lsernau/grh-lean/blob/72fead102f663d5c7a407964a13dc0cf9d977207/grh-compact/src/New/Endgame/PrimeError.lean\#L69-L111}{{\leanfont isPolar\_primeError}}}{PrimeError.lean}{69--111}
\crow{entireLnezeroofprimeErrorle}{\href{https://github.com/lsernau/grh-lean/blob/72fead102f663d5c7a407964a13dc0cf9d977207/grh-compact/src/New/Endgame/OneSidedCriterion.lean\#L28-L36}{{\leanfont entireL\_ne\_zero\_of\_primeError\_le}}}{OneSidedCriterion.lean}{28--36}
\crow{eventuallymomentle}{\href{https://github.com/lsernau/grh-lean/blob/72fead102f663d5c7a407964a13dc0cf9d977207/grh-compact/src/New/Endgame/OneSidedCriterion.lean\#L41-L67}{{\leanfont IsPolar.\allowbreak{}eventually\_moment\_le}}}{OneSidedCriterion.lean}{41--67}
\crow{existswindowsubmulle}{\href{https://github.com/lsernau/grh-lean/blob/72fead102f663d5c7a407964a13dc0cf9d977207/grh-compact/src/New/Endgame/OneSidedCriterion.lean\#L73-L99}{{\leanfont IsPolar.\allowbreak{}exists\_window\_sub\_mul\_le}}}{OneSidedCriterion.lean}{73--99}
\crow{mainHeight}{\href{https://github.com/lsernau/grh-lean/blob/72fead102f663d5c7a407964a13dc0cf9d977207/grh-compact/src/New/Endgame/ShortAndEnough.lean\#L32-L32}{{\leanfont Data.\allowbreak{}mainHeight}}}{ShortAndEnough.lean}{32--32}
\crow{mainAsk}{\href{https://github.com/lsernau/grh-lean/blob/72fead102f663d5c7a407964a13dc0cf9d977207/grh-compact/src/New/Endgame/ShortAndEnough.lean\#L36-L37}{{\leanfont Data.\allowbreak{}mainAsk}}}{ShortAndEnough.lean}{36--37}
\crow{remainder}{\href{https://github.com/lsernau/grh-lean/blob/72fead102f663d5c7a407964a13dc0cf9d977207/grh-compact/src/New/Endgame/ShortAndEnough.lean\#L42-L42}{{\leanfont Data.\allowbreak{}remainder}}}{ShortAndEnough.lean}{42--42}
\crow{remaindereq}{\href{https://github.com/lsernau/grh-lean/blob/72fead102f663d5c7a407964a13dc0cf9d977207/grh-compact/src/New/Endgame/ShortAndEnough.lean\#L86-L112}{{\leanfont Data.\allowbreak{}remainder\_eq}}}{ShortAndEnough.lean}{86--112}
\crow{heightPerRowle}{\href{https://github.com/lsernau/grh-lean/blob/72fead102f663d5c7a407964a13dc0cf9d977207/grh-compact/src/New/Endgame/ShortAndEnough.lean\#L48-L70}{{\leanfont Data.\allowbreak{}heightPerRow\_le}}}{ShortAndEnough.lean}{48--70}
\crow{absmainHeightaddremainderle}{\href{https://github.com/lsernau/grh-lean/blob/72fead102f663d5c7a407964a13dc0cf9d977207/grh-compact/src/New/Endgame/ShortAndEnough.lean\#L74-L82}{{\leanfont Data.\allowbreak{}abs\_mainHeight\_add\_remainder\_le}}}{ShortAndEnough.lean}{74--82}
\crow{scaledInclusion}{\href{https://github.com/lsernau/grh-lean/blob/72fead102f663d5c7a407964a13dc0cf9d977207/grh-compact/src/New/Endgame/ShortAndEnough.lean\#L116-L116}{{\leanfont Data.\allowbreak{}scaledInclusion}}}{ShortAndEnough.lean}{116--116}
\crow{projectedPoincare}{\href{https://github.com/lsernau/grh-lean/blob/72fead102f663d5c7a407964a13dc0cf9d977207/grh-compact/src/New/Endgame/ShortAndEnough.lean\#L120-L121}{{\leanfont Data.\allowbreak{}projectedPoincare}}}{ShortAndEnough.lean}{120--121}
\crow{integralRamified}{\href{https://github.com/lsernau/grh-lean/blob/72fead102f663d5c7a407964a13dc0cf9d977207/grh-compact/src/New/Endgame/ShortAndEnough.lean\#L125-L126}{{\leanfont Data.\allowbreak{}integralRamified}}}{ShortAndEnough.lean}{125--126}
\crow{IndexLe}{\href{https://github.com/lsernau/grh-lean/blob/72fead102f663d5c7a407964a13dc0cf9d977207/grh-compact/src/New/Endgame/ShortAndEnough.lean\#L130-L132}{{\leanfont Data.\allowbreak{}IndexLe}}}{ShortAndEnough.lean}{130--132}
\crow{existsratlogeqandpadicnormeq}{\href{https://github.com/lsernau/grh-lean/blob/72fead102f663d5c7a407964a13dc0cf9d977207/grh-compact/src/New/Endgame/ShortAndEnough.lean\#L21-L26}{{\leanfont exists\_rat\_log\_eq\_and\_padic\_norm\_eq}}}{ShortAndEnough.lean}{21--26}
\crow{shortforms}{\href{https://github.com/lsernau/grh-lean/blob/72fead102f663d5c7a407964a13dc0cf9d977207/grh-compact/src/New/Endgame/ShortAndEnough.lean\#L138-L168}{{\leanfont Data.\allowbreak{}short\_forms}}}{ShortAndEnough.lean}{138--168}
\crow{enoughforms}{\href{https://github.com/lsernau/grh-lean/blob/72fead102f663d5c7a407964a13dc0cf9d977207/grh-compact/src/New/Endgame/ShortAndEnough.lean\#L173-L178}{{\leanfont Data.\allowbreak{}enough\_forms}}}{ShortAndEnough.lean}{173--178}
\crow{heightLevolumeratiosq}{\href{https://github.com/lsernau/grh-lean/blob/72fead102f663d5c7a407964a13dc0cf9d977207/grh-compact/src/New/Endgame/LongForms.lean\#L19-L32}{{\leanfont Data.\allowbreak{}heightLe\_volume\_ratio\_sq}}}{LongForms.lean}{19--32}
\crow{heightLevolumerows}{\href{https://github.com/lsernau/grh-lean/blob/72fead102f663d5c7a407964a13dc0cf9d977207/grh-compact/src/New/Endgame/LongForms.lean\#L39-L58}{{\leanfont Data.\allowbreak{}heightLe\_volume\_rows}}}{LongForms.lean}{39--58}
\crow{integralformslong}{\href{https://github.com/lsernau/grh-lean/blob/72fead102f663d5c7a407964a13dc0cf9d977207/grh-compact/src/New/Endgame/LongForms.lean\#L63-L71}{{\leanfont Data.\allowbreak{}integral\_forms\_long}}}{LongForms.lean}{63--71}
\crow{existsfamilycovolumeSqle}{\href{https://github.com/lsernau/grh-lean/blob/72fead102f663d5c7a407964a13dc0cf9d977207/grh-compact/src/New/Endgame/LongForms.lean\#L78-L92}{{\leanfont Data.\allowbreak{}exists\_family\_covolumeSq\_le}}}{LongForms.lean}{78--92}}
% the notations of the defined functions: one slot for each explicit argument of the Lean declaration
\newcommand{\chiK}[1]{\chi_{#1}}
\newcommand{\quadTwist}[2]{#1\chi_{#2}}
\newcommand{\entireL}[1]{\tilde L(\,\cdot\,,#1)}
\newcommand{\chowla}[1]{\operatorname{cs}(#1)}
\newcommand{\lambdaClass}[3]{\lambda_{#1,#2}(#3)}
\newcommand{\lambdaSum}[4]{A_{#1,#2}(#3,#4)}
\newcommand{\logDerivMean}[3]{G_{#1,#2}(#3)}
\newcommand{\primeError}[3]{R_{#1,#2}(#3)}
\newcommand{\isPolar}[2]{\operatorname{Polar}(#1,#2)}
\newcommand{\polarWords}[2]{$#1$ is polar with residue function~$#2$}
\newcommand{\mainHeight}[1]{h_{\mathrm{main}}(#1)}
\newcommand{\mainAsk}[1]{\mathrm{Ask}^{\mathrm{main}}(#1)}
\newcommand{\remainder}[1]{\mathrm{rem}(#1)}
\newcommand{\scaledInclusion}[1]{\operatorname{inc}(#1)}
\newcommand{\projectedPoincare}[3]{\mathcal{P}_{#1}(#2)#3}
\newcommand{\integralRamified}[1]{\mathrm{Col}^{\mathrm{ram}}_{\mathrm{int}}(#1)}
\newcommand{\IndexLe}[3]{\operatorname{ind}_{#1}(#2)\le #3}
\newcommand{\indexWords}[3]{$#2$ has $#1$-index at most~$#3$}
\newcommand{\UnitMembers}{29}
\newcommand{\UnitInNote}{29}
\newcommand{\UnitProofs}{18}
\newcommand{\UnitSteps}{104}
\newcommand{\UnitAllSteps}{121}


\linespread{1.08}                       % 11 pt on 14.7 pt
\setlength{\parindent}{0pt}
\setlength{\parskip}{5.5pt plus 1.5pt minus 0.5pt}
\setlength{\abovedisplayskip}{7.5pt plus 2pt minus 3pt}
\setlength{\belowdisplayskip}{7.5pt plus 2pt minus 3pt}
\setlength{\abovedisplayshortskip}{3pt plus 2pt}
\setlength{\belowdisplayshortskip}{6pt plus 2pt minus 2pt}
\emergencystretch=1.5em
\clubpenalty=10000 \widowpenalty=10000 \displaywidowpenalty=10000
\relpenalty=7000 \binoppenalty=9000
\raggedbottom
\pagestyle{fancy}\fancyhf{}\renewcommand{\headrulewidth}{0pt}
\fancyfoot[C]{\sffamily\fontsize{7.6}{9}\selectfont\color{inkfaint}\thepage\,/\,\pageref*{LastPage}}
\setlist{topsep=2pt,partopsep=0pt,parsep=0pt,itemsep=4pt,leftmargin=2.1em,labelsep=0.55em}

\newcommand{\Q}{\mathbb{Q}}
\newcommand{\Z}{\mathbb{Z}}
\newcommand{\F}{\mathbb{F}}
\newcommand{\len}{\operatorname{len}}
\newcommand{\cok}{\operatorname{cok}}
\newcommand{\End}{\operatorname{End}}
\newcommand{\Hom}{\operatorname{Hom}}
\newcommand{\Gal}{\operatorname{Gal}}
\newcommand{\diag}{\operatorname{diag}}
\newcommand{\kw}[1]{\emph{#1}}
\newcommand{\m}[1]{\mbox{$#1$}}        % a short formula that must not be broken across lines
\newcommand{\head}[1]{\par\addvspace{17pt plus 4pt minus 2pt}\Needspace{6\baselineskip}%
  \noindent{\fontsize{13.2}{16}\selectfont\bfseries\color{ink}#1}\par\nopagebreak\vspace{-1pt}}
\newcommand{\sect}[2]{\refstepcounter{section}\label{#2}\par\addvspace{19pt plus 4pt minus 2pt}\Needspace{7\baselineskip}%
  \noindent\llap{{\fontsize{13.2}{16}\selectfont\bfseries\color{inkfaint}\thesection}\hspace{9pt}}%
  {\fontsize{13.2}{16}\selectfont\bfseries\color{ink}#1}\par\nopagebreak\vspace{-1pt}}
\newcommand{\subhead}[1]{\par\addvspace{9pt plus 2pt}\Needspace{4\baselineskip}\noindent\textbf{#1}\enspace\ignorespaces}
% where a statement lives in the Lean archive
\newcommand{\leantag}{{\sffamily\bfseries\fontsize{6.3}{8}\selectfont\color{leancol}\addfontfeatures{LetterSpace=7}LEAN}}
\newcommand{\sep}{\unskip\nobreak\hspace{5pt}{\color{inkfaint}\textperiodcentered}\hspace{5pt plus 2pt}\ignorespaces}
\newcommand{\leanline}[1]{\par\nopagebreak\vspace{-2.5pt}{\fontsize{8.9}{12.4}\selectfont\color{inksoft}\raggedright
  \noindent\llap{\leantag\hspace{9pt}}#1\par}\vspace{1pt}}
% a displayed statement in a tinted box: #1 colour, #2 text
\newcommand{\statement}[2]{\par\addvspace{7pt}\begingroup\setlength{\fboxsep}{8.5pt}\setlength{\fboxrule}{0.6pt}%
  \noindent\fcolorbox{#1!30!white}{#1!3!white}{\parbox{\dimexpr\linewidth-2\fboxsep-2\fboxrule\relax}{#2}}%
  \endgroup\par\addvspace{7pt}}
\newcommand{\nt}[2]{\par\noindent\hangindent=1.1em\hangafter=1 #1\enspace{\color{inksoft}#2}\par\vspace{1pt}}
\newcommand{\legdot}[1]{\tikz[baseline=-0.62ex]\fill[#1] (0,0) circle (2.5pt);\hspace{3pt}}
\newcommand{\crow}[4]{\csname CS@#1\endcsname & #2 & \textcolor{inksoft}{#3} & \textcolor{inksoft}{#4}\\}
\newcommand{\leanquoteline}[2]{\par\noindent\hangindent=\dimexpr#1\fontdimen2\font+2.2em\relax\hangafter=1
  \hspace*{#1\fontdimen2\font}#2}
\newcommand{\quoted}[1]{\par\vspace{3pt}{\ttfamily\fontsize{7.5}{10.4}\selectfont\color{leancol}\raggedright
  \setlength{\parskip}{0pt}\leanquote{#1}\par}\vspace{4pt}}
\newcommand{\bx}[3]{{\sffamily\bfseries\fontsize{7.3}{9.2}\selectfont\color{#1!82!black}#2}\\[1.5pt]#3}
\newcommand{\cs}[2]{\expandafter\newcommand\csname CS@#1\endcsname{#2}}

% theorems
\newtheoremstyle{plainx}{7pt plus 2pt}{7pt plus 2pt}{\itshape}{}{\bfseries}{.}{0.5em}{}
\newtheoremstyle{uprightx}{7pt plus 2pt}{7pt plus 2pt}{}{}{\bfseries}{.}{0.5em}{}
\theoremstyle{plainx}
\newtheorem{theorem}{Theorem}[section]
\newtheorem{lemma}[theorem]{Lemma}
\newtheorem{proposition}[theorem]{Proposition}
\newtheorem{corollary}[theorem]{Corollary}
\theoremstyle{uprightx}
\newtheorem{definition}[theorem]{Definition}
\newtheorem{setting}[theorem]{Setting}          % hypotheses shared by several statements; each of them cites it by number
\newtheorem{remark}[theorem]{Remark}
\newtheorem{example}[theorem]{Example}
\renewcommand{\qedsymbol}{{\setlength{\fboxsep}{0pt}\setlength{\fboxrule}{0.5pt}\fbox{\rule{0pt}{0.46em}\hspace{0.46em}}}}
% what holds throughout the note: the ambient objects only (implicit and instance arguments), never a hypothesis
\newenvironment{conventions}{\par\addvspace{5pt}\noindent\textbf{Conventions.}\enspace\ignorespaces}{\par\addvspace{3pt}}
\makeatletter
\renewenvironment{proof}[1][Proof]{\par\pushQED{\qed}\normalfont\topsep0pt\partopsep0pt
  \trivlist\item[\hskip\labelsep\itshape #1.]\ignorespaces}{\popQED\endtrivlist\@endpefalse}
\makeatother

% ------------------------------------------------------------------------------------- macros of this note
% (the notations of the defined functions of the artifact are made by the build, in data.tex)
\newcommand{\C}{\mathbb{C}}
\newcommand{\RR}{\mathbb{R}}
\newcommand{\NN}{\mathbb{N}}
\newcommand{\ee}{\mathrm{e}}
\newcommand{\dd}{\mathrm{d}}
\newcommand{\re}{\operatorname{Re}}
\newcommand{\im}{\operatorname{Im}}
\newcommand{\ord}{\operatorname{ord}}
\newcommand{\AutC}{\operatorname{Aut}(\C)}
\newcommand{\one}{\mathbf{1}}
\newcommand{\Eu}{\mathrm{C}_{\mathrm{Eu}}}
\newcommand{\Ltv}[2]{\tilde L(#1,#2)}
\newcommand{\Ltd}[2]{\frac{\tilde L'}{\tilde L}(#1,#2)}
\newcommand{\kap}[2]{\kappa_{#1,#2}}
\newcommand{\Ph}[3]{\Phi_{#1,#2}(#3)}
\newcommand{\res}[3]{\operatorname{res}_{#1,#2}(#3)}
\newcommand{\cE}{\mathcal{E}}
\newcommand{\scx}[1]{x(#1)}
\newcommand{\scxo}[1]{x_1(#1)}
\newcommand{\scb}[1]{b(#1)}
\newcommand{\Dm}[1]{m(#1)}
\newcommand{\Dr}[1]{r(#1)}
% Part II: objects of other units, always written with the data
\newcommand{\tops}[1]{\mathrm{tops}(#1)}
\newcommand{\thick}[1]{\mathrm{thick}(#1)}
\newcommand{\budget}[1]{\mathrm{budget}(#1)}
\newcommand{\wD}[2]{w_{#1}(#2)}
\newcommand{\wbar}[1]{\bar{w}_{#1}}
\newcommand{\Ask}[3]{\mathrm{Ask}_{#1}(#2,#3)}
\newcommand{\jet}[2]{\mathrm{jet}_{#1}(#2)}
\newcommand{\cost}[1]{\mathsf{U}(#1)}
\newcommand{\hrow}[1]{h_{0}(#1)}
\newcommand{\Row}[1]{\mathrm{Row}(#1)}
\newcommand{\Col}[1]{\mathrm{Col}(#1)}
\newcommand{\Colram}[1]{\mathrm{Col}^{\mathrm{ram}}(#1)}
\newcommand{\GD}[1]{\mathsf{G}(#1)}
\newcommand{\PF}[1]{\mathrm{PF}(#1)}
\newcommand{\Poin}[1]{\operatorname{Poin}_{#1}}
\newcommand{\gpt}[2]{\varpi_{#1}(#2)}
\newcommand{\Th}[2]{\Theta_{#1,#2}}
\newcommand{\Thf}[2]{\Theta^{\flat}_{#1,#2}}
\newcommand{\AAf}[2]{\mathcal{A}_{#1,#2}}
\newcommand{\nsq}[2]{\|#1\|^{2}_{#2}}
\newcommand{\covol}[2]{\operatorname{covol}^{2}_{#2}(#1)}
\newcommand{\zz}[2]{\mathfrak{z}_{#1}(#2)}
\newcommand{\gain}[2]{\mathrm{gain}_{#1}(#2)}
\newcommand{\err}[2]{\mathrm{err}_{#1}(#2)}
\newcommand{\nonsplit}[1]{\mathrm{nonsplit}(#1)}
\newcommand{\Ch}[2]{\operatorname{Ch}_{#1}(#2)}
\newcommand{\vol}[3]{\operatorname{vol}_{#1}(#2,#3)}
\newcommand{\cst}[2]{\operatorname{cst}^{\mathrm{#1}}_{#2}}
\newcommand{\GamD}[1]{\mathrm{per}_{#1}}
\newcommand{\rr}{\mathfrak{r}}
\newcommand{\iv}[1]{[#1\text{ principal}]}
\newcounter{inp}
\newcommand{\inputhead}[2]{\refstepcounter{inp}\textbf{Input~\theinp} (#1; #2)\textbf{.}\enspace\ignorespaces}
\newcommand{\src}[1]{\par\vspace{2pt}{\fontsize{9.2}{12.4}\selectfont\color{inksoft}#1\par}}

% where each declaration of the concordance is treated: \cs{key}{place}, the place naming a \label or a lettered item
\cs{chiK}{Definition~\ref{d:chi}}
\cs{quadTwist}{Definition~\ref{d:chi}}
\cs{quadTwistapply}{Definition~\ref{d:chi}}
\cs{entireL}{Definition~\ref{d:Lt}}
\cs{chowlaSelbergConst}{Definition~\ref{d:cs}}
\cs{lambdaClass}{Definition~\ref{d:lambda}}
\cs{lambdaSum}{Definition~\ref{d:A}}
\cs{logDerivMean}{Definition~\ref{d:G}}
\cs{primeError}{Definition~\ref{d:R}}
\cs{IsPolar}{Definition~\ref{d:polar}}
\cs{entireLnezero}{Input~\ref{i:nonvanishing}}
\cs{LSeriesHasSummulvonMangoldt}{Input~\ref{i:logderiv}}
\cs{tendstosubmullogDeriventireL}{Input~\ref{i:logderiv}}
\cs{quadTwistneone}{Input~\ref{i:twist}}
\cs{landau}{Input~\ref{i:landau}}
\cs{existsentireLquadTwistnezero}{Input~\ref{i:density}}
\cs{laplacepow}{Input~\ref{i:laplace}}
\cs{integrableOnlaplaceofabsle}{Input~\ref{i:laplace}}
\cs{summabletaylorofremovable}{Input~\ref{i:taylor}}
\cs{lambdaClassnonnegandabsle}{Lemma~\ref{l:lambdabound}}
\cs{abslambdaSumexple}{Lemma~\ref{l:Abound}}
\cs{primeErroreqsum}{Lemma~\ref{l:Rsum}}
\cs{relogDerivMeanone}{Lemma~\ref{l:constant}}
\cs{isPolarprimeError}{Theorem 1}
\cs{entireLnezeroofprimeErrorle}{Theorem 2}
\cs{eventuallymomentle}{Proposition~\ref{p:moments}}
\cs{existswindowsubmulle}{Theorem 3}
\cs{existsratlogeqandpadicnormeq}{Lemma~\ref{l:rational}}
\cs{mainHeight}{Definition~\ref{d:hmain}}
\cs{mainAsk}{Definition~\ref{d:askmain}}
\cs{remainder}{Definition~\ref{d:rem}}
\cs{heightPerRowle}{Proposition~\ref{p:level}}
\cs{absmainHeightaddremainderle}{Lemma~\ref{l:crude}}
\cs{remaindereq}{Theorem 4}
\cs{scaledInclusion}{Definition~\ref{d:inc}}
\cs{projectedPoincare}{Definition~\ref{d:proj}}
\cs{integralRamified}{Definition~\ref{d:intram}}
\cs{IndexLe}{Definition~\ref{d:index}}
\cs{shortforms}{Theorem A}
\cs{enoughforms}{Theorem B}
\cs{heightLevolumeratiosq}{Lemma~\ref{l:ratio}}
\cs{heightLevolumerows}{Lemma~\ref{l:volrows}}
\cs{integralformslong}{Theorem C}
\cs{existsfamilycovolumeSqle}{Corollary~\ref{c:AB}}

\begin{document}
\color{ink}

{\raggedright
{\sffamily\bfseries\fontsize{7.2}{9}\selectfont\color{leancol}\addfontfeatures{LetterSpace=8}READBACK NOTE \,·\, CHAPTER II\par}
\vspace{5pt}
{\fontsize{21.5}{25.5}\selectfont\bfseries From a one-sided bound for primes to a zero-free half-plane\par}
\vspace{5pt}
{\fontsize{12.4}{16}\selectfont\itshape\color{inksoft} One unit of the Lean artifact \ArtifactName, read back into
mathematics\par}}
\vspace{8pt}
{\color{rulegrey}\hrule height 0.6pt}
\vspace{7pt}
\thispagestyle{empty}

\textbf{Summary.} The unit read here closes the argument of the formal proof, and it has two parts. Part~I is analytic
number theory. For a prime $d\equiv 7\pmod 8$ and a modulus $M$ prime to~$d$ it studies the remainder
$\primeError{d}{M}{x}$ of a weighted sum over the prime powers $\equiv 1\pmod M$ that sees only the primes not split in
$\Q(\sqrt{-d})$. The Laplace transform of $v\mapsto\primeError{d}{M}{\ee^{v}}$ is meromorphic. In the half-plane
$\re s>0$ its poles lie at zeros of the Dirichlet $L$-functions modulo~$M$ and of their twists by the Legendre symbol
modulo~$d$, and its residue at such a point~$\rho$ is $1/(\varphi(M)\rho(1-\rho))$ times the difference of two sums
of orders of vanishing at~$\rho$: that of the $L$-functions modulo~$M$ minus that of their twists; here $\varphi$ is
Euler's function (Theorem~1). By Landau's lemma, a one-sided bound
$\primeError{d}{M}{x}\le x^{\beta-1}$, valid for all primes $d$ of one residue class, then forces $L(s,\chi)\ne 0$ for
$\re s>\beta$ and every character $\chi$ modulo~$M$ (Theorem~2); a companion finds, in every window
$[v_0,v_0^{2}]$, a point where a combination of two such remainders is at most $\ee^{-v_0}$ (Theorem~3). Part~II belongs
to the geometric side of the proof: it writes the remainder of the data as a combination of the
$\primeError{d}{M}{x}$ (Theorem~4), and it carries out the last steps of the three statements on theta lifts that the
authors call Theorems~A, B and~C.

\statement{leancol}{\fontsize{9.6}{13.2}\selectfont\textbf{What this note is.} The artifact \ArtifactName{} is a
body of Lean code: the compact edition of a formal proof of the Riemann hypothesis and of its generalisation to
Dirichlet $L$-functions. Its authors report that it compiles and that the kernel of Lean accepts the final theorems
from the three standard axioms; a replay by the owner of this write-up is recorded with the write-up, not here, and
nothing was compiled for this note. No mathematician has read the argument through, so the headline is a
\emph{claim}, and this note does not say that the Riemann hypothesis is proved. It says what one unit of the formal
text asserts and how it argues. The note treats \UnitInNote{} of the \UnitMembers{} declarations of the unit.
\emph{Proved here in full}: everything in Part~I (Theorems 1 to~3 and the lemmas), from seven results of the background
library that are quoted with their statements (Inputs \ref{i:nonvanishing} to~\ref{i:taylor}: among them Landau's
lemma, the non-vanishing of $L$-functions on $\re s\ge1$, and a zero-free box for quadratic twists); and Theorem~4,
from the description of the data. \emph{Deduced here from inputs that carry the substance}: Theorems A, B and~C,
whose proofs in this unit are short deductions from results of other units, stated as Inputs \ref{i:sizes}
to~\ref{i:hadamard} with the unit that owns each. Every statement of the unit, and every quoted statement
of the background library, was compared by the writer with its Lean text, argument by argument. The Lean text read is
\PinText; Appendix~A is the concordance of \ConcDecls{} declarations.}

\sect{The statements of Part I}{s:intro}
The objects first, in one paragraph; the definitions are in §\ref{s:setting}. Let $d$ be a prime number and $M\ge 1$ a
natural number. Write $\Lambda$ for von Mangoldt's function, $\varphi$ for Euler's function and $\chiK{d}$ for the
quadratic character modulo~$d$, the Legendre symbol when $d$ is odd. The \kw{weight} of a natural number $\nu$ is
$\lambdaClass{d}{M}{\nu}=\Lambda(\nu)\bigl(1-\chiK{d}(\nu)\bigr)$ if $\nu\equiv 1\pmod M$, and $0$ otherwise. For
$d\equiv 3\pmod 4$ it is carried by the powers of~$d$ and by the odd powers of the primes that are inert in
$\Q(\sqrt{-d})$, in the residue class $1$ modulo~$M$. The \kw{remainder} is, for real $x\ge 1$,
\[
\primeError{d}{M}{x}=\sum_{1\le\nu\le x}\lambdaClass{d}{M}{\nu}\Bigl(\frac1\nu-\frac1x\Bigr)-\frac{\log x}{\varphi(M)}-\kap{d}{M},
\]
where the constant $\kap{d}{M}$ is formed from the logarithmic derivatives at $s=1$ of the $L$-functions of the
Dirichlet characters $\xi$ modulo~$M$ and of their twists $\quadTwist{\xi}{d}$ (Definitions \ref{d:G} and~\ref{d:R}).
For a Dirichlet character $\chi$ we write $\Ltv{s}{\chi}$ for the entire function $L(s,\chi)$, or $(s-1)L(s,\chi)$ if
$\chi$ is principal, and $\ord_s f$ for the order of vanishing at~$s$ of an entire function $f\ne 0$. A function
$\cE\colon\RR\to\RR$ is \kw{polar with residue function} $\omega\colon\C\to\C$ (Definition~\ref{d:polar}) if
$\cE(v)=O(\ee^{v})$ and the transform $\int_0^\infty\cE(v)\,\ee^{(1-s)v}\,\dd v$ continues to a meromorphic function
$\Phi$ on~$\C$ that has no pole in $\re s\ge 1$ and satisfies $(s-\rho)\Phi(s)\to\omega(\rho)$ as $s\to\rho$, for every
$\rho$ with $\re\rho>0$.

\statement{arith}{\phantomsection\label{t:1}\textbf{Theorem 1} (zeros as poles)\textbf{.}\enspace Let $d$ be a prime
number with $d\equiv 7\pmod 8$ and let $M\ge 1$ be a natural number with $\gcd(M,d)=1$. For a complex number $s$ put
\[
\res{d}{M}{s}=\frac{1}{\varphi(M)\,s(1-s)}\sum_{\xi}\Bigl(\ord_{s}\entireL{\xi}-\ord_{s}\entireL{\quadTwist{\xi}{d}}\Bigr)
\quad\text{for }s\ne 0,1,
\]
the sum over the Dirichlet characters $\xi$ modulo~$M$, and put $\res{d}{M}{0}=\res{d}{M}{1}=0$. Then, as a function of the real variable~$v$,
\polarWords{v\mapsto\primeError{d}{M}{\ee^{v}}}{s\mapsto\res{d}{M}{s}}.}

\statement{arith}{\phantomsection\label{t:2}\textbf{Theorem 2} (the one-sided criterion)\textbf{.}\enspace Let $M\ge 1$
be a natural number, $c$ an element of $\Z/M$ and $\beta$ a real number. Assume that $\beta>\tfrac12$, that $c$ is a
unit of $\Z/M$, and that every prime number $d$ with $d\equiv c\pmod M$ has the following three properties:
$d\equiv 7\pmod 8$; $\gcd(M,d)=1$; and
\[\primeError{d}{M}{x}\le x^{\beta-1}\qquad\text{for all sufficiently large real }x.\]
Then for every complex number $\rho$ with $\re\rho>\beta$ and every Dirichlet character $\chi$ modulo~$M$,
\[\Ltv{\rho}{\chi}\ne 0.\]}

\statement{arith}{\phantomsection\label{t:3}\textbf{Theorem 3} (a good point in every window)\textbf{.}\enspace Let
$\cE_1,\cE_2\colon\RR\to\RR$ and $\omega_1,\omega_2\colon\C\to\C$ be functions such that
\polarWords{\cE_1}{\omega_1} and \polarWords{\cE_2}{\omega_2}. Then there is a real number $v_\ast$ such that for every
real $v_0\ge v_\ast$ and every real $X$ with $0\le X\le\ee^{v_0}$ there is a real number $v_1$ with
$v_0\le v_1\le v_0^{2}$ and
\[\cE_1(v_1)-X\,\cE_2(v_1)\le\ee^{-v_0}.\]}

In Theorem~2 the conclusion gives $L(\rho,\chi)\ne 0$, the point $\rho=1$ excepted when $\chi$ is principal. The three
theorems are used as follows in the last unit of the proof (Chapter~I of the write-up). The geometric part of the
proof bounds $\primeError{d}{N_0}{x}$ from above by $x^{8\varepsilon-1/2}$ for every $\varepsilon\in(0,\tfrac1{100})$,
for a multiple $N_0$ of $18480$ and every prime $d\equiv N_0/2-1\pmod{N_0}$; Theorems 1 and~3 enter that deduction at
one point, the choice of the second scale. Theorem~2 with $M=N_0$ and $c=N_0/2-1$ then turns the bound into
$\Ltv{\rho}{\chi}\ne 0$ for $\re\rho>\tfrac12$. The bound is one-sided, and it is needed for every prime of a class:
§\ref{s:criterion} explains both points. Part~II (§§\ref{s:data} to~\ref{s:C}) needs the vocabulary of the data and of
the forms, and its statements (Theorems 4, A, B and~C) stand there.

\sect{The setting of Part I}{s:setting}
\begin{conventions}\phantomsection\label{c:conventions}
$\NN=\{0,1,2,\dots\}$. For a natural number $M\ge 1$, a \kw{Dirichlet character modulo~$M$} is a multiplicative
map $\Z/M\to\C$ that vanishes off the units; it is read as a function on $\Z$, and it is \kw{principal} if it is $1$
on the units. Sums over $\xi$ run over all Dirichlet characters modulo~$M$, and $\varphi(M)$ is their number. $\one$ is
the Dirichlet character modulo~$1$. $\Lambda$ is von Mangoldt's function, with $\Lambda(0)=0$; $\log$ is the natural
logarithm, $\ee$ its base, and $(z)_{+}=\max(z,0)$ for real~$z$. For an entire function $f\ne0$ and $s\in\C$, $\ord_s f\in\NN$ is the
order of vanishing of $f$ at~$s$. A function is \kw{analytic at a point} if it is holomorphic in a neighbourhood of
it. \kw{For all sufficiently large $x$} means: for all $x\ge x_\ast$, for some real~$x_\ast$. Integrals are Lebesgue
integrals; a family of complex numbers is \kw{summable} if it is absolutely summable. In Part~I the letter $\nu$ is a natural
number, the variable of the prime sums, and in Part~II an integer; it is never a character of the ideles (Convention~0.6). A bare $x$ is a
free real variable (Convention~0.3); in Part~II the scales of the data are written $\scx{D}$, $\scxo{D}$, $\scb{D}$.
\end{conventions}

\begin{definition}[the quadratic character and the twists]\label{d:chi}
Let $d$ be a prime number. The Dirichlet character $\chiK{d}$ modulo~$d$ is the quadratic character of the field
$\Z/d$, with values in $\{-1,0,1\}$: $\chiK{d}(n)=0$ if $d\mid n$, $\chiK{d}(n)=1$ if $n$ is a non-zero square
modulo~$d$, and $\chiK{d}(n)=-1$ otherwise. For odd $d$ it is the Legendre symbol, and for $d\equiv 3\pmod 4$ it is the
quadratic character $\chi_K$ of the field $K=\Q(\sqrt{-d})$. Let moreover $M\ge 1$ be a natural number and $\xi$ a
Dirichlet character modulo~$M$. The \kw{twist} $\quadTwist{\xi}{d}$ is the Dirichlet character modulo $Md$ with
$(\quadTwist{\xi}{d})(n)=\xi(n)\,\chiK{d}(n)$ for every integer~$n$.
\end{definition}
\leanline{\lean{chiK}\sep \lean{quadTwist}\sep \lean{quadTwistapply}}

\begin{definition}[the entire $L$-function]\label{d:Lt}
Let $M\ge1$ be a natural number and $\chi$ a Dirichlet character modulo~$M$. The function
$\entireL{\chi}\colon\C\to\C$, with values written $\Ltv{s}{\chi}$, is the Dirichlet $L$-function
$s\mapsto L(s,\chi)$ if $\chi$ is not principal, and the entire function that equals $(s-1)L(s,\chi)$ for $s\ne 1$
if $\chi$ is principal (its value at $s=1$ is then $\prod_{p\mid M}(1-1/p)$, the residue of $L(s,\chi)$). In both cases it is entire and not identically zero. Its derivative is $\tilde L'$, and
$\Ltd{s}{\chi}$ is the quotient of the two values at~$s$.
\end{definition}
\leanline{\lean{entireL}}

\begin{definition}[the Chowla--Selberg constant]\label{d:cs}
Let $d$ be a prime number. Then
\[\chowla{d}=\Eu-\re\Ltd{1}{\quadTwist{\one}{d}},\]
where $\Eu$ is Euler's constant and $\quadTwist{\one}{d}$ is the twist of the character modulo~$1$, that is, the
character $n\mapsto\chiK{d}(n)$ modulo~$d$.
\end{definition}
\leanline{\lean{chowlaSelbergConst}}

\begin{definition}[the weight]\label{d:lambda}
Let $d$ be a prime number and let $M$ and $\nu$ be natural numbers. Then
\[
\lambdaClass{d}{M}{\nu}=\Lambda(\nu)\bigl(1-\chiK{d}(\nu)\bigr)\quad\text{if }\nu\equiv 1\pmod M,\qquad
\lambdaClass{d}{M}{\nu}=0\quad\text{otherwise.}
\]
\end{definition}
\leanline{\lean{lambdaClass}}

So $\lambdaClass{d}{M}{\nu}\ne 0$ only if $\nu=p^{e}$ is a power of a prime, $\nu\equiv 1\pmod M$, and either $p=d$
(the value is $\log d$) or $\chiK{d}(p)=-1$ and $e$ is odd (the value is $2\log p$).

\begin{definition}[the weighted sums]\label{d:A}
Let $d$ be a prime number, let $M$ and $i$ be natural numbers and $x$ a real number. Then
\[\lambdaSum{d}{M}{i}{x}=x^{-i}\sum_{1\le\nu\le x}\lambdaClass{d}{M}{\nu}\,\nu^{\,i-1},\]
the sum over the natural numbers $\nu$ with $1\le\nu\le x$; it is empty, and the value is $0$, for $x<1$.
\end{definition}
\leanline{\lean{lambdaSum}}

\begin{definition}[the mean of the logarithmic derivatives]\label{d:G}
Let $d$ be a prime number, $M\ge 1$ a natural number and $s$ a complex number. Then
\[
\logDerivMean{d}{M}{s}=\frac{1}{\varphi(M)}\Bigl(\sum_{\xi}\Bigl(\Ltd{s}{\quadTwist{\xi}{d}}-\Ltd{s}{\xi}\Bigr)-1\Bigr),
\qquad \kap{d}{M}=\re\logDerivMean{d}{M}{1},
\]
the sum over the Dirichlet characters $\xi$ modulo~$M$. As a function of~$s$ it is meromorphic on~$\C$, and it is
analytic at every point where none of the functions $\entireL{\xi}$ and $\entireL{\quadTwist{\xi}{d}}$ vanishes.
\end{definition}
\leanline{\lean{logDerivMean}}

\begin{definition}[the remainder]\label{d:R}
Let $d$ be a prime number, $M\ge 1$ a natural number and $x$ a real number. Then
\[
\begin{aligned}
\primeError{d}{M}{x}&=\lambdaSum{d}{M}{0}{x}-\lambdaSum{d}{M}{1}{x}-\frac{\log x}{\varphi(M)}-\kap{d}{M}\\
&=\sum_{1\le\nu\le x}\lambdaClass{d}{M}{\nu}\Bigl(\frac1\nu-\frac1x\Bigr)-\frac{\log x}{\varphi(M)}-\kap{d}{M},
\end{aligned}
\]
with $\kap{d}{M}=\re\logDerivMean{d}{M}{1}$ as in Definition~\ref{d:G}.
\end{definition}
\leanline{\lean{primeError}}

\begin{definition}[polar functions]\label{d:polar}
Let $\cE\colon\RR\to\RR$ and $\omega\colon\C\to\C$ be functions. We say that \polarWords{\cE}{\omega}, and write
$\isPolar{\cE}{\omega}$, if $\cE$ is Borel measurable and there are a real number~$C$ and a function
$\Phi\colon\C\to\C$, meromorphic on~$\C$ and defined at every point (at a pole its value is arbitrary), with the
following four properties.
\begin{enumerate}[label=(P\arabic*)]
\item $|\cE(v)|\le C\,\ee^{v}$ for every real $v\ge 0$.
\item $\Phi(s)=\int_0^\infty\cE(v)\,\ee^{(1-s)v}\,\dd v$ for every $s\in\C$ with $\re s>2$.
\item The function $s\mapsto(s-1)\Phi(s)$, whose value at $s=1$ is $0$ whatever $\Phi(1)$ is, is analytic at every
point of the closed half-plane $\re s\ge 1$.
\item For every $\rho\in\C$ with $\re\rho>0$: $(s-\rho)\Phi(s)\to\omega(\rho)$ as $s\to\rho$, $s\ne\rho$.
\end{enumerate}
\end{definition}
\leanline{\lean{IsPolar}}

\begin{remark}[how to read the definition]\label{r:polar}
By (P1) and the measurability, the integral in (P2) converges absolutely for $\re s>2$, so $\Phi$ is the meromorphic
continuation of the Laplace transform of $v\mapsto\cE(v)\ee^{v}$, and it is unique. (P4) says that the poles of $\Phi$
in the half-plane $\re s>0$ are at most simple, and that $\omega(\rho)$ is the residue at~$\rho$; so
$\omega(\rho)=0$ exactly when $\Phi$ is analytic at~$\rho$ (after correcting its value there). The clause of (P3)
on the value at $s=1$ is essential: (P3) says that $\Phi$ is analytic at
every point $s\ne 1$ of $\re s\ge1$, and that near $s=1$ it is the quotient by $s-1$ of an analytic function that
vanishes at~$1$: the singularity at $s=1$ is removable. Hence a polar function has $\omega(\rho)=0$ for $\re\rho\ge 1$.
Read without the clause, for a meromorphic function with its removable singularities filled in, (P3) would allow a
simple pole of $\Phi$ at $s=1$; the constant function~$1$, with $\Phi(s)=1/(s-1)$, would then be polar, and
Theorem~3 would be false for the constant functions $1$ and~$0$.
The values of $\omega$ on $\re\rho\le0$ play no role.
\end{remark}

\begin{remark}[what the remainder measures; not part of the formal text]\label{r:heuristic}
By the orthogonality of characters, for $\re s>1$ and $\gcd(M,d)=1$, $d\ne 2$,
\[
\sum_{\nu\ge1}\lambdaClass{d}{M}{\nu}\,\nu^{-s}=\logDerivMean{d}{M}{s}+\frac{s}{\varphi(M)(s-1)}
\]
(this is step~3 of the proof of Theorem~1). The only singularity on $\re s\ge 1$ is the simple pole at $s=1$ with
residue $1/\varphi(M)$, which comes from the principal character. Accordingly one expects
$\sum_{\nu\le x}\lambdaClass{d}{M}{\nu}/\nu=\log x/\varphi(M)+\text{constant}+o(1)$ and
$x^{-1}\sum_{\nu\le x}\lambdaClass{d}{M}{\nu}=1/\varphi(M)+o(1)$, and the constant $\kap{d}{M}$ is the one for which
$\primeError{d}{M}{x}$ has no constant term: in the language of Definition~\ref{d:polar}, the one for which the
transform has no pole at $s=1$. A zero $\rho$ of some $\entireL{\xi}$ with $0<\re\rho<1$ gives the transform
a pole with residue $\ord_\rho\entireL{\xi}/(\varphi(M)\rho(1-\rho))$, hence, formally, a term of size $x^{\re\rho-1}$ in
$\primeError{d}{M}{x}$; a zero of a twist $\entireL{\quadTwist{\xi}{d}}$ does the same with the opposite sign. So the
bound $\primeError{d}{M}{x}=O(x^{\beta-1})$ is what one expects if there is no zero with $\re\rho>\beta$, and the
generalised Riemann hypothesis corresponds to $\beta=\tfrac12$. Theorem~2 goes in the other direction, and it asks for
an upper bound only.
\end{remark}

\sect{Quoted results}{s:inputs}
Part~I rests on seven results that the artifact proves in its background library, from Mathlib, and we quote all
seven. Six of them are classical or elementary. The seventh, Input~\ref{i:density}, is not: the theorem that its
source names treats real characters, and the formal statement also covers the quadratic twists of non-real
characters, by an estimate for each fixed character that the library proves itself; its box says exactly what is
proved and what the cited theorem gives. Each box gives the statement as the formal text has it, the unit that owns it, and
what is known about its source: the comparison with the literature that the second reader of this write-up made
for three of them, and otherwise the entry of the catalogue of quoted results. \emph{Reference to be confirmed}
means that there is no reference with a confirmed place.

\statement{ana}{\inputhead{non-vanishing on $\re s\ge1$}{unit~107}\label{i:nonvanishing}
Let $M\ge 1$ be a natural number, $\chi$ a Dirichlet character modulo~$M$ and $s$ a complex number with
$\re s\ge 1$. Then $\Ltv{s}{\chi}\ne0$.
\src{Classical: Hadamard and de la Vall\'ee Poussin on the line $\re s=1$, Dirichlet at $s=1$. In the formal text
it is a direct consequence of two theorems of Mathlib, at the version that the artifact pins, after the principal
case is split off: for non-principal $\chi$, $L(s,\chi)\ne0$ on $\re s\ge1$; for principal $\chi$ and $s\ne1$ both
factors of $(s-1)L(s,\chi)$ are non-zero; and at $s=1$ the entire function has the value $\prod_{p\mid M}(1-1/p)\ne0$.
Relation: the theorem itself. The comparison with Mathlib at the pin was made by the second reader of this
write-up.}}
\leanline{\lean{entireLnezero}}

\statement{ana}{\inputhead{the logarithmic derivative}{unit~107}\label{i:logderiv}
Let $M\ge 1$ be a natural number, $\chi$ a Dirichlet character modulo~$M$ and $s$ a complex number.

\emph{(a)} If $\re s>1$, the family $\bigl(\chi(\nu)\Lambda(\nu)\nu^{-s}\bigr)_{\nu\ge1}$ is summable, with sum
\[\sum_{\nu\ge 1}\chi(\nu)\Lambda(\nu)\,\nu^{-s}=\frac{\iv{\chi}}{s-1}-\Ltd{s}{\chi},\]
where the bracket $\iv{\chi}$ is the number $1$ if $\chi$ is principal and $0$ otherwise.

\emph{(b)} $\ord_s\entireL{\chi}=0$ if and only if $\Ltv{s}{\chi}\ne0$; and
$(s'-s)\,\Ltd{s'}{\chi}\to\ord_s\entireL{\chi}$ as $s'\to s$, $s'\ne s$.
\src{Elementary: (a) is the logarithmic derivative of the Euler product, (b) the residue of the logarithmic
derivative of an analytic function. The catalogue names Davenport, \emph{Multiplicative Number Theory}, for (a),
without a place: reference to be confirmed.}}
\leanline{\lean{LSeriesHasSummulvonMangoldt}\sep \lean{tendstosubmullogDeriventireL}}

\statement{ana}{\inputhead{twists are not principal}{unit~107}\label{i:twist}
Let $M$ be a natural number, $\xi$ a Dirichlet character modulo~$M$ and $d$ a prime number with $d\neq 2$ and
$\gcd(M,d)=1$. Then the twist $\quadTwist{\xi}{d}$ is not the principal character modulo~$Md$.
\src{Elementary (the Chinese remainder theorem gives $n$ with $n\equiv1\pmod M$ and $\chiK{d}(n)=-1$).}}
\leanline{\lean{quadTwistneone}}

\statement{ana}{\inputhead{Landau's lemma}{unit~107}\label{i:landau}
Let $\cE\colon\RR\to\RR$ and $\omega\colon\C\to\C$ be functions such that \polarWords{\cE}{\omega}, and let
$\beta$ be a real number with $\beta>0$. Assume that $\re\omega(s)\ge 0$ for every real number $s>\beta$, and that
$\cE(v)\le\ee^{(\beta-1)v}$ for all sufficiently large real~$v$. Then $\omega(\rho)=0$ for every complex number
$\rho$ with $\re\rho>\beta$.
\src{Relation to the literature, as compared by the second reader of this write-up from the pinned Lean text and
the primary texts (the writer has not consulted them): a corollary of D.~V.~Widder, \emph{The Laplace Transform},
Chapter~II, §5, Theorem~5b, p.~58 [Wi], with a residue argument. Widder's theorem concerns the
Laplace--Stieltjes transform of a monotone function with a finite abscissa of convergence, and gives a singularity at
the real point of the abscissa; it applies to the primitive of the non-negative measurable function
$v\mapsto C_1\ee^{\beta v}-\cE(v)\ee^{v}$, for a suitable constant~$C_1$. What the formal statement adds is the step at that real point, which
compares the sign of the residue with positivity (Remark~\ref{r:landau}). The header of the formal file names
H.~L.~Montgomery and R.~C.~Vaughan, \emph{Multiplicative Number Theory~I}, §15.1, Lemma~15.1, p.~463 [MV], as the
source; that lemma is the same theorem for a density that is Riemann integrable on bounded intervals and of
eventually constant sign, whereas $\cE$ is only assumed measurable here, so it does not apply as it stands.}}
\leanline{\lean{landau}}

\statement{ana}{\inputhead{a zero-free box for quadratic twists}{unit~127}\label{i:density}
Let $M\ge 1$ be a natural number, $\beta$ and $z$ real numbers and $c$ an element of $\Z/M$. Assume that
$\beta>\tfrac12$, that $c$ is a unit of $\Z/M$, and that every prime number $d$ with $d\equiv c\pmod M$ satisfies
$d\equiv 7\pmod 8$ and $\gcd(M,d)=1$. Then there is a prime number $d$ with $d\equiv c\pmod M$ such that
\[\Ltv{s}{\quadTwist{\xi}{d}}\ne 0\]
for every Dirichlet character $\xi$ modulo~$M$ and every complex number $s$ with $\re s\ge\beta$ and $|\im s|\le z$.
\src{The number $z$ is any real number; for $z<0$ the condition on $s$ is empty, and the statement still gives a
prime. \emph{What the formal library proves}, in two steps that the writer read in the Lean text. First, an
estimate for one character: for every Dirichlet character $\xi$ modulo~$M$, real or not, every real number $\Delta>0$
and every real~$z$, the number of the primes $d\le x$ with $d\equiv7\pmod 8$ and $\gcd(M,d)=1$ for which
$\entireL{\quadTwist{\xi}{d}}$ has a zero $s$ with $\re s\ge\tfrac12+9\Delta$ and $|\im s|\le z$ is $O(x^{1-\Delta/2})$ as
$x\to\infty$. The proof passes to the primitive character that induces~$\xi$, and uses a three-lines argument and a
smoothed mean square of sums of Jacobi symbols. Second, the bounds for the $\varphi(M)$ characters $\xi$ are added
and set against the divergence of $\sum 1/p$ over the primes $p\equiv c\pmod M$: if every prime of the class failed
for some~$\xi$, that sum would converge. \emph{What the literature gives}, as compared by the second reader of this
write-up from the pinned Lean text and the primary text (the writer has not consulted it): the header of the formal
file names M.~Jutila, On mean values of Dirichlet polynomials with real characters, Acta Arith.\ 27 (1975),
191--198; its Theorem~2, equation~(6), p.~193 [Ju], bounds the number of zeros in a box, summed over real
characters. For a non-real character $\xi$ the twists $\quadTwist{\xi}{d}$ are not real, so that theorem alone does
not give the first step, and it does not contain the second. Input~\ref{i:density} is therefore not a classical
theorem in the form used here: it is the combination of a variant of a classical zero-density estimate, for which
no published place has been confirmed, with the divergence of a sum over primes. The catalogue of quoted results
records it in this way.}}
\leanline{\lean{existsentireLquadTwistnezero}}

\statement{ana}{\inputhead{Laplace integrals}{unit~107}\label{i:laplace}
\emph{(a)} Let $s$ be a complex number with $\re s>0$ and $i$ a natural number. Then
$\int_0^\infty v^{i}\,\ee^{-sv}\,\dd v=i!/s^{i+1}$.

\emph{(b)} Let $\cE\colon\RR\to\RR$ be a Borel measurable function, $C$ a real number with $|\cE(v)|\le C\,\ee^{v}$
for every real $v\ge 0$, and $s$ a complex number with $\re s>2$. Then $v\mapsto\cE(v)\,\ee^{(1-s)v}$ is integrable
on $(0,\infty)$.
\src{Elementary: (a) is the integral of the Gamma function.}}
\leanline{\lean{laplacepow}\sep \lean{integrableOnlaplaceofabsle}}

\statement{ana}{\inputhead{Taylor series across a removable singularity}{unit~107}\label{i:taylor}
Let $\Phi\colon\C\to\C$ be a function, $\rr_0$ a real number and $s_0,s_1,s_2$ complex numbers. Assume that
the function $s\mapsto(s-s_1)\Phi(s)$ is analytic at every point of the open disc $|s-s_0|<\rr_0$, and that
$s_0\neq s_1$, $|s_1-s_0|<\rr_0$ and $|s_2-s_0|<\rr_0$. Then the family
$\bigl((s_2-s_0)^{i}\,\Phi^{(i)}(s_0)/i!\bigr)_{i\ge 0}$ is summable, where $\Phi^{(i)}$ is the $i$-th complex
derivative.
\src{Elementary. The function $s\mapsto(s-s_1)\Phi(s)$ vanishes at $s_1$, because it is a product with the factor
$s-s_1$ whatever the value $\Phi(s_1)$ is; being analytic, it is divisible by $s-s_1$. So off $s_1$ the function
$\Phi$ agrees with a function analytic in the disc, with the same derivatives at $s_0$, and the Taylor series of
that function at $s_0$ converges absolutely in the disc. Analyticity of the product without the vanishing would
allow a simple pole of $\Phi$ at~$s_1$.}}
\leanline{\lean{summabletaylorofremovable}}

\begin{remark}[from the classical lemma to Input~\ref{i:landau}; our derivation]\label{r:landau}
In the formal library Input~\ref{i:landau} is proved directly, by the Taylor expansion at a real point that proves
the classical lemma. Here is how it follows from the classical statement in this form (its relation to Widder's
theorem is described in the box of Input~\ref{i:landau}): \emph{if $\cE_0\ge 0$ on $(0,\infty)$ and the
Laplace integral $\int_0^\infty\cE_0(v)\ee^{-sv}\dd v$ converges for some real~$s$ but not for all, then the function
that it defines to the right of its abscissa of convergence~$x_0$ has no analytic continuation to a neighbourhood of
the point~$x_0$.} Let $C$ and $\Phi$ be as in Definition~\ref{d:polar}, let $\cE(v)\le\ee^{(\beta-1)v}$ for
$v\ge v_\ast\ge0$, and put $C_1=\max(1,|C|\ee^{2v_\ast})$. Then
$\cE_0(v)=C_1\,\ee^{\beta v}-\cE(v)\ee^{v}$ is $\ge0$ for all $v>0$: for $v\ge v_\ast$ because
$\cE(v)\ee^{v}\le\ee^{\beta v}$, and for $v<v_\ast$ because $\cE(v)\ee^{v}\le|C|\ee^{2v_\ast}\le C_1\ee^{\beta v}$. Its
Laplace integral converges for $\re s>\max(2,\beta)$, where by (P2) it equals $C_1/(s-\beta)-\Phi(s)$, a function
meromorphic on~$\C$. Suppose that the abscissa of convergence $x_0$ is $>\beta$. To the right of~$x_0$ the integral
is analytic, so by the identity theorem it equals $C_1/(s-\beta)-\Phi(s)$ on all of $\re s>x_0$; for real $s>x_0$ it is
$\ge0$. Multiplying by $s-x_0>0$ and letting $s$ decrease to~$x_0$, (P4) gives $-\omega(x_0)\ge0$. With the hypothesis
$\re\omega(x_0)\ge0$ this forces $\omega(x_0)=0$, so $C_1/(s-\beta)-\Phi(s)$ is analytic at~$x_0$, against the classical
lemma. Hence $x_0\le\beta$: the integral is analytic in $\re s>\beta$ and equals $C_1/(s-\beta)-\Phi(s)$ there, so
$\Phi$ has no pole in $\re s>\beta$, and $\omega(\rho)=0$ for $\re\rho>\beta$ by (P4). Note that the hypothesis does not
ask that $\Phi$ be analytic at the real points to the right of~$\beta$: poles there whose residue has non-negative
real part are allowed, and the conclusion says that there are none. This is the form in which the lemma is used
below.
\end{remark}

\sect{Elementary properties of the weight and of the remainder}{s:elementary}

\begin{lemma}[bounds for the weight]\label{l:lambdabound}
Let $d$ be a prime number and let $M$ and $\nu$ be natural numbers. Then
$0\le\lambdaClass{d}{M}{\nu}$ and $|\lambdaClass{d}{M}{\nu}|\le 2\nu$.
\end{lemma}
\begin{proof}
$0\le\Lambda(\nu)\le\log\nu\le\nu$ for $\nu\ge1$, $\Lambda(0)=0$, and $1-\chiK{d}(\nu)\in\{0,1,2\}$.
\end{proof}
\leanline{\lean{lambdaClassnonnegandabsle}}

\begin{lemma}[bound for the weighted sums]\label{l:Abound}
Let $d$ be a prime number, let $M$ and $i$ be natural numbers and $v$ a real number. Then
$|\lambdaSum{d}{M}{i}{\ee^{v}}|\le 2\,\ee^{v}$.
\end{lemma}
\begin{proof}
Put $x=\ee^{v}$. By Lemma~\ref{l:lambdabound}, $|\lambdaClass{d}{M}{\nu}|\,\nu^{\,i-1}\le 2\nu^{\,i}\le 2x^{i}$ for
$1\le\nu\le x$, and the sum in Definition~\ref{d:A} has at most $x$ terms. So
$|\lambdaSum{d}{M}{i}{x}|\le x^{-i}\cdot x\cdot 2x^{i}=2x$.
\end{proof}
\leanline{\lean{abslambdaSumexple}}

\begin{lemma}[the remainder as a sum of positive parts]\label{l:Rsum}
Let $d$ be a prime number, let $M\ge 1$ and $a$ be natural numbers, and let $x$ be a real number with $0<x$ and
$x\le a$. Then
\[
\primeError{d}{M}{x}=\frac1a\sum_{\nu=1}^{a}\lambdaClass{d}{M}{\nu}\Bigl(\frac a\nu-\frac ax\Bigr)_{+}-\frac{\log x}{\varphi(M)}-\kap{d}{M}.
\]
\end{lemma}
\begin{proof}
$a\ge x>0$, so $a\ge 1$. For $1\le\nu\le a$ the number $\frac1a(\frac a\nu-\frac ax)_{+}=(\frac1\nu-\frac1x)_{+}$ is
$\frac1\nu-\frac1x$ if $\nu\le x$ and $0$ if $\nu>x$. So the sum is
$\sum_{1\le\nu\le x}\lambdaClass{d}{M}{\nu}(\frac1\nu-\frac1x)=\lambdaSum{d}{M}{0}{x}-\lambdaSum{d}{M}{1}{x}$, and the
claim is Definition~\ref{d:R}.
\end{proof}
\leanline{\lean{primeErroreqsum}}

This identity is the reason for the shape of the remainder. In Part~II the prime sums of the data have the form
$\sum_{\nu\le a}\lambdaClass{d}{M}{\nu}(z/\nu-Y)_{+}$ with a real \emph{order} $z$ and a real \emph{weight} $Y$; with
$z=a$ and $Y=a/x$ this is $a$ times the sum of the lemma.

\begin{lemma}[the constant for the modulus $1$]\label{l:constant}
Let $d$ be a prime number. Then $\re\logDerivMean{d}{1}{1}=-1-\chowla{d}$; in other words
$\kap{d}{1}=-1-\chowla{d}$.
\end{lemma}
\begin{proof}
Modulo~$1$ there is one Dirichlet character, $\one$, and $\varphi(1)=1$; so
$\logDerivMean{d}{1}{s}=\Ltd{s}{\quadTwist{\one}{d}}-\Ltd{s}{\one}-1$. The function $\entireL{\one}$ is
$s\mapsto(s-1)\zeta(s)$, where $\zeta$ is the Riemann zeta function; from the expansion
$\zeta(s)=1/(s-1)+\Eu+O(s-1)$ at $s=1$ its value at $1$ is $1$ and its derivative at $1$ is $\Eu$, so
$\Ltd{1}{\one}=\Eu$. Hence $\re\logDerivMean{d}{1}{1}=\re\Ltd{1}{\quadTwist{\one}{d}}-\Eu-1=-\chowla{d}-1$ by
Definition~\ref{d:cs}.
\end{proof}
\leanline{\lean{relogDerivMeanone}}

\sect{The transform of the remainder: proof of Theorem 1}{s:polar}
Let $d$ and $M$ be as in Theorem~1: $d$ is a prime with $d\equiv7\pmod 8$, $M\ge1$ and $\gcd(M,d)=1$. Put
\[
\Ph{d}{M}{s}=\frac{1}{s-1}\Bigl(\frac{\logDerivMean{d}{M}{s}}{s}-\kap{d}{M}\Bigr)\qquad(s\in\C,\ s\ne1),
\tag{5.1}
\]
and $\Ph{d}{M}{1}=0$; the value at $s=1$ plays no role. We verify Definition~\ref{d:polar} for $\cE(v)=\primeError{d}{M}{\ee^{v}}$ and $\omega(s)=\res{d}{M}{s}$, with the
constant $C=|\kap{d}{M}|+5$ and the function $\Phi=\Phi_{d,M}$. The proof has seven steps.

\subhead{Step 1: analyticity.} By Input~\ref{i:nonvanishing} none of the functions $\entireL{\xi}$ and
$\entireL{\quadTwist{\xi}{d}}$, $\xi$ modulo~$M$, has a zero in $\re s\ge1$. So $s\mapsto\logDerivMean{d}{M}{s}$ is
analytic at every point of $\re s\ge1$ (Definition~\ref{d:G}), and so is the function
$s\mapsto\logDerivMean{d}{M}{s}/s-\kap{d}{M}$. As a quotient of meromorphic functions, $\Phi_{d,M}$ is meromorphic
on~$\C$.

\subhead{Step 2: measurability and (P1).} For $i=0,1$ the function $v\mapsto\lambdaSum{d}{M}{i}{\ee^{v}}$ is the
product of $\ee^{-iv}$ and a step function, and $v\mapsto v/\varphi(M)+\kap{d}{M}$ is continuous; so
$v\mapsto\primeError{d}{M}{\ee^{v}}$ is Borel measurable. For $v\ge0$, Lemma~\ref{l:Abound} gives
\[
|\primeError{d}{M}{\ee^{v}}|\le 2\ee^{v}+2\ee^{v}+\frac{v}{\varphi(M)}+|\kap{d}{M}|\le\bigl(|\kap{d}{M}|+5\bigr)\,\ee^{v},
\]
because $v\le\ee^{v}$, $\varphi(M)\ge1$ and $1\le\ee^{v}$.

\subhead{Step 3: the Dirichlet series of the weight.} For $\re s>1$ the family
$(\lambdaClass{d}{M}{\nu}\nu^{-s})_{\nu\ge1}$ is summable, and
\[
\sum_{\nu\ge1}\lambdaClass{d}{M}{\nu}\,\nu^{-s}=\logDerivMean{d}{M}{s}+\frac{s}{\varphi(M)\,(s-1)}.
\tag{5.2}
\]
Indeed, by the orthogonality of characters the indicator of the class $\nu\equiv1\pmod M$ is
$\varphi(M)^{-1}\sum_\xi\xi(\nu)$, so
$\lambdaClass{d}{M}{\nu}=\varphi(M)^{-1}\sum_\xi\bigl(\xi(\nu)-(\quadTwist{\xi}{d})(\nu)\bigr)\Lambda(\nu)$ by
Definitions \ref{d:lambda} and~\ref{d:chi}. Input~\ref{i:logderiv}(a), for the characters $\xi$ modulo~$M$ and
$\quadTwist{\xi}{d}$ modulo~$Md$, gives
\[
\sum_{\nu\ge1}\lambdaClass{d}{M}{\nu}\,\nu^{-s}
=\frac1{\varphi(M)}\sum_\xi\Bigl(\frac{\iv{\xi}-\iv{\quadTwist{\xi}{d}}}{s-1}-\Ltd{s}{\xi}+\Ltd{s}{\quadTwist{\xi}{d}}\Bigr).
\]
Here $\iv{\quadTwist{\xi}{d}}=0$ for every $\xi$, by Input~\ref{i:twist}: $d\ne2$ because $d\equiv7\pmod 8$, and
$\gcd(M,d)=1$. And $\sum_\xi\iv{\xi}=1$. So the right side is
$\logDerivMean{d}{M}{s}+\varphi(M)^{-1}\bigl(1+\frac1{s-1}\bigr)$, which is~(5.2). This is the only place where the
hypothesis $d\equiv7\pmod8$ enters the proof.

\subhead{Step 4: the transforms of the weighted sums.} Let $\re s>2$ and $i\in\{0,1\}$. By Lemma~\ref{l:Abound} and
Input~\ref{i:laplace}(b) the function $v\mapsto\lambdaSum{d}{M}{i}{\ee^{v}}\ee^{(1-s)v}$ is integrable on
$(0,\infty)$, and
\[
(s+i-1)\int_0^\infty\lambdaSum{d}{M}{i}{\ee^{v}}\,\ee^{(1-s)v}\,\dd v=\sum_{\nu\ge1}\lambdaClass{d}{M}{\nu}\,\nu^{-s}.
\tag{5.3}
\]
To see this, substitute $x=\ee^{v}$: the integral is
\[
\int_1^\infty\lambdaSum{d}{M}{i}{x}\,x^{-s}\,\dd x
=\int_1^\infty\Bigl(\sum_{1\le\nu\le x}\lambdaClass{d}{M}{\nu}\,\nu^{\,i-1}\Bigr)\,x^{-(s+i-1)-1}\,\dd x .
\]
The coefficients
$\lambdaClass{d}{M}{\nu}\nu^{\,i-1}$ are $\ge0$, and their sum up to~$x$ is $x^{i}\lambdaSum{d}{M}{i}{x}\le2x^{1+i}$ by
Lemma~\ref{l:Abound}. Partial summation, in the form that Mathlib provides, says that a Dirichlet series whose
coefficients have absolute partial sums $O(x^{1+i})$ is, at every $s'$ with $\re s'>1+i$, equal to $s'$ times the
integral over $(1,\infty)$ of its partial sums against $x^{-s'-1}$. With $s'=s+i-1$, whose real part exceeds $1+i$
because $\re s>2$, this is~(5.3).

\subhead{Step 5: (P2).} Let $\re s>2$. By Input~\ref{i:laplace}(a) at the point $s-1$,
$\int_0^\infty\ee^{(1-s)v}\dd v=\frac1{s-1}$ and $\int_0^\infty v\,\ee^{(1-s)v}\dd v=\frac1{(s-1)^2}$. By Definition~\ref{d:R},
$\primeError{d}{M}{\ee^{v}}=\lambdaSum{d}{M}{0}{\ee^{v}}-\lambdaSum{d}{M}{1}{\ee^{v}}-v/\varphi(M)-\kap{d}{M}$; each of
the four terms, times $\ee^{(1-s)v}$, is integrable on $(0,\infty)$. So by (5.3) for $i=0$ and $i=1$, and by (5.2),
\begin{align*}
\int_0^\infty\primeError{d}{M}{\ee^{v}}\,\ee^{(1-s)v}\,\dd v
&=\Bigl(\frac1{s-1}-\frac1s\Bigr)\Bigl(\logDerivMean{d}{M}{s}+\frac{s}{\varphi(M)(s-1)}\Bigr)\\
&\qquad-\frac{1}{\varphi(M)(s-1)^2}-\frac{\kap{d}{M}}{s-1}\\
&=\frac{\logDerivMean{d}{M}{s}}{s(s-1)}-\frac{\kap{d}{M}}{s-1}=\Ph{d}{M}{s}.
\end{align*}
The term $\log x/\varphi(M)$ of the remainder has cancelled the double pole at $s=1$ that the pole of the
Dirichlet series produces.

\subhead{Step 6: the constant is real, and (P3).} We claim that $\logDerivMean{d}{M}{1}=\kap{d}{M}$, that is, that
$\logDerivMean{d}{M}{1}$ is real. For real $s>2$ the number $\Ph{d}{M}{s}$ is the integral of a real function, by
Step~5; so $\logDerivMean{d}{M}{s}/s-\kap{d}{M}=(s-1)\Ph{d}{M}{s}$ is real. By Step~1 the function
$s\mapsto\im\bigl(\logDerivMean{d}{M}{s}/s-\kap{d}{M}\bigr)$ of the real variable~$s$ is real analytic on a
neighbourhood of the half-line
$[1,\infty)$; it vanishes on $(2,\infty)$, hence on $[1,\infty)$, by the identity theorem. At $s=1$ this gives
$\im\logDerivMean{d}{M}{1}=0$, and $\re\logDerivMean{d}{M}{1}=\kap{d}{M}$ by definition. Now (P3): for $s\ne1$,
$(s-1)\Ph{d}{M}{s}=\logDerivMean{d}{M}{s}/s-\kap{d}{M}$, and at $s=1$ the left side is $0$ and the right side is
$\logDerivMean{d}{M}{1}-\kap{d}{M}=0$. So $s\mapsto(s-1)\Ph{d}{M}{s}$ is the function of Step~1, which is analytic at
every point of $\re s\ge1$. This is the role of the constant $\kap{d}{M}$ in the remainder: it removes the pole
of the transform at $s=1$.

\subhead{Step 7: (P4).} Let $\re\rho>0$. If $\rho=1$, then $(s-1)\Ph{d}{M}{s}\to0=\res{d}{M}{1}$ by Step~6. Let
$\rho\ne1$. By Input~\ref{i:logderiv}(b), applied to each $\xi$ and to each $\quadTwist{\xi}{d}$,
\[
(s-\rho)\,\logDerivMean{d}{M}{s}\ \longrightarrow\ \frac1{\varphi(M)}\sum_\xi\bigl(\ord_\rho\entireL{\quadTwist{\xi}{d}}-\ord_\rho\entireL{\xi}\bigr)
\qquad(s\to\rho,\ s\ne\rho),
\]
the constant $-1$ in $\logDerivMean{d}{M}{s}$ contributing nothing. Since $s\to\rho\ne0$, $s-1\to\rho-1\ne0$ and
$(s-\rho)\kap{d}{M}\to0$, (5.1) gives
\[
(s-\rho)\,\Ph{d}{M}{s}\ \longrightarrow\ \frac{1}{\varphi(M)\,\rho(\rho-1)}\sum_\xi\bigl(\ord_\rho\entireL{\quadTwist{\xi}{d}}-\ord_\rho\entireL{\xi}\bigr)=\res{d}{M}{\rho}.
\]
This proves Theorem~1. \hfill\qedsymbol
\leanline{\lean{isPolarprimeError}}

\sect{The one-sided criterion: proof of Theorem 2}{s:criterion}
Let $M$, $c$, $\beta$ be as in Theorem~2, let $\rho\in\C$ with $\re\rho>\beta$, and let $\chi$ be a Dirichlet character
modulo~$M$. If $\re\rho\ge1$, then $\Ltv{\rho}{\chi}\ne0$ by Input~\ref{i:nonvanishing}. Assume $\re\rho<1$.

\subhead{Step 1: the prime $d$.} The hypotheses of Input~\ref{i:density} hold for $M$, $\beta$, $c$ and
$z=|\im\rho|$: they are part of the hypotheses of Theorem~2. So there is a prime $d\equiv c\pmod M$ with
\[
\Ltv{s}{\quadTwist{\xi}{d}}\ne0\qquad\text{for all }\xi\text{ modulo }M\text{ and all }s\text{ with }\re s\ge\beta,\ |\im s|\le|\im\rho|.
\tag{6.1}
\]
By hypothesis this $d$ satisfies $d\equiv7\pmod 8$, $\gcd(M,d)=1$, and
$\primeError{d}{M}{x}\le x^{\beta-1}$ for all sufficiently large~$x$.

\subhead{Step 2: the polar function.} By Theorem~1,
\polarWords{v\mapsto\primeError{d}{M}{\ee^{v}}}{s\mapsto\res{d}{M}{s}}. We check the two hypotheses of
Input~\ref{i:landau} for it, with the exponent $\beta>\tfrac12>0$.

\emph{The one-sided bound.} If $v$ is large, so is $\ee^{v}$, and then
$\primeError{d}{M}{\ee^{v}}\le(\ee^{v})^{\beta-1}=\ee^{(\beta-1)v}$.

\emph{The sign of the residues at real points.} Let $s$ be a real number with $s>\beta$. If $s\ge1$, no
$\entireL{\xi}$ and
no $\entireL{\quadTwist{\xi}{d}}$ vanishes at~$s$ (Input~\ref{i:nonvanishing}), all orders at $s$ are $0$
(Input~\ref{i:logderiv}(b)), and $\res{d}{M}{s}=0$. If $\beta<s<1$, the point $s$ lies in the box of (6.1), so
$\ord_s\entireL{\quadTwist{\xi}{d}}=0$ for all $\xi$, and
\[\res{d}{M}{s}=\frac{1}{\varphi(M)\,s(1-s)}\sum_\xi\ord_s\entireL{\xi}\ \ge\ 0,\]
because $s(1-s)>0$.

\subhead{Step 3: conclusion.} Input~\ref{i:landau} gives $\res{d}{M}{\rho}=0$. The point $\rho$ lies in the box of
(6.1), so again $\ord_\rho\entireL{\quadTwist{\xi}{d}}=0$ for all $\xi$; and $\rho(1-\rho)\ne0$ because
$0<\re\rho<1$. Hence $\sum_\xi\ord_\rho\entireL{\xi}=0$. The summands are natural numbers, so each is $0$; in
particular $\ord_\rho\entireL{\chi}=0$, that is, $\Ltv{\rho}{\chi}\ne0$ by Input~\ref{i:logderiv}(b). \hfill\qedsymbol
\leanline{\lean{entireLnezeroofprimeErrorle}}

\begin{remark}[the anatomy of the criterion]\label{r:anatomy}
\emph{Where positivity comes from.} Landau's lemma needs a transform whose poles to the right of~$\beta$, if real,
have residues of one sign. In $\res{d}{M}{s}$ the zeros of the $\entireL{\xi}$ count positively and those of the
twists $\entireL{\quadTwist{\xi}{d}}$ negatively. The twists are present because the weight
$\lambdaClass{d}{M}{\nu}$ carries the factor $1-\chiK{d}(\nu)$; this factor is what the geometric part of the proof
can bound, since it selects the primes that do not split in~$K=\Q(\sqrt{-d})$. Input~\ref{i:density} removes the
negative terms from a box of the height of~$\rho$.

\emph{Why every prime of a class.} The prime $d$ is not chosen in advance: Input~\ref{i:density} produces it, and
it depends on the height $|\im\rho|$ of the point at which a zero is to be excluded. So the one-sided bound must be
known for every prime of the class~$c$.

\emph{Why an upper bound suffices.} The lemma is applied to the non-negative function
$C_1\ee^{\beta v}-\primeError{d}{M}{\ee^{v}}\ee^{v}$ (Remark~\ref{r:landau}); only the inequality
$\primeError{d}{M}{x}\le x^{\beta-1}$ is used, and no lower bound.

\emph{Real zeros.} Zeros of the $\entireL{\xi}$ on the real segment $(\beta,1)$ are not excluded by hypothesis: they
would give poles with positive residue, which Input~\ref{i:landau} admits, and the conclusion excludes them. Real
zeros of the twists in $(\beta,1)$ are excluded by (6.1). The header of the formal file notes that the classical
forms of this argument assume that there are no real zeros, and names as related work M.~Suzuki, \emph{On variants
of Chebyshev's conjecture} (2024), §§15.1--15.2 of [MV], and Landau's paper of 1905 on a theorem of Chebyshev; we
have not compared the statements.
\end{remark}

\sect{A good point in every window: proof of Theorem 3}{s:window}
The proof of Theorem~3 uses the analyticity of the transform on $\re s\ge1$ through the moments
$\int_0^\infty v^{i}\ee^{-3v}\cE(v)\,\dd v$, which are the derivatives of $\Phi$ at $s=4$ up to sign.

\begin{proposition}[moments of a polar function]\label{p:moments}
Let $\cE\colon\RR\to\RR$ and $\omega\colon\C\to\C$ be functions such that \polarWords{\cE}{\omega}. Then there is a
real number $\rr_0>3$ such that for every real number $\rr$ with $3<\rr<\rr_0$ there is a real number
$v_\ast$ such that for every real $v_0\ge v_\ast$:
\begin{enumerate}[label=(\roman*)]
\item $|\cE(v)|\le v_0\,\ee^{v}$ for every real $v\ge0$;
\item for every natural number $i$ the function $v\mapsto v^{i}\ee^{-3v}\cE(v)$ is integrable on $(0,\infty)$, and
\[\Bigl|\int_0^\infty v^{i}\,\ee^{-3v}\,\cE(v)\,\dd v\Bigr|\le\frac{i!\;v_0}{\rr^{i}}.\]
\end{enumerate}
\end{proposition}
\begin{proof}
Let $C$ and $\Phi$ be as in Definition~\ref{d:polar}. Only the measurability of $\cE$ and (P1), (P2), (P3) are used.

\emph{A disc of analyticity.} The closed disc $|s-4|\le3$ lies in the half-plane $\re s\ge1$, and it contains the
point $s=1$. By (P3) the function $s\mapsto(s-1)\Phi(s)$, which vanishes at $s=1$, is analytic at each of its
points. The set of points at which a function is analytic is open
and the disc is compact, so there is a real number $\rr_0>3$ such that $s\mapsto(s-1)\Phi(s)$ is analytic at every
point of the open disc $|s-4|<\rr_0$. Let $3<\rr<\rr_0$. Input~\ref{i:taylor}, with $s_0=4$, $s_1=1$, $s_2=4+\rr$
and the radius $\rr_0$, shows that the Taylor series of $\Phi$ at $4$ converges absolutely at $4+\rr$,
beyond the point $1$:
\[C_{\rr}=\sum_{i\ge0}\frac{\rr^{i}\,|\Phi^{(i)}(4)|}{i!}<\infty.\]

\emph{Differentiation under the integral sign.} We claim that for every natural number $i$ and every $s\in\C$ with
$\re s>2$ the function $v\mapsto(-v)^{i}\cE(v)\ee^{(1-s)v}$ is integrable on $(0,\infty)$, and
\[\Phi^{(i)}(s)=\int_0^\infty(-v)^{i}\,\cE(v)\,\ee^{(1-s)v}\,\dd v.\tag{7.1}\]
For $i=0$ this is Input~\ref{i:laplace}(b) and (P2). Let the claim hold for~$i$, let $\re s>2$, and choose
$\Delta>0$ with $2\Delta<\re s-2$. For $v>0$ and $s'\in\C$ with $\re s'>2+2\Delta$,
\[
\bigl|(-v)^{i+1}\cE(v)\,\ee^{(1-s')v}\bigr|\le\frac1{\Delta}\,\bigl|(-v)^{i}\cE(v)\,\ee^{(1-(2+\Delta))v}\bigr|,
\]
because $\Delta v\le\ee^{\Delta v}$ and $\Delta+1-\re s'\le-1-\Delta$. The right side is
integrable on $(0,\infty)$, by the claim for $i$ at the point $2+\Delta$, and it does not depend on~$s'$. So
the integral in (7.1) may be differentiated under the integral sign on the half-plane $\re s'>2+2\Delta$,
which is a neighbourhood of~$s$: its derivative at $s$ is the integral for $i+1$, which converges absolutely. Since
the integral for $i$ equals $\Phi^{(i)}$ on $\re s'>2$, this is the claim for $i+1$.

\emph{Conclusion.} At $s=4$, (7.1) says that $v\mapsto v^{i}\ee^{-3v}\cE(v)$ is integrable on $(0,\infty)$ and that
$\bigl|\int_0^\infty v^{i}\ee^{-3v}\cE(v)\,\dd v\bigr|=|\Phi^{(i)}(4)|\le i!\,C_{\rr}/\rr^{i}$. So (i) and (ii) hold for
every $v_0\ge v_\ast=\max(C,C_{\rr})$, (i) by (P1).
\end{proof}
\leanline{\lean{eventuallymomentle}}

\subhead{Proof of Theorem 3.} Apply Proposition~\ref{p:moments} to $(\cE_1,\omega_1)$ and to $(\cE_2,\omega_2)$, and
choose a real number $\rr$ with $3<\rr<4$ that is below the bound $\rr_0$ of each of the two. Let $v_\ast$ be so large that the conclusions
(i) and (ii) hold for $\cE_1$ and for $\cE_2$ whenever $v_0\ge v_\ast$, and that
$v_\ast\ge 99/\log(\rr/3)$ and $v_\ast\ge9$. Let $v_0\ge v_\ast$ and $0\le X\le\ee^{v_0}$, and put
$i=\lceil v_0^{2}/9\rceil$. Since $v_0+1\le\ee^{v_0}$ and $X\le\ee^{v_0}$,
\[v_0+v_0X+1\le\ee^{2v_0}.\tag{7.2}\]
Suppose, for a contradiction, that $\cE_1(v)-X\cE_2(v)>\ee^{-v_0}$ for every $v$ in the window $[v_0,v_0^{2}]$.
Consider, for $v>0$,
\[
I(v)=\ee^{7v_0-\rr v}+\ee^{2v_0-v_0^{2}-v}+\ee^{-3v}\bigl(\cE_1(v)-X\,\cE_2(v)-\ee^{-v_0}\bigr).
\]

\emph{$I(v)\ge0$ for all $v>0$.} In the window all three terms are positive. Outside it, by (i) for $\cE_1$ and
$\cE_2$ and by (7.2), $|\cE_1(v)-X\cE_2(v)-\ee^{-v_0}|\le(v_0+Xv_0+1)\ee^{v}\le\ee^{2v_0+v}$, so the third term is
at least $-\ee^{2v_0-2v}$. If $v<v_0$, then $7v_0-\rr v\ge2v_0-2v$, because $(\rr-2)v<2v_0$; so the first term
compensates. If $v>v_0^{2}$, then $2v_0-v_0^{2}-v\ge2v_0-2v$; so the second term compensates.

\emph{The contradiction.} Multiply by $v^{i}$ and integrate over $(0,\infty)$. By Input~\ref{i:laplace}(a), and by
(ii) for $\cE_1$ and $\cE_2$,
\[
0\le\int_0^\infty v^{i}I(v)\,\dd v\le\frac{i!\,\ee^{7v_0}}{\rr^{i+1}}+i!\,\ee^{2v_0-v_0^{2}}+\frac{i!\,(v_0+Xv_0)}{\rr^{i}}-\frac{i!\,\ee^{-v_0}}{3^{i+1}} .
\]
With (7.2) and after division by $i!$,
\[
\frac{\ee^{-v_0}}{3^{i+1}}<\frac{\ee^{7v_0}}{\rr^{i+1}}+\ee^{2v_0-v_0^{2}}+\frac{\ee^{2v_0}}{\rr^{i}}.\tag{7.3}
\]
But each term on the right is at most a third of the left side. Indeed $i\ge v_0^{2}/9$ and $v_0\log(\rr/3)\ge99$
give $(\rr/3)^{i}\ge\ee^{11v_0}$, and $\ee^{11v_0}\ge3\ee^{8v_0}$ and $\ee^{11v_0}\ge9\ee^{3v_0}$: this settles the
first and the third term. For the second term one needs $(i+2)\log3+3v_0\le v_0^{2}$; it holds because
$i+2\le v_0^{2}/9+3$, $\log3<\tfrac{11}{10}$ and $v_0\ge9$. So the right side of (7.3) is at most the left side, a
contradiction. Hence some $v_1\in[v_0,v_0^{2}]$ has $\cE_1(v_1)-X\cE_2(v_1)\le\ee^{-v_0}$. \hfill\qedsymbol
\leanline{\lean{existswindowsubmulle}}

\begin{remark}[what the window lemma replaces]\label{r:window}
The weight $v^{i}\ee^{-3v}$ has its mass near $v=i/3\approx v_0^{2}/27$, inside the window, and the factor
$\rr^{-i}$ with $\rr>3$ in Proposition~\ref{p:moments} says that the moments of a polar function are smaller, by a
factor exponential in~$i$, than those of a function of constant sign and size~$1$. So a polar function cannot stay
above $\ee^{-v_0}$ on a whole window. In the proof of the one-sided bound (Chapter~I) the theorem is applied to
$\cE_1(v)=\primeError{d}{N_0}{\ee^{v}}$ and $\cE_2(v)=\primeError{d}{1}{\ee^{v}}$, which are polar by Theorem~1, and
$\ee^{v_1}$ becomes the second scale of the data. A pointwise bound for these remainders, that is, a prime number
theorem in progressions with a remainder term, is thus not needed; the header of the formal file names the work of
Pintz and of Kaczorowski on sign changes of remainder terms as related, and we have not compared it.
\end{remark}

\sect{Part II. The data, and what is taken from other units}{s:data}
From here on the note needs the vocabulary of the geometric part of the proof: the data, the rows and columns, the
theta lifts. These objects belong to other units, and the chapters on those units define them. This section and the
Settings of §§\ref{s:AB} and~\ref{s:C} say what is used of them. The results of the present unit are then proved in
full from the properties that are listed; the properties themselves are \emph{taken for granted}, each with the unit
that owns it, in the words of the reader's dossier, whose statements were compared with the Lean text.

\begin{setting}[parameters and data; unit~152]\label{st:data}
\emph{Parameters} are a pair $P=(N_0,d)$: a positive multiple $N_0$ of $18480$ and a prime number $d$ with
$d\equiv N_0/2-1\pmod{N_0}$. \emph{Admissible data} for $P$ are a record~$D$ of unit~152. Among its components are
three positive real numbers $\scx{D}$, $\scxo{D}$, $\scb{D}$, the \emph{scales} (the record stores their logarithms), two natural
numbers $\Dm{D}$ and $\Dr{D}$, and a finite set $\thick{D}$ of infinite places of the field $E(D)$ defined below. To the
parameters $P$ and data $D$ the following objects are attached.

\emph{The field and its places.} $K=\Q(\sqrt{-d})$. The natural number $N_1(D)$ is a multiple of $N_0$ with the same
prime divisors. The field
$E(D)$ is $K$ with the roots of unity of order $N_1(D)$ adjoined. Its infinite places form a principal homogeneous
space under $\Gal(E(D)/K)\cong(\Z/N_1(D))^{\times}$; the isomorphism is restriction to the roots of unity, and it is
onto because $K$ is not contained in the field that they generate over~$\Q$: the prime $d$ ramifies in~$K$ and does
not divide $N_1(D)$. A base place, fixed once by the artifact, identifies the places with this group, and so makes
them a finite commutative group; $f(D)$ is their number. Reduction modulo~$N_0$ gives each place $\tau$ a
\emph{class} in $(\Z/N_0)^{\times}$, multiplicatively in~$\tau$. For an integer $\nu$ and a place $\tau$, the place
$\nu\cdot\tau$ is the product of $\tau$ and the place identified with $\nu$ modulo $N_1(D)$ if $\nu$ is prime to
$N_1(D)$, and $\nu\cdot\tau=\tau$ otherwise; if $\nu$ is prime to $N_1(D)$, the class of $\nu\cdot\tau$ is $\nu$ times
the class of~$\tau$.

\emph{Tops and thick places.} $\tops{D}$ is the set of the places of class $1$ modulo~$N_0$. Every place of
$\thick{D}$ has class $-1$ modulo~$N_0$. And $f(D)=\varphi(N_0)\cdot\#\tops{D}$.

\emph{Sizes.} $k(D)=\Dr{D}+1-\Dm{D}$, $t(D)=5\Dr{D}$, $a(D)$ is the least multiple of $N_1(D)$ that is at least $\scx{D}\,\scb{D}$, and
$\budget{D}=f(D)\,a(D)\,(\log \scb{D})^{3}$.

\emph{Prime sums.} The \emph{nominal weight} of a place $\tau$ is $\wD{D}{\tau}=a(D)/\scx{D}$ if $\tau\in\thick{D}$ and
$\wD{D}{\tau}=a(D)/\scxo{D}$ otherwise; it is positive, and $\wbar{D}$ is the mean of $\log\wD{D}{\tau}$ over all $f(D)$
infinite places~$\tau$. For a real number $z$ and a place~$\tau$,
\[
\Ask{D}{z}{\tau}=\sum_{\nu=1}^{a(D)}\lambdaClass{d}{1}{\nu}\,\Bigl(\frac z\nu-\wD{D}{(-\nu)\cdot\tau}\Bigr)_{+}.
\]

\emph{Periods, the cost of one direction, the level.} $g(D)$ is the product of two primes of the data, and
the \emph{period term} is $\GamD{D}=\log\bigl(g(D)\,\mathrm{rd}_{E(D)}/2\pi\bigr)$, with $\mathrm{rd}_{E(D)}$ the root
discriminant of $E(D)$. For a
natural number $i$ the \emph{jet size} is $\jet{D}{i}=\log(i!)-i\,\bigl(\chowla{d}-\GamD{D}\bigr)$. The \emph{cost of
one direction} $\cost{D}$ is a real number: a sum over all infinite places $\tau$ of a jet size minus the prime sum
$\Ask{D}{\wD{D}{\tau}}{\tau}$; its formula is in Chapter~III, and its value plays no role here. The \emph{level}
$N(D)$ is a prime number (unit~141). Finally (unit~145)
\[
\begin{aligned}
\hrow{D}={}&\Dm{D}\Bigl(\sum_{\tau\in\tops{D}}\bigl(\jet{D}{a(D)}-\Ask{D}{a(D)}{\tau}\bigr)+k(D)\,\cost{D}\Bigr)\\
&-\#\tops{D}\cdot a(D)\,\log N(D).
\end{aligned}
\]
\end{setting}

The formal text writes $\chowla{d}-\GamD{D}$ as $\log(\Omega/g(D))$ with the period
$\Omega=\ee^{\chowla{d}}\cdot2\pi/\mathrm{rd}_{E(D)}$. The next two boxes are what Part~II takes from units 152 and~141
about sizes.

\statement{ana}{\inputhead{comparison of sizes}{unit~152}\label{i:sizes}
In Setting~\ref{st:data}:
\begin{enumerate}[label=(S\arabic*),start=0]
\item $a(D)\ge1$, $\scx{D}\le a(D)$, $\scxo{D}\le a(D)$, $\Dm{D}\,k(D)\ne0$, $\scb{D}>0$ and $\budget{D}>0$;
\item $\#\tops{D}\cdot\bigl(\Dm{D}\,(1+\tfrac12\log a(D))+a(D)\,(\GamD{D}+\tfrac{15}{2}\log \scb{D})\bigr)\le\budget{D}$;
\item if $|\Ask{D}{z}{\tau}|\le2\,a(D)\,|z|$ for all real $z$ and all places $\tau$, then
$|\mainHeight{D}|+\Dm{D}\,|\remainder{D}|\le \scb{D}^{2}\scx{D}^{99}$ (Definitions \ref{d:hmain} and~\ref{d:rem});
\item $\scx{D}^{99}+5\,\budget{D}\le\log \scb{D}\cdot\budget{D}$;
\item $\scx{D}^{99}+\scxo{D}^{2}\scb{D}^{2}+4\Dm{D}^{2}\,t(D)\,\scx{D}^{999}\le\tfrac14\,\scb{D}^{3}$;
\item $\Dm{D}^{2}\cdot2\,t(D)\,\scx{D}^{999}\le2\log \scb{D}\cdot\budget{D}$;
\item $\scb{D}\le\log \scb{D}\cdot\budget{D}$.
\end{enumerate}
\src{These are instances of one mechanism of unit~152. Every quantity attached to the data carries a bound of the
form $\exp$ of a linear form in $\log \scx{D}$, $\log \scxo{D}$, $\log \scb{D}$ and $\log\log \scb{D}$, found from the bounds of its factors,
and two such bounds are compared through the conditions on the scales, among them $(\log \scb{D})^{100}\le \scx{D}$,
$\scx{D}^{10^{6}}\le \scxo{D}$ and $\scxo{D}^{10^{6}}\le \scb{D}$. At (S2), (S3) and (S4) the source does not write out the inequality that
it passes to this mechanism; the forms given here are those of the dossier, and the proofs below show that they are
what is needed.}}

\statement{ana}{\inputhead{the level}{unit~141}\label{i:level}
In Setting~\ref{st:data}: $N(D)$ is a prime number, and
$\log N(D)\ge \Dm{D}\,\bigl(\wbar{D}+\GamD{D}\bigr)-\GamD{D}-\tfrac{15}{2}\log \scb{D}$.
\src{The level is chosen, in a window of length $\sqrt{\log \scb{D}}$, above this target.}}

\sect{The remainder of the data: Theorem 4}{s:rem}

\begin{definition}[the main amount]\label{d:hmain}
Let $P$ be parameters and $D$ admissible data for~$P$, with components $\scx{D}$, $\scxo{D}$, $\scb{D}$, $\Dm{D}$, $\Dr{D}$
(Setting~\ref{st:data}). Then $\mainHeight{D}=\Dm{D}\,k(D)\,\cost{D}$.
\end{definition}
\leanline{\lean{mainHeight}}

\begin{definition}[the expected value of a prime sum]\label{d:askmain}
Let $P$ be parameters and $D$ admissible data for~$P$, with components $\scx{D}$, $\scxo{D}$, $\scb{D}$, $\Dm{D}$, $\Dr{D}$
(Setting~\ref{st:data}). Then
\[\mainAsk{D}=a(D)\,\bigl(\log a(D)-1-\chowla{d}-\wbar{D}\bigr)=a(D)\,\bigl(\log a(D)+\kap{d}{1}-\wbar{D}\bigr),\]
the second form by Lemma~\ref{l:constant}.
\end{definition}
\leanline{\lean{mainAsk}}

\begin{definition}[the remainder of the data]\label{d:rem}
Let $P$ be parameters and $D$ admissible data for~$P$, with components $\scx{D}$, $\scxo{D}$, $\scb{D}$, $\Dm{D}$, $\Dr{D}$
(Setting~\ref{st:data}). Then
\[\remainder{D}=\sum_{\tau\in\tops{D}}\bigl(\Ask{D}{a(D)}{\tau}-\mainAsk{D}\bigr).\]
\end{definition}
\leanline{\lean{remainder}}

\statement{arith}{\phantomsection\label{t:4}\textbf{Theorem 4} (the remainder of the data in terms of the remainders
of Part~I)\textbf{.}\enspace Let $P=(N_0,d)$ be parameters and $D$ admissible data for~$P$, with components $\scx{D}$,
$\scxo{D}$, $\scb{D}$, $\Dm{D}$, $\Dr{D}$ (Setting~\ref{st:data}). Then
\[
\begin{aligned}
\remainder{D}=a(D)\,\Bigl(&\#\thick{D}\cdot\bigl(\primeError{d}{N_0}{\scx{D}}-\primeError{d}{N_0}{\scxo{D}}\bigr)\\
&+\#\tops{D}\cdot\primeError{d}{1}{\scxo{D}}\Bigr).
\end{aligned}
\]}

\begin{proof}
Write $a=a(D)$. By (S0), $a\ge1$, $0<\scx{D}\le a$ and $0<\scxo{D}\le a$.

\emph{Step 1: counting.} For every natural number $\nu$,
\[
\#\{\tau\in\tops{D}:(-\nu)\cdot\tau\in\thick{D}\}=
\begin{cases}\#\thick{D}&\text{if }\nu\equiv1\pmod{N_0},\\ 0&\text{otherwise.}\end{cases}
\tag{9.1}
\]
Let $\nu\equiv1\pmod{N_0}$. Then $\nu$ is prime to $N_0$, hence to $N_1(D)$, which has the same prime divisors. So
$\tau\mapsto(-\nu)\cdot\tau$ is a bijection of the set of places that multiplies the classes by $-\nu\equiv-1$: it maps
$\tops{D}$ onto the set of the places of class $-1$ modulo~$N_0$, which contains $\thick{D}$, and every thick place
is met exactly once. Let $\nu\not\equiv1\pmod{N_0}$. If $\nu$ is prime to $N_1(D)$, the places $(-\nu)\cdot\tau$ with
$\tau\in\tops{D}$ have class $-\nu\ne-1$, and none is thick. If $\nu$ is not prime to $N_1(D)$, then
$(-\nu)\cdot\tau=\tau$ by the convention of Setting~\ref{st:data}, and no place of $\tops{D}$ is thick, because the
classes $1$ and $-1$ modulo $N_0$ differ. This proves (9.1). By Definition~\ref{d:lambda},
$\lambdaClass{d}{N_0}{\nu}$ equals $\lambdaClass{d}{1}{\nu}$ if $\nu\equiv1\pmod{N_0}$ and is $0$ otherwise; so (9.1)
gives, for every natural number $\nu$,
\[
\#\{\tau\in\tops{D}:(-\nu)\cdot\tau\in\thick{D}\}\cdot\lambdaClass{d}{1}{\nu}=\#\thick{D}\cdot\lambdaClass{d}{N_0}{\nu}.
\tag{9.2}
\]

\emph{Step 2: the sum of the prime sums over the tops.} For $\tau\in\tops{D}$ and $1\le\nu\le a$ the weight
$\wD{D}{(-\nu)\cdot\tau}$ is $a/\scx{D}$ if $(-\nu)\cdot\tau$ is thick and $a/\scxo{D}$ otherwise. Summing over $\tau$ and using
(9.2),
\begin{align*}
\sum_{\tau\in\tops{D}}\Ask{D}{a}{\tau}
={}&\#\thick{D}\sum_{\nu=1}^{a}\lambdaClass{d}{N_0}{\nu}\Bigl(\bigl(\tfrac a\nu-\tfrac ax\bigr)_{+}-\bigl(\tfrac a\nu-\tfrac a{\scxo{D}}\bigr)_{+}\Bigr)\\
&+\#\tops{D}\sum_{\nu=1}^{a}\lambdaClass{d}{1}{\nu}\bigl(\tfrac a\nu-\tfrac a{\scxo{D}}\bigr)_{+}.
\end{align*}
By Lemma~\ref{l:Rsum}, for $M\in\{N_0,1\}$ and a real number $z$ with $0<z\le a$,
$\sum_{\nu=1}^{a}\lambdaClass{d}{M}{\nu}(\frac a\nu-\frac az)_{+}=a\bigl(\primeError{d}{M}{z}+\log z/\varphi(M)+\kap{d}{M}\bigr)$.
We use it for $M=N_0$ at $z=\scx{D}$ and $z=\scxo{D}$, where the constant $\kap{d}{N_0}$ drops out of the difference, and
for $M=1$ at $z=\scxo{D}$, where $\varphi(1)=1$:
\begin{align*}
\sum_{\tau\in\tops{D}}\Ask{D}{a}{\tau}
={}&a\Bigl(\#\thick{D}\bigl(\primeError{d}{N_0}{\scx{D}}-\primeError{d}{N_0}{\scxo{D}}\bigr)\\
&\qquad{}+\#\tops{D}\,\primeError{d}{1}{\scxo{D}}\Bigr)\\
&+a\Bigl(\#\thick{D}\,\frac{\log \scx{D}-\log \scxo{D}}{\varphi(N_0)}\\
&\qquad{}+\#\tops{D}\,\bigl(\log \scxo{D}+\kap{d}{1}\bigr)\Bigr).
\end{align*}

\emph{Step 3: the expected value.} Since $\#\thick{D}$ places have weight $a/\scx{D}$ and the other $f(D)-\#\thick{D}$
have weight $a/\scxo{D}$, $\log a-\wbar{D}=\log \scxo{D}+\frac{\#\thick{D}}{f(D)}(\log \scx{D}-\log \scxo{D})$. With
$f(D)=\varphi(N_0)\cdot\#\tops{D}$ and Definition~\ref{d:askmain},
\[
\begin{aligned}
\#\tops{D}\cdot\mainAsk{D}=a\Bigl(&\#\thick{D}\,\frac{\log \scx{D}-\log \scxo{D}}{\varphi(N_0)}\\
&+\#\tops{D}\,\bigl(\log \scxo{D}+\kap{d}{1}\bigr)\Bigr).
\end{aligned}
\]

\emph{Step 4.} By Definition~\ref{d:rem}, $\remainder{D}=\sum_{\tau\in\tops{D}}\Ask{D}{a}{\tau}-\#\tops{D}\cdot\mainAsk{D}$,
and Steps 2 and~3 give the claim: the terms $\log \scx{D}$, $\log \scxo{D}$ and $\kap{d}{1}$ cancel.
\end{proof}
\leanline{\lean{remaindereq}}

Step~1 is the place where the progression $1$ modulo $N_0$ enters the proof: the thick places sit in one class,
and the tops are the stabiliser of that class. The convention for the integers $\nu$ that are not prime to $N_1(D)$
does not affect the prime sums: for such $\nu$ the weight $\lambdaClass{d}{1}{\nu}$ is $0$, because every prime
that divides $N_1(D)$ splits in~$K$ (unit~152).

\begin{proposition}[the level pays the excess]\label{p:level}
Let $P$ be parameters and $D$ admissible data for~$P$, with components $\scx{D}$, $\scxo{D}$, $\scb{D}$, $\Dm{D}$, $\Dr{D}$
(Setting~\ref{st:data}). Then
\[\hrow{D}\le\mainHeight{D}-\Dm{D}\cdot\remainder{D}+\budget{D}.\]
\end{proposition}
\begin{proof}
Write $a=a(D)$. By the definitions of $\hrow{D}$, $\mainHeight{D}$ and $\remainder{D}$,
\[
\begin{aligned}
\hrow{D}=\mainHeight{D}-\Dm{D}\cdot\remainder{D}&\\
{}+\#\tops{D}\cdot\Bigl(\Dm{D}\,\bigl(\jet{D}{a}-\mainAsk{D}\bigr)&-a\,\log N(D)\Bigr).
\end{aligned}\tag{9.3}
\]
\emph{The excess of one row.} By the definitions of the jet size and of $\mainAsk{D}$, in which the constant
$\chowla{d}$ cancels,
$\jet{D}{a}-\mainAsk{D}=\log(a!)-a\log a+a+a\,(\wbar{D}+\GamD{D})$. Stirling's bound
$\log(a!)\le a\log a-a+1+\tfrac12\log a$, valid for $a\ge1$, gives
\[\jet{D}{a}-\mainAsk{D}\le a\,(\wbar{D}+\GamD{D})+1+\tfrac12\log a.\]
\emph{The level.} By Input~\ref{i:level}, $a\log N(D)\ge a\,\Dm{D}\,(\wbar{D}+\GamD{D})-a\,(\GamD{D}+\tfrac{15}{2}\log \scb{D})$.
So the term $\Dm{D}\,a\,(\wbar{D}+\GamD{D})$ cancels in the bracket of (9.3), which is at most
$\Dm{D}\,(1+\tfrac12\log a)+a\,(\GamD{D}+\tfrac{15}{2}\log \scb{D})$. The claim follows from (S1).
\end{proof}
\leanline{\lean{heightPerRowle}}

The number $\hrow{D}$ will bound the logarithm of the size of a chain of $\Dm{D}$ theta lifts of a short column
(Input~\ref{i:size}). The proposition says that the level $N(D)$, whose logarithm is subtracted in $\hrow{D}$, has
been chosen so that it cancels the excess of size over expected divisibility of the $\Dm{D}$ lifts; what remains is the
main amount, less $\Dm{D}$ times the remainder of the data.

\begin{lemma}[a crude bound]\label{l:crude}
Let $P$ be parameters and $D$ admissible data for~$P$, with components $\scx{D}$, $\scxo{D}$, $\scb{D}$, $\Dm{D}$, $\Dr{D}$
(Setting~\ref{st:data}). Then
\[|\mainHeight{D}|+\Dm{D}\cdot|\remainder{D}|\le \scb{D}^{2}\,\scx{D}^{99}.\]
\end{lemma}
\begin{proof}
For every real number $z$ and every place $\tau$, $|\Ask{D}{z}{\tau}|\le2\,a(D)\,|z|$: the sum has $a(D)$ terms, and
$0\le\lambdaClass{d}{1}{\nu}\,(z/\nu-Y)_{+}\le2\nu\cdot|z|/\nu$ for every weight $Y\ge0$, by Lemma~\ref{l:lambdabound}.
The claim is then (S2). In the formal text the comparison of sizes also uses that $|\chowla{d}|$ is at most a
constant of the parameters, which the conditions on the scales make negligible.
\end{proof}
\leanline{\lean{absmainHeightaddremainderle}}

The bound is far from sharp, and it is only used to pay for a small number of missing dimensions in the proof of
the main inequality. To see that there is room (this is not the formal argument): unit~152 has $\Dm{D}\le2\sqrt{\scx{D}}$,
$\Dr{D}\le3\scx{D}^{5}$ and $f(D)\le \scx{D}^{11}(\log \scb{D})^{1/5}$ up to a constant of the parameters, and $a(D)$ is about $\scx{D}\,\scb{D}$; so the
$f(D)$ prime sums in $\cost{D}$, each at most $2a(D)^{2}/\scx{D}$, give $|\mainHeight{D}|$ the order $\scx{D}^{18}\scb{D}^{2}$ at
most, and $\Dm{D}\,|\remainder{D}|$ has the order $\scx{D}^{14}\scb{D}^{2}$ at most.

\sect{Theorems A and B}{s:AB}

\begin{setting}[forms and lifts; units 102, 140, 141, 143, 145, 147, 148]\label{st:forms}
Let $P$ be parameters and $D$ admissible data for~$P$, with components $\scx{D}$, $\scxo{D}$, $\scb{D}$, $\Dm{D}$, $\Dr{D}$, as in
Setting~\ref{st:data}. Other units attach to $D$ the following objects.

\emph{Columns} (unit~145). $\Col{D}$ is a finite-dimensional complex vector space of algebraic modular forms on a
definite unitary group of rank $k(D)$ over $E(D)$; its complex dimension $\dim\Col{D}$ is a different number from the
rank $k(D)$, and both occur below. The \emph{ramifying operator} is a linear map
$\operatorname{ram}_D\colon\Col{D}\to\Col{D}$, and its image $\Colram{D}$ is the space of \emph{ramified columns}.

\emph{Rows} (units 140, 141). $\GD{D}$ is a finite group. $\PF{D}$ is a complex vector space of cusp forms on an
indefinite unitary group of rank~$\Dr{D}$ over $E(D)$, the \emph{forms of principal level}, on which $\GD{D}$ acts
linearly. The space of \emph{rows} $\Row{D}\subseteq\PF{D}$ is the image of a fixed idempotent of the group algebra of
$\GD{D}$; it is finite-dimensional. For $\gamma\in\GD{D}$ the linear map $\gpt{D}{\gamma}\colon\PF{D}\to\Row{D}$
translates a form by $\gamma$ and applies $\#\GD{D}$ times the idempotent.

\emph{Poincaré series} (unit~148). A linear map $\Poin{D}\colon\Col{D}\to\PF{D}$.

\emph{Twists and lifts} (units 141, 143, 147). A \emph{twist} of $D$ is a character $\vartheta$ of a certain kind,
with a finite set $Q(\vartheta)$ of primes of $E(D)$ where it ramifies; an \emph{$\Dm{D}$-fold twist} is a twist whose
powers $\vartheta^{2},\dots,\vartheta^{\Dm{D}}$ are twists with the same set; the trivial character $1$ is an $\Dm{D}$-fold
twist with $Q(1)=\emptyset$. A prime number $\ell$ is \emph{prime to~$\vartheta$} if it is prime to the norms of the
primes in $Q(\vartheta)$. For a twist $\vartheta$, the \emph{chain of lifts} is a linear map
$\Thf{D}{\vartheta}\colon\PF{D}\to\AAf{D}{\vartheta}$, the composite of $\Dm{D}$ theta lifts, into a finite-dimensional
complex vector space $\AAf{D}{\vartheta}$ of algebraic modular forms on a definite unitary group, which carries a
positive definite hermitian product. For an $\Dm{D}$-fold twist $\vartheta$, $\Th{D}{\vartheta}\colon\Row{D}\to\AAf{D}{\vartheta}$
is the restriction of $\Thf{D}{\vartheta}$ to the rows, and
$\zz{D}{\vartheta}=\det\bigl(\Th{D}{\vartheta}^{\ast}\Th{D}{\vartheta}\bigr)$.

\emph{Conjugates and sizes} (unit~102). The columns and the elements of the spaces $\AAf{D}{\vartheta}$ are functions
with values in complex vector spaces that have a rational structure and a norm; so for $\sigma\in\AutC$ a form $F$
has a conjugate $\sigma\circ F$. We say that $F$ is \emph{pointwise at most $\ee^{B}$ at~$\sigma$}, for a real
number~$B$, if every value of $\sigma\circ F$ has squared norm at most $\ee^{B}$. For $F,F_1,\dots,F_n$ in
$\AAf{D}{\vartheta}$, $\nsq{F}{\sigma}$ is the Petersson norm of $\sigma\circ F$, a finite sum over classes of squared
norms of values, and $\covol{F_1,\dots,F_n}{\sigma}$ is the absolute value of the determinant of the matrix of
Petersson products of the $\sigma\circ F_i$.

\emph{Integrality} (unit~102). For a prime number $\ell$ and a ring homomorphism $j\colon\C\to\C_\ell$, $O_j$ is the
ring of the complex numbers $H$ with $|j(H)|_\ell\le1$. A column, or an element of some $\AAf{D}{\vartheta}$, is
\emph{$j$-integral} if all its values are, for a lattice fixed in unit~102; the $j$-integral elements form an
$O_j$-module, and if $F$ is $j$-integral and $|j(H)|_\ell\le1$ for a complex number $H$, then $HF$ is $j$-integral.
A column is \emph{integral} if it is $j$-integral for every prime $\ell$ and every $j\colon\C\to\C_\ell$.
\end{setting}

\statement{ana}{\inputhead{the seesaw identity}{unit~141}\label{i:seesaw}
In Settings \ref{st:data} and~\ref{st:forms} there is a non-zero complex number $H$ such that for every twist
$\vartheta$ of $D$ there is a good datum $\Psi$ for~$\vartheta$ with
\[\Thf{D}{\vartheta}\bigl(\Poin{D}(H\,\psi)\bigr)=\Ch{D}{\vartheta,\Psi}(\psi)\qquad\text{for every }\psi\in\Colram{D}.\]
Here a \emph{good datum} for $\vartheta$ is a notion of unit~145, and
$\Ch{D}{\vartheta,\Psi}\colon\Col{D}\to\AAf{D}{\vartheta}$ is the chain of $\Dm{D}$ theta lifts between definite spaces that
$\Psi$ defines (unit~147). The formal text fixes one such number~$H$, by choice; we write $\cst{see}{D}$ for it.
\src{The identity moves the lifts of Poincaré series to the definite side, where sizes and divisibility are
computed; it is the subject of Chapter~V.}}

\begin{definition}[the scaled inclusion]\label{d:inc}
Let $P$ be parameters and $D$ admissible data for~$P$, with components $\scx{D}$, $\scxo{D}$, $\scb{D}$, $\Dm{D}$, $\Dr{D}$ (Settings
\ref{st:data} and~\ref{st:forms}). The linear map $\scaledInclusion{D}\colon\Colram{D}\to\Col{D}$ is the inclusion
multiplied by the constant $\cst{see}{D}$ of Input~\ref{i:seesaw}.
\end{definition}
\leanline{\lean{scaledInclusion}}

\begin{definition}[the Poincaré maps]\label{d:proj}
Let $P$ be parameters and $D$ admissible data for~$P$, with components $\scx{D}$, $\scxo{D}$, $\scb{D}$, $\Dm{D}$, $\Dr{D}$ (Settings
\ref{st:data} and~\ref{st:forms}), and let $\gamma\in\GD{D}$. The linear map
$\projectedPoincare{D}{\gamma}{}\colon\Colram{D}\to\Row{D}$ is the composite
$\gpt{D}{\gamma}\circ\Poin{D}\circ\scaledInclusion{D}$; its value at a ramified column $\psi$ is written
$\projectedPoincare{D}{\gamma}{\psi}$.
\end{definition}
\leanline{\lean{projectedPoincare}}

\begin{definition}[integral ramified columns]\label{d:intram}
Let $P$ be parameters and $D$ admissible data for~$P$, with components $\scx{D}$, $\scxo{D}$, $\scb{D}$, $\Dm{D}$, $\Dr{D}$ (Settings
\ref{st:data} and~\ref{st:forms}). Then $\integralRamified{D}$ is the set of the ramified columns
$\psi\in\Colram{D}$ that are integral: $j$-integral for every prime number $\ell$ and every ring homomorphism
$j\colon\C\to\C_\ell$.
\end{definition}
\leanline{\lean{integralRamified}}

\begin{definition}[the index of a family of rows]\label{d:index}
Let $P$ be parameters and $D$ admissible data for~$P$, with components $\scx{D}$, $\scxo{D}$, $\scb{D}$, $\Dm{D}$, $\Dr{D}$ (Settings
\ref{st:data} and~\ref{st:forms}). Let $n$ be a natural number, $y=(y_1,\dots,y_n)$ a family of rows in $\Row{D}$,
and $\delta$ a real number. We say that \indexWords{D}{y}{\delta}, and write $\IndexLe{D}{y}{\delta}$, if there are
real numbers $\eta_\ell$, $\ell\in\NN$, with $\sum_{\ell\in J}\eta_\ell\le\delta$ for every finite set $J\subset\NN$,
such that for every prime number $\ell$, every ring homomorphism $j\colon\C\to\C_\ell$ and every $\Dm{D}$-fold twist
$\vartheta$ of $D$ to which $\ell$ is prime, there is an $n\times n$ matrix $u=(u_{ii'})$ with entries in $O_j$
such that $|j(\det u)|_\ell\ge\ee^{-\eta_\ell}$ and such that
$\Th{D}{\vartheta}\bigl(\sum_{i'}u_{ii'}\,y_{i'}\bigr)$ is $j$-integral for $i=1,\dots,n$.
\end{definition}
\leanline{\lean{IndexLe}}

The index bounds, prime by prime and for all twists at once, the denominators of the lifts of the family: after a
change of the family by a matrix of small determinant, the lifts are integral.

\statement{ana}{\inputhead{short integral columns span}{unit~145}\label{i:short}
In Settings \ref{st:data} and~\ref{st:forms}, $\Colram{D}$ is spanned over $\C$ by the ramified columns that are
integral and pointwise at most $\ee^{\budget{D}}$ at every $\sigma\in\AutC$.}

\statement{ana}{\inputhead{the size of the chain}{units 145, 150}\label{i:size}
In Settings \ref{st:data} and~\ref{st:forms}, let $\nonsplit{D}$ be the set of the prime numbers $\ell<\scxo{D}$ with
$\chiK{d}(\ell)\ne1$, and let $\gain{D}{\ell}$, for a natural number~$\ell$, be the real number that unit~150 calls the
\emph{gain} at~$\ell$: the divisibility by $\ell$ that is asked of a chain of lifts. Let $\Psi$ be a good datum for
the trivial twist, $\sigma\in\AutC$, $\psi\in\Colram{D}$ and $B$ a real number. If $\psi$ is pointwise at most
$\ee^{B}$ at~$\sigma$, then $\Ch{D}{1,\Psi}(\psi)$ is pointwise at most
\[\exp\Bigl(B+\hrow{D}+2\sum_{\ell\in\nonsplit{D}}\gain{D}{\ell}\,\log\ell+\budget{D}\Bigr)\quad\text{at }\sigma.\]}

\statement{ana}{\inputhead{passing to the rows}{unit~143}\label{i:goodpart}
In Settings \ref{st:data} and~\ref{st:forms}, let $F\in\PF{D}$ and $\gamma\in\GD{D}$.

\emph{(a)} Let $\sigma\in\AutC$ and let $B$ be a real number. If $\Thf{D}{1}(F)$ is pointwise at most $\ee^{B}$
at~$\sigma$, then $\nsq{\Th{D}{1}(\gpt{D}{\gamma}F)}{\sigma}\le\exp(\scx{D}^{99}+B)$.

\emph{(\scb{D})} Let $\vartheta$ be a twist of $D$, $\ell$ a prime number, $j\colon\C\to\C_\ell$ a ring homomorphism and
$H$ a complex number. If $H\cdot\Thf{D}{\vartheta}(F)$ is $j$-integral, then so is
$H\cdot\Thf{D}{\vartheta}(\gpt{D}{\gamma}F)$.}

\statement{ana}{\inputhead{the divisibility of the chain}{units 145, 150}\label{i:div}
In Settings \ref{st:data} and~\ref{st:forms}, unit~145 defines for every natural number $\ell$ a real number
$\err{D}{\ell}$, the \emph{allowed loss} at~$\ell$, with these properties.

\emph{(a)} $\err{D}{\ell}=0$ unless $\ell\in\nonsplit{D}$ or $\ell=N(D)$; and
$\sum_{\ell\in J}\err{D}{\ell}\log\ell\le\budget{D}$ for every finite set $J\subset\NN$.

\emph{(\scb{D})} Let $\ell$ be a prime number, $j\colon\C\to\C_\ell$ a ring homomorphism, and $e$ an integer with
$e<\gain{D}{\ell}-\err{D}{\ell}+1$ if $\ell\in\nonsplit{D}$, and $e<-\err{D}{\ell}+1$ otherwise. Then there are a
natural number $i_0\le\err{D}{\ell}\cdot\dim\Col{D}$ and an $O_j$-submodule $\mathcal{S}$ of $\Col{D}$ such that:
every finite family $\psi_1,\dots,\psi_n$ of $j$-integral columns admits an $n\times n$ matrix $u$ over $O_j$ whose
determinant divides $\ell^{\,i_0}$ in $O_j$ and for which all $\sum_{i'}u_{ii'}\psi_{i'}$ lie in $\mathcal{S}$; and
$\ell^{-e}\cdot\Ch{D}{\vartheta,\Psi}(\psi)$ is $j$-integral for every twist $\vartheta$ of $D$ to which $\ell$ is
prime, every good datum $\Psi$ for~$\vartheta$ and every $\psi\in\mathcal{S}\cap\Colram{D}$.}

\statement{ana}{\inputhead{the index of a span}{unit~149}\label{i:span}
In Settings \ref{st:data} and~\ref{st:forms}, let $\ell$ be a prime number, $j\colon\C\to\C_\ell$ a ring
homomorphism, $i_0$ a natural number and $q$ a rational number. Let $\mathcal{M}\subseteq\Colram{D}$ and
$\mathcal{N}\subseteq\Row{D}$ be $O_j$-submodules. Assume that for every $\gamma\in\GD{D}$ every finite family
$\psi_1,\dots,\psi_n$ in $\mathcal{M}$ admits an $n\times n$ matrix $u$ over $O_j$ whose determinant divides
$\ell^{\,i_0}$ and for which all $q\,\projectedPoincare{D}{\gamma}{\bigl(\sum_{i'}u_{ii'}\psi_{i'}\bigr)}$ lie in
$\mathcal{N}$. Then for every finite set $\mathcal{T}\subseteq\GD{D}$, every finite family $y_1,\dots,y_n$ in the
$O_j$-span of $\mathcal{N}$ and of the rows $q\,\projectedPoincare{D}{\gamma}{\psi}$ with $\gamma\in\mathcal{T}$,
$\psi\in\mathcal{M}$, admits an $n\times n$ matrix $u$ over $O_j$ whose determinant divides
$\ell^{\,i_0\cdot\#\mathcal{T}}$ and for which all $\sum_{i'}u_{ii'}y_{i'}$ lie in~$\mathcal{N}$.
\src{A special case of a lemma on modules over a commutative ring: the index of a sum of images is at most the
product of the indices.}}

\begin{lemma}[a rational number with prescribed valuations]\label{l:rational}
Let $J$ be a finite set of natural numbers and $e\colon\NN\to\Z$, $\ell\mapsto e_\ell$, a function. Assume that every
element of $J$ is a prime number and that $e_\ell=0$ for every $\ell\notin J$. Then there is a rational number $q$
with $q>0$, $\log q=\sum_{\ell\in J}e_\ell\log\ell$, and $|q|_\ell=|\ell|_\ell^{\,e_\ell}$ for every prime
number~$\ell$, where $|\cdot|_\ell$ is the absolute value of $\C_\ell$.
\end{lemma}
\begin{proof}
Take $q=\prod_{\ell\in J}\ell^{\,e_\ell}$. For a prime $\ell\in J$ the other factors are $\ell$-adic units; for a prime
$\ell\notin J$ all factors are, and $e_\ell=0$.
\end{proof}
\leanline{\lean{existsratlogeqandpadicnormeq}}

\statement{arith}{\phantomsection\label{t:A}\textbf{Theorem A} (short forms)\textbf{.}\enspace Let $P$ be parameters
and $D$ admissible data for~$P$, with components $\scx{D}$, $\scxo{D}$, $\scb{D}$, $\Dm{D}$, $\Dr{D}$ (Settings \ref{st:data}
and~\ref{st:forms}). There is a rational number $q>0$ with the following two properties.

\emph{(a) Size.} There is a subset $S$ of $\integralRamified{D}$ that spans $\Colram{D}$ over~$\C$, such that for all
$\psi\in S$, all $\gamma\in\GD{D}$ and all $\sigma\in\AutC$,
\[
\nsq{\Th{D}{1}\bigl(q\,\projectedPoincare{D}{\gamma}{\psi}\bigr)}{\sigma}\le\exp\bigl(\mainHeight{D}-\Dm{D}\cdot\remainder{D}+\log \scb{D}\cdot\budget{D}\bigr).
\]

\emph{(\scb{D}) Index.} For every finite set $\mathcal{T}\subseteq\GD{D}$, every natural number $n$ and every family
$y=(y_1,\dots,y_n)$ of rows, each of the form $q\,\projectedPoincare{D}{\gamma}{\psi}$ with $\gamma\in\mathcal{T}$ and
$\psi\in\integralRamified{D}$: \indexWords{D}{y}{\#\mathcal{T}\cdot\dim\Col{D}\cdot\budget{D}}.}

\begin{proof}
\emph{The number $q$.} Let $J=\nonsplit{D}\cup\{N(D)\}$, a finite set of primes (Input~\ref{i:level}), and for a
natural number $\ell$ let
\[e_\ell=\bigl\lfloor\err{D}{\ell}-\gain{D}{\ell}\bigr\rfloor\ \text{ if }\ell\in\nonsplit{D},\qquad
e_\ell=\bigl\lfloor\err{D}{\ell}\bigr\rfloor\ \text{ otherwise.}\]
Then $e_\ell=0$ for $\ell\notin J$ by Input~\ref{i:div}(a). Let $q$ be the rational number of Lemma~\ref{l:rational}:
$q>0$, $\log q=\sum_{\ell\in J}e_\ell\log\ell$, and $|q|_\ell=|\ell|_\ell^{\,e_\ell}$ for every prime~$\ell$. For every
twist $\vartheta$ choose a good datum $\Psi_\vartheta$ as in Input~\ref{i:seesaw}. By Definitions \ref{d:inc}
and~\ref{d:proj}, for $\gamma\in\GD{D}$, $\psi\in\Colram{D}$ and an $\Dm{D}$-fold twist~$\vartheta$,
\[
\projectedPoincare{D}{\gamma}{\psi}=\gpt{D}{\gamma}\bigl(\Poin{D}(\cst{see}{D}\,\psi)\bigr),\qquad
\Thf{D}{\vartheta}\bigl(\Poin{D}(\cst{see}{D}\,\psi)\bigr)=\Ch{D}{\vartheta,\Psi_\vartheta}(\psi).\tag{10.1}
\]

\emph{(a).} Let $S$ be the set of the ramified columns that are integral and pointwise at most
$\ee^{\budget{D}}$ at every $\sigma\in\AutC$. Then $S\subseteq\integralRamified{D}$, and $S$ spans $\Colram{D}$ by
Input~\ref{i:short}. Let $\psi\in S$, $\gamma\in\GD{D}$ and $\sigma\in\AutC$. Since $\sigma(q)=q$, the column $q\psi$
is pointwise at most $q^{2}\ee^{\budget{D}}=\exp(2\log q+\budget{D})$ at~$\sigma$. By (10.1) and Input~\ref{i:size},
$\Thf{D}{1}(\Poin{D}(\cst{see}{D}\,q\psi))=\Ch{D}{1,\Psi_1}(q\psi)$ is pointwise at most
$\exp\bigl(2\log q+\budget{D}+\hrow{D}+2\sum_{\ell\in\nonsplit{D}}\gain{D}{\ell}\log\ell+\budget{D}\bigr)$ at~$\sigma$,
and by Input~\ref{i:goodpart}(a) the number $\nsq{\Th{D}{1}(q\,\projectedPoincare{D}{\gamma}{\psi})}{\sigma}$ is at
most $\exp$ of $\scx{D}^{99}$ plus this exponent. Now
\begin{align*}
2\log q=2\sum_{\ell\in J}e_\ell\log\ell
&\le2\sum_{\ell\in J}\err{D}{\ell}\log\ell-2\sum_{\ell\in\nonsplit{D}}\gain{D}{\ell}\log\ell\\
&\le2\,\budget{D}-2\sum_{\ell\in\nonsplit{D}}\gain{D}{\ell}\log\ell,
\end{align*}
by Input~\ref{i:div}(a). The gains cancel: this is what $q$ is for. With
$\hrow{D}\le\mainHeight{D}-\Dm{D}\cdot\remainder{D}+\budget{D}$ (Proposition~\ref{p:level}) the exponent is at most
$\mainHeight{D}-\Dm{D}\cdot\remainder{D}+\scx{D}^{99}+5\,\budget{D}$, and (S3) gives~(a).

\emph{(\scb{D}).} Let $\mathcal{T}$, $n$ and $y$ be as in~(\scb{D}), and put
$\eta_\ell=\#\mathcal{T}\cdot\dim\Col{D}\cdot\err{D}{\ell}\log\ell$ for $\ell\in\NN$. Then
$\sum_{\ell\in J'}\eta_\ell\le\#\mathcal{T}\cdot\dim\Col{D}\cdot\budget{D}$ for every finite set $J'\subset\NN$, by
Input~\ref{i:div}(a). Let $\ell$ be a prime number, $j\colon\C\to\C_\ell$ a ring homomorphism and $\vartheta$ an
$\Dm{D}$-fold twist to which $\ell$ is prime. The integer $e=-e_\ell$ satisfies the hypothesis of Input~\ref{i:div}(\scb{D}),
because $-\lfloor z\rfloor<-z+1$ for real~$z$. Let $i_0$ and $\mathcal{S}$ be as given there, and let $\mathcal{M}$
be the $O_j$-module of the $j$-integral ramified columns, $\mathcal{N}$ the $O_j$-module of the rows $y'$ for which
$\Th{D}{\vartheta}(y')$ is $j$-integral. We check the hypothesis of Input~\ref{i:span}. Let $\gamma\in\GD{D}$ and let
$\psi_1,\dots,\psi_n$ be in $\mathcal{M}$. Input~\ref{i:div}(\scb{D}) gives a matrix $u$ over $O_j$ whose determinant
divides $\ell^{\,i_0}$ such that every $\psi'_i=\sum_{i'}u_{ii'}\psi_{i'}$ lies in $\mathcal{S}\cap\Colram{D}$ and
$\ell^{\,e_\ell}\Ch{D}{\vartheta,\Psi_\vartheta}(\psi'_i)$ is $j$-integral. Since $|q|_\ell=|\ell^{\,e_\ell}|_\ell$, also
$q\,\Ch{D}{\vartheta,\Psi_\vartheta}(\psi'_i)=q\,\Thf{D}{\vartheta}(\Poin{D}(\cst{see}{D}\psi'_i))$ is $j$-integral, by
(10.1); and then so is $q\,\Thf{D}{\vartheta}\bigl(\gpt{D}{\gamma}\Poin{D}(\cst{see}{D}\psi'_i)\bigr)
=\Th{D}{\vartheta}\bigl(q\,\projectedPoincare{D}{\gamma}{\psi'_i}\bigr)$, by Input~\ref{i:goodpart}(\scb{D}). So
$q\,\projectedPoincare{D}{\gamma}{\psi'_i}\in\mathcal{N}$. Each $y_i$ is of the form
$q\,\projectedPoincare{D}{\gamma}{\psi}$ with $\gamma\in\mathcal{T}$ and $\psi\in\integralRamified{D}\subseteq\mathcal{M}$.
So Input~\ref{i:span} gives an $n\times n$ matrix $u$ over $O_j$ whose determinant divides
$\ell^{\,i_0\cdot\#\mathcal{T}}$ in $O_j$, with $\Th{D}{\vartheta}(\sum_{i'}u_{ii'}y_{i'})$ $j$-integral for all~$i$.
Finally $|j(\det u)|_\ell\ge|\ell|_\ell^{\,i_0\cdot\#\mathcal{T}}=\exp(-i_0\cdot\#\mathcal{T}\cdot\log\ell)\ge\ee^{-\eta_\ell}$,
because $i_0\le\err{D}{\ell}\cdot\dim\Col{D}$. This is Definition~\ref{d:index}.
\end{proof}
\leanline{\lean{shortforms}}

\statement{ana}{\inputhead{columns and rows exist}{units 141, 148}\label{i:nonzero}
In Settings \ref{st:data} and~\ref{st:forms}: $\Col{D}\ne0$ and $\Row{D}\ne0$.}

\statement{ana}{\inputhead{the filling theorem, and few translates}{units 148, 135}\label{i:filling}
In Settings \ref{st:data} and~\ref{st:forms}:

\emph{(a)} For every non-zero complex number $H$ the family of linear maps
$\gpt{D}{\gamma}\circ\Poin{D}\circ(H\cdot\operatorname{ram}_D)\colon\Col{D}\to\Row{D}$, $\gamma\in\GD{D}$, is
\emph{spread with constant $\scb{D}^{-2}$}: for every subspace $W$ of $\Row{D}$ there is $\gamma\in\GD{D}$ such that the
preimage of $W$ under the map of index $\gamma$ has dimension at most
$\bigl(\tfrac12(1+\dim W/\dim\Row{D})+\scb{D}^{-2}\bigr)\cdot\dim\Col{D}$.

\emph{(\scb{D})} If a family of linear maps from $\Col{D}\ne0$ to $\Row{D}$, indexed by $\GD{D}$, is spread with constant
$\scb{D}^{-2}$, then there is a finite set $\mathcal{T}\subseteq\GD{D}$ with
$\#\mathcal{T}\cdot\dim\Col{D}\le5\,\dim\Row{D}\cdot\log(1+\scb{D}^{2})$ such that the sum of the images of the maps of
index $\gamma\in\mathcal{T}$ has dimension at least $(1-4\scb{D}^{-2})\dim\Row{D}$.
\src{(a) is the filling theorem of unit~148, the subject of Chapter~VI. (\scb{D}) is a special case of a lemma of
unit~135 on families of linear maps between finite-dimensional spaces, proved by a greedy choice of translates.}}

\statement{arith}{\phantomsection\label{t:B}\textbf{Theorem B} (enough forms)\textbf{.}\enspace Let $P$ be parameters
and $D$ admissible data for~$P$, with components $\scx{D}$, $\scxo{D}$, $\scb{D}$, $\Dm{D}$, $\Dr{D}$ (Settings \ref{st:data}
and~\ref{st:forms}). Then $\dim\Col{D}>0$ and $\dim\Row{D}>0$, and there is a finite set $\mathcal{T}\subseteq\GD{D}$
with
\begin{align*}
\#\mathcal{T}\cdot\dim\Col{D}&\le5\,\dim\Row{D}\cdot\log(1+\scb{D}^{2})\qquad\text{and}\\
\Bigl(1-\frac4{\scb{D}^{2}}\Bigr)\dim\Row{D}&\le\dim\sum_{\gamma\in\mathcal{T}}\projectedPoincare{D}{\gamma}{\bigl(\Colram{D}\bigr)} .
\end{align*}}

\begin{proof}
The dimensions are positive by Input~\ref{i:nonzero}. Apply Input~\ref{i:filling}(a) with $H=\cst{see}{D}$, and then
(\scb{D}), which is allowed because $\scb{D}>0$ by (S0): there is a finite set $\mathcal{T}\subseteq\GD{D}$ with the first
inequality, such that the images of the maps
$\gpt{D}{\gamma}\circ\Poin{D}\circ(\cst{see}{D}\cdot\operatorname{ram}_D)$, $\gamma\in\mathcal{T}$, span dimension at
least $(1-4/\scb{D}^{2})\dim\Row{D}$. For a column $\psi_0$ the element $\operatorname{ram}_D(\psi_0)$ is a ramified column,
and $\gpt{D}{\gamma}\Poin{D}(\cst{see}{D}\operatorname{ram}_D(\psi_0))=\projectedPoincare{D}{\gamma}{\operatorname{ram}_D(\psi_0)}$
by Definitions \ref{d:inc} and~\ref{d:proj}. So each of these images lies in
$\projectedPoincare{D}{\gamma}{(\Colram{D})}$, and the second inequality follows.
\end{proof}
\leanline{\lean{enoughforms}}

\sect{Theorem C, and the family given by A and B}{s:C}

\begin{setting}[heights and fit families; units 135, 138]\label{st:heights}
Let $P$ and $D$ be as in Settings \ref{st:data} and~\ref{st:forms}. Families of complex numbers and of twists are
indexed by the pairs $(\ell,\iota)$ of a prime number $\ell$ and a field isomorphism
$\iota\colon\C\to\overline{\Q}_\ell$ onto an algebraic closure of the field of $\ell$-adic numbers, with its absolute
value $|\cdot|_\ell$. A family $Z=(Z_{\ell,\iota})$ of complex numbers has \emph{height at most~$h$}, for a real
number~$h$, if there are a non-zero complex number $H$, a real number $\alpha$ and real numbers $\eta_\ell$,
$\ell\in\NN$, such that $|\sigma(H)|\le\ee^{\alpha}$ for all $\sigma\in\AutC$,
$|\iota(H\,Z_{\ell,\iota})|_\ell\le\ee^{\eta_\ell}$ for all pairs $(\ell,\iota)$, and
$\alpha+\sum_{\ell<i,\ \ell\text{ prime}}\eta_\ell\le h$ for all sufficiently large natural numbers~$i$ (unit~135). A
family $\theta=(\theta_{\ell,\iota})$ of $\Dm{D}$-fold twists of $D$ may be \emph{fit for~$D$}; this property is defined in
unit~138 and bounds, for each $\ell$, the congruences of the lifts at the twists $\theta_{\ell,\iota}$. Of its
definition we use only that $\ell$ is prime to $\theta_{\ell,\iota}$ for every pair $(\ell,\iota)$; otherwise it enters
through Inputs \ref{i:volume}, \ref{i:lift} and~\ref{i:crude}.
\end{setting}

\statement{ana}{\inputhead{heights}{unit~135}\label{i:heights}
For families $Z$, $Z'$ of complex numbers indexed by the pairs $(\ell,\iota)$ and real numbers $h$, $h'$:

\emph{(a)} if $Z$ has height at most $h$ and $Z'$ has height at most $h'$, then the family of the products
$Z_{\ell,\iota}Z'_{\ell,\iota}$ has height at most $h+h'$; the same holds for finite products, and the family of the
$n$-th powers of $Z$ has height at most $n\,h$;

\emph{(\scb{D})} if $Z$ has height at most $h$, $h\le h'$, and $|\iota(Z'_{\ell,\iota})|_\ell\le|\iota(Z_{\ell,\iota})|_\ell$ for
all pairs, then $Z'$ has height at most $h'$;

\emph{(c)} if the constant family $1$ has height at most $h$, then $h\ge0$.
\src{(c) is the product formula, applied to the norm of the auxiliary number~$H$. Chapter~IV.}}

\statement{ana}{\inputhead{fit families exist}{unit~138}\label{i:fit}
In Settings \ref{st:data}, \ref{st:forms} and~\ref{st:heights} there is a family $\theta=(\theta_{\ell,\iota})$ of
$\Dm{D}$-fold twists of $D$ that is fit for~$D$.
\src{The subject of Chapter~VIII.}}

\statement{ana}{\inputhead{the volume of the chain}{units 143, 147}\label{i:volume}
In Settings \ref{st:data}, \ref{st:forms} and~\ref{st:heights}, say that a complex-valued function
$\vartheta\mapsto\operatorname{cst}(\vartheta)$ on the $\Dm{D}$-fold twists of $D$ \emph{costs $n$ squares}, for a natural
number~$n$, if for every family $\theta$ that is fit for~$D$ the family
$(\ell,\iota)\mapsto|\operatorname{cst}(\theta_{\ell,\iota})/\operatorname{cst}(1)|^{-2}$ has height at most
$n\cdot2\,t(D)\,\scx{D}^{999}$. There is a function $\cst{vol}{D}$ that costs $\Dm{D}^{2}$ squares such that for every $\Dm{D}$-fold
twist~$\vartheta$,
\[
\zz{D}{\vartheta}\ne0\qquad\text{and}\qquad
\zz{D}{\vartheta}=\bigl|\cst{vol}{D}(\vartheta)\bigr|^{\dim\Row{D}}\cdot\prod_{j=1}^{\Dm{D}}\vol{D}{\vartheta}{j},
\]
where $\vol{D}{\vartheta}{j}$ is the volume of the $j$-th lift of the chain at the twist $\vartheta^{j}$
(unit~143).}

\statement{ana}{\inputhead{the height of one lift}{unit~143}\label{i:lift}
In Settings \ref{st:data}, \ref{st:forms} and~\ref{st:heights}, let $\theta$ be a family of $\Dm{D}$-fold twists that
is fit for~$D$, and let $j$ be a natural number with $1\le j\le \Dm{D}$. Then the family
$(\ell,\iota)\mapsto\vol{D}{1}{j}/\vol{D}{\theta_{\ell,\iota}}{j}$ has height at most
\[\dim\Row{D}\cdot k(D)\cdot\Bigl(\log \scb{D}\cdot\frac{\budget{D}}{\Dm{D}\,k(D)}-\cost{D}\Bigr).\]
\src{Each lift repays one direction: this is where $\cost{D}$, and so the main amount, enters the lower bound.
Chapter~VII.}}

\statement{ana}{\inputhead{crude rows}{unit~147}\label{i:crude}
In Settings \ref{st:data}, \ref{st:forms} and~\ref{st:heights} there is a function $\cst{cr}{D}$ on the $\Dm{D}$-fold
twists of $D$, with non-zero complex values, that costs $2\Dm{D}^{2}$ squares, such that $\PF{D}$ is spanned over $\C$ by
the forms $F$ with these two properties: for every $\Dm{D}$-fold twist $\vartheta$, every prime number $\ell$ that is
prime to~$\vartheta$ and every ring homomorphism $j\colon\C\to\C_\ell$, the form
$\bigl(\cst{cr}{D}(\vartheta)/\cst{cr}{D}(1)\bigr)\cdot\Thf{D}{\vartheta}(F)$ is $j$-integral; and $\Thf{D}{1}(F)$ is
pointwise at most $\exp(\scxo{D}^{2}\scb{D}^{2})$ at every $\sigma\in\AutC$. We call such forms \emph{crude}.}

\statement{ana}{\inputhead{one determinant}{unit~135}\label{i:onedet}
In Settings \ref{st:data}, \ref{st:forms} and~\ref{st:heights}, let $\theta=(\theta_{\ell,\iota})$ be a family of
$\Dm{D}$-fold twists of $D$ and let $y_1,\dots,y_{n'}$ be a basis of $\Row{D}$. Let $B$, $\delta$, $h$ be real numbers,
$n''$ a natural number, $(\eta_\ell)_{\ell\in\NN}$ real numbers and $(H_{\ell,\iota})$ a family of non-zero complex
numbers. Assume:
\begin{enumerate}[label=(\arabic*)]
\item $\Th{D}{1}$ is injective, and $|\sigma(\det(\langle\Th{D}{1}y_i,\Th{D}{1}y_{i'}\rangle)_{i,i'})|\le\ee^{B}$ for
every $\sigma\in\AutC$;
\item the family $(\ell,\iota)\mapsto|H_{\ell,\iota}|^{-2}$ has height at most~$h$;
\item $\sum_{\ell\in J}\eta_\ell\le\delta$ for every finite set $J\subset\NN$;
\item for every pair $(\ell,\iota)$ and every ring homomorphism $j\colon\C\to\C_\ell$: if
$F,F'\in\AAf{D}{\theta_{\ell,\iota}}$, $F'$ is $j$-integral and $F$ is integral for the composite of $j$ with complex
conjugation, then $|j(\langle F,F'\rangle)|_\ell\le1$;
\item for every pair $(\ell,\iota)$ and every ring homomorphism $j\colon\C\to\C_\ell$ there is a complex
$n'\times n'$ matrix $u$ with $\ee^{-\eta_\ell}\,|j(H_{\ell,\iota})|_\ell^{\,n''}\le|j(\det u)|_\ell$ such that all
$\Th{D}{\theta_{\ell,\iota}}(\sum_{i'}u_{ii'}y_{i'})$ are $j$-integral.
\end{enumerate}
Then the family $(\ell,\iota)\mapsto\zz{D}{\theta_{\ell,\iota}}/\zz{D}{1}$ has height at most $B+2\delta+n''\,h$.
\src{A special case: the lemma of unit~135 is stated for an arbitrary family of linear maps from one space into
inner product spaces. Chapter~IV.}}

\statement{ana}{\inputhead{integral forms pair integrally}{units 102, 152, 147}\label{i:pairing}
In Settings \ref{st:data} and~\ref{st:forms}, hypothesis (4) of Input~\ref{i:onedet} holds for every twist: for a
twist~$\vartheta$, a prime number $\ell$, a ring homomorphism $j\colon\C\to\C_\ell$ and $F,F'\in\AAf{D}{\vartheta}$
with $F'$ $j$-integral and $F$ integral for the composite of $j$ with complex conjugation,
$|j(\langle F,F'\rangle)|_\ell\le1$.
\src{Unit~102 proves this with a factor: a common multiple of the orders of the stabilisers that occur in the
Petersson product. Units 152 and~147 show that at the level of $\AAf{D}{\vartheta}$ all these stabilisers are
trivial.}}

\statement{ana}{\inputhead{Hadamard and Fischer}{unit~102}\label{i:hadamard}
In Settings \ref{st:data} and~\ref{st:forms}, let $\vartheta$ be a twist, $\sigma\in\AutC$, $B$ and $V$ real
numbers, and $F_1,\dots,F_n,F'_1,\dots,F'_{n''}\in\AAf{D}{\vartheta}$.

\emph{(a)} If $\nsq{F'_i}{\sigma}\le\ee^{B}$ for all $i$, then $\covol{F'_1,\dots,F'_{n''}}{\sigma}\le\ee^{n''B}$.

\emph{(\scb{D})} If moreover $\covol{F_1,\dots,F_n}{\sigma}\le\ee^{V}$, then
$\covol{F_1,\dots,F_n,F'_1,\dots,F'_{n''}}{\sigma}\le\ee^{V+n''B}$.

\emph{(c)} If $\covol{F_1,\dots,F_n}{\sigma'}\le\ee^{V}$ for every $\sigma'\in\AutC$, then
$|\sigma(\det(\langle F_i,F_{i'}\rangle)_{i,i'})|\le\ee^{V}$.}

\begin{lemma}[the height of the inverse ratio of volumes]\label{l:ratio}
Let $P$ be parameters and $D$ admissible data for~$P$, with components $\scx{D}$, $\scxo{D}$, $\scb{D}$, $\Dm{D}$, $\Dr{D}$ (Settings
\ref{st:data}, \ref{st:forms} and~\ref{st:heights}), and let $\theta=(\theta_{\ell,\iota})$ be a family of $\Dm{D}$-fold
twists of $D$ that is fit for~$D$. Then the family
$(\ell,\iota)\mapsto\zz{D}{1}^{2}/\zz{D}{\theta_{\ell,\iota}}^{2}$ has height at most
\[2\,\dim\Row{D}\cdot\bigl(2\log \scb{D}\cdot\budget{D}-\mainHeight{D}\bigr).\]
\end{lemma}
\begin{proof}
By Input~\ref{i:volume}, for every pair $(\ell,\iota)$,
\[
\frac{\zz{D}{1}^{2}}{\zz{D}{\theta_{\ell,\iota}}^{2}}
=\Bigl(\prod_{j=1}^{\Dm{D}}\frac{\vol{D}{1}{j}}{\vol{D}{\theta_{\ell,\iota}}{j}}\Bigr)^{2}\cdot
\Bigl(\Bigl|\frac{\cst{vol}{D}(\theta_{\ell,\iota})}{\cst{vol}{D}(1)}\Bigr|^{-2}\Bigr)^{\dim\Row{D}}.
\]
By Inputs \ref{i:lift} and~\ref{i:heights}(a) the product over $j$ has height at most $\Dm{D}$ times the bound of
Input~\ref{i:lift}, which is $\dim\Row{D}\cdot(\log \scb{D}\cdot\budget{D}-\mainHeight{D})$ by Definition~\ref{d:hmain}; here
$\Dm{D}\,k(D)\ne0$ by (S0). The function $\cst{vol}{D}$ costs $\Dm{D}^{2}$ squares, so the last family in brackets has height
at most $\Dm{D}^{2}\cdot2\,t(D)\,\scx{D}^{999}\le2\log \scb{D}\cdot\budget{D}$, by (S5) and Input~\ref{i:heights}(\scb{D}). By
Input~\ref{i:heights}(a) the height of the product is at most
$2\dim\Row{D}\,(\log \scb{D}\cdot\budget{D}-\mainHeight{D})+\dim\Row{D}\cdot2\log \scb{D}\cdot\budget{D}$.
\end{proof}
\leanline{\lean{heightLevolumeratiosq}}

\begin{setting}[an almost spanning family of small index and small covolume]\label{st:family}
Let $P$ be parameters and $D$ admissible data for~$P$, with components $\scx{D}$, $\scxo{D}$, $\scb{D}$, $\Dm{D}$, $\Dr{D}$ (Settings
\ref{st:data} and~\ref{st:forms}). Let $n$ be a natural number, $y=(y_1,\dots,y_n)$ a family of rows in $\Row{D}$,
and $\delta$ and $V$ real numbers, such that:
\begin{enumerate}[label=(\roman*)]
\item the family $y$ is linearly independent over~$\C$;
\item $(1-4/\scb{D}^{2})\cdot\dim\Row{D}\le n$;
\item \indexWords{D}{y}{\delta};
\item $\covol{\Th{D}{1}y_1,\dots,\Th{D}{1}y_n}{\sigma}\le\ee^{V}$ for every $\sigma\in\AutC$.
\end{enumerate}
\end{setting}

\begin{lemma}[the height of the ratio of volumes]\label{l:volrows}
In Setting~\ref{st:family}, let $\theta=(\theta_{\ell,\iota})$ be a family of $\Dm{D}$-fold twists of $D$ that is fit
for~$D$ (Setting~\ref{st:heights}). Then the family $(\ell,\iota)\mapsto\zz{D}{\theta_{\ell,\iota}}/\zz{D}{1}$ has
height at most $V+2\delta+\dim\Row{D}\cdot \scb{D}$.
\end{lemma}
\begin{proof}
\emph{Completion to a basis.} Let $\cst{cr}{D}$ be as in Input~\ref{i:crude}. The map $\gpt{D}{1}$, for the unit
element of $\GD{D}$, maps $\PF{D}$ onto $\Row{D}$, so the rows $\gpt{D}{1}F$ with $F$ crude span $\Row{D}$. Since $y$
is linearly independent (i), there are $n''=\dim\Row{D}-n$ crude forms $F_1,\dots,F_{n''}$ such that
$y_1,\dots,y_n,\gpt{D}{1}F_1,\dots,\gpt{D}{1}F_{n''}$ is a basis of $\Row{D}$. We apply Input~\ref{i:onedet} to this
basis, to the family $\theta$, and to $H_{\ell,\iota}=\cst{cr}{D}(\theta_{\ell,\iota})/\cst{cr}{D}(1)$.

\emph{Hypothesis (1).} $\Th{D}{1}$ is injective because $\zz{D}{1}\ne0$ (Input~\ref{i:volume}). Each $F_i$ is crude,
so $\nsq{\Th{D}{1}(\gpt{D}{1}F_i)}{\sigma}\le\exp(\scx{D}^{99}+\scxo{D}^{2}\scb{D}^{2})$ for every $\sigma$, by
Input~\ref{i:goodpart}(a). With (iv) and Input~\ref{i:hadamard}(\scb{D}), (c), the conjugates of the Gram determinant of
the images of the basis under $\Th{D}{1}$ have absolute value at most $\ee^{B}$ with
$B=V+n''\,(\scx{D}^{99}+\scxo{D}^{2}\scb{D}^{2})$.

\emph{Hypotheses (2), (3), (4).} The function $\cst{cr}{D}$ costs $2\Dm{D}^{2}$ squares and $\theta$ is fit, so (2) holds
with $h=2\Dm{D}^{2}\cdot2\,t(D)\,\scx{D}^{999}$. Let $(\eta_\ell)$ be as in Definition~\ref{d:index} for the family $y$, by
(iii); this is (3). And (4) is Input~\ref{i:pairing}.

\emph{Hypothesis (5).} Let $(\ell,\iota)$ be a pair and $j\colon\C\to\C_\ell$ a ring homomorphism. The prime $\ell$ is
prime to $\theta_{\ell,\iota}$, because $\theta$ is fit. By (iii) there is an $n\times n$ matrix $u_0$ over $O_j$ with
$|j(\det u_0)|_\ell\ge\ee^{-\eta_\ell}$ such that all
$\Th{D}{\theta_{\ell,\iota}}(\sum_{i'}(u_0)_{ii'}y_{i'})$ are $j$-integral. Let $u$ be the block matrix with the blocks
$u_0$ and $H_{\ell,\iota}$ times the identity matrix of size~$n''$. Then
$|j(\det u)|_\ell\ge\ee^{-\eta_\ell}|j(H_{\ell,\iota})|_\ell^{\,n''}$. The new rows of the second block are the
$H_{\ell,\iota}\,\gpt{D}{1}F_i$, and $\Th{D}{\theta_{\ell,\iota}}(H_{\ell,\iota}\,\gpt{D}{1}F_i)
=H_{\ell,\iota}\Thf{D}{\theta_{\ell,\iota}}(\gpt{D}{1}F_i)$ is $j$-integral, by the first property of crude forms and
Input~\ref{i:goodpart}(\scb{D}).

\emph{Conclusion.} Input~\ref{i:onedet} gives the height bound
$V+2\delta+n''\,\bigl(\scx{D}^{99}+\scxo{D}^{2}\scb{D}^{2}+4\Dm{D}^{2}t(D)\scx{D}^{999}\bigr)$. By (ii), $n''\le4\scb{D}^{-2}\dim\Row{D}$, and by (S4)
the bracket is at most $\tfrac14\scb{D}^{3}$. So the bound is at most $V+2\delta+\dim\Row{D}\cdot \scb{D}$, and
Input~\ref{i:heights}(\scb{D}) concludes.
\end{proof}
\leanline{\lean{heightLevolumerows}}

\statement{arith}{\phantomsection\label{t:C}\textbf{Theorem C} (integral forms are long)\textbf{.}\enspace Let $P$ be
parameters and $D$ admissible data for~$P$, with components $\scx{D}$, $\scxo{D}$, $\scb{D}$, $\Dm{D}$, $\Dr{D}$ (Settings \ref{st:data}
and~\ref{st:forms}). Let $n$ be a natural number, $y=(y_1,\dots,y_n)$ a family of rows in $\Row{D}$, and $\delta$ and
$V$ real numbers. Assume that the family $y$ is linearly independent over~$\C$, that
$(1-4/\scb{D}^{2})\cdot\dim\Row{D}\le n$, that \indexWords{D}{y}{\delta}, and that
$\covol{\Th{D}{1}y_1,\dots,\Th{D}{1}y_n}{\sigma}\le\ee^{V}$ for every $\sigma\in\AutC$. Then
\[\dim\Row{D}\cdot\bigl(\mainHeight{D}-3\log \scb{D}\cdot\budget{D}\bigr)\le V+2\delta.\]}

\begin{proof}
The hypotheses are those of Setting~\ref{st:family}. Let $\theta$ be a family of $\Dm{D}$-fold twists that is fit for~$D$
(Input~\ref{i:fit}). For every pair $(\ell,\iota)$ the numbers $\zz{D}{1}$ and $\zz{D}{\theta_{\ell,\iota}}$ are not
zero (Input~\ref{i:volume}), so
\[\Bigl(\frac{\zz{D}{\theta_{\ell,\iota}}}{\zz{D}{1}}\Bigr)^{2}\cdot\frac{\zz{D}{1}^{2}}{\zz{D}{\theta_{\ell,\iota}}^{2}}=1.\]
By Lemma~\ref{l:volrows} and Input~\ref{i:heights}(a) the first factor has height at most
$2\,(V+2\delta+\dim\Row{D}\cdot\scb{D})$; by Lemma~\ref{l:ratio} the second has height at most
\[2\dim\Row{D}\,\bigl(2\log \scb{D}\cdot\budget{D}-\mainHeight{D}\bigr).\]
So the constant family $1$ has height at most the sum, and
by the product formula, Input~\ref{i:heights}(c),
\[\begin{aligned}0\le{}&2\,(V+2\delta+\dim\Row{D}\cdot \scb{D})\\&+2\dim\Row{D}\cdot\bigl(2\log \scb{D}\cdot\budget{D}-\mainHeight{D}\bigr).\end{aligned}\]
With $\scb{D}\le\log \scb{D}\cdot\budget{D}$, which is (S6), this is the claim.
\end{proof}
\leanline{\lean{integralformslong}}

\begin{corollary}[Theorems A and B together]\label{c:AB}
Let $P$ be parameters and $D$ admissible data for~$P$, with components $\scx{D}$, $\scxo{D}$, $\scb{D}$, $\Dm{D}$, $\Dr{D}$ (Settings
\ref{st:data} and~\ref{st:forms}). There are a natural number $n$ and a family $y=(y_1,\dots,y_n)$ of rows in
$\Row{D}$ such that: the family $y$ is linearly independent over~$\C$; $(1-4/\scb{D}^{2})\cdot\dim\Row{D}\le n$; there is a
real number $\delta\le5\,\dim\Row{D}\cdot\log(1+\scb{D}^{2})\cdot\budget{D}$ such that \indexWords{D}{y}{\delta}; and for
every $\sigma\in\AutC$,
\[
\begin{gathered}
\covol{\Th{D}{1}y_1,\dots,\Th{D}{1}y_n}{\sigma}\\
\le\exp\Bigl(n\cdot\bigl(\mainHeight{D}-\Dm{D}\cdot\remainder{D}+\log \scb{D}\cdot\budget{D}\bigr)\Bigr).
\end{gathered}
\]
\end{corollary}
\begin{proof}
Let $q$ and $S$ be as in Theorem~A, and $\mathcal{T}$ as in Theorem~B. Since $S$ spans $\Colram{D}$ and $q\ne0$, the
rows $q\,\projectedPoincare{D}{\gamma}{\psi}$ with $\gamma\in\mathcal{T}$, $\psi\in S$ span the space
$\sum_{\gamma\in\mathcal{T}}\projectedPoincare{D}{\gamma}{(\Colram{D})}$. Choose a basis $y_1,\dots,y_n$ of this space
among them. Then $y$ is linearly independent and $n\ge(1-4/\scb{D}^{2})\dim\Row{D}$ by Theorem~B. By Theorem~A(\scb{D}), with
$S\subseteq\integralRamified{D}$, the family $y$ has $D$-index at most
$\delta=\#\mathcal{T}\cdot\dim\Col{D}\cdot\budget{D}$, and
$\delta\le5\dim\Row{D}\cdot\log(1+\scb{D}^{2})\cdot\budget{D}$ by Theorem~B and $\budget{D}>0$. Finally each
$\nsq{\Th{D}{1}y_i}{\sigma}$ is at most $\exp(\mainHeight{D}-\Dm{D}\cdot\remainder{D}+\log \scb{D}\cdot\budget{D})$ by
Theorem~A(a), and Hadamard's inequality, Input~\ref{i:hadamard}(a), bounds the covolume.
\end{proof}
\leanline{\lean{existsfamilycovolumeSqle}}

\sect{Where the results are used}{s:use}
All consumers are in the last unit of the proof (unit~154, Chapter~I), which has three steps.

\emph{The main inequality} $\Dm{D}\cdot\remainder{D}\le(\log \scb{D})^{2}\cdot\budget{D}$. Take the family of
Corollary~\ref{c:AB} and apply Theorem~C to it, with $V=n\,(\mainHeight{D}-\Dm{D}\cdot\remainder{D}+\log \scb{D}\cdot\budget{D})$.
The main amount occurs on both sides, once with the factor $\dim\Row{D}$ and once with the factor~$n$;
Lemma~\ref{l:crude} pays for the difference, since $\dim\Row{D}-n\le4\scb{D}^{-2}\dim\Row{D}$. After division by
$\dim\Row{D}$ the main amount has cancelled, and what remains is the inequality.

\emph{The one-sided bound} $\primeError{d}{N_0}{x}\le x^{8\varepsilon-1/2}$ for large~$x$. Theorem~4 solves for
$\primeError{d}{N_0}{x}$, and leaves the combination
$\primeError{d}{N_0}{\scxo{D}}-X\,\primeError{d}{1}{\scxo{D}}$ with $X=\#\tops{D}/\#\thick{D}$ to be bounded from above.
Theorem~1, for the moduli $N_0$ and~$1$, and Theorem~3 provide a point $v_1$ in the window $[v_0,v_0^{2}]$,
$v_0=10^{6}\log x$, at which this combination is at most $\ee^{-v_0}$, and the second scale is moved to
$\ee^{v_1}$; the main inequality bounds the rest.

\emph{The zero-free half-plane.} Theorem~2 with $M=N_0$ and $c=N_0/2-1$.

Inside the unit, Lemmas \ref{l:lambdabound} to~\ref{l:constant} serve Theorems 1 and~4 and Lemma~\ref{l:crude};
Proposition~\ref{p:moments} serves Theorem~3; Proposition~\ref{p:level} and Lemma~\ref{l:rational} serve Theorem~A;
Lemmas \ref{l:ratio} and~\ref{l:volrows} serve Theorem~C.

\sect{Notes on the formal text}{s:notes}

\subhead{Conventions of Lean that the statements hide.} Division by zero gives zero. So the residue function of
Theorem~1 is defined at $s=0$ and $s=1$, with the value~$0$, as stated; at $s=1$ the numerator vanishes anyway, by
Input~\ref{i:nonvanishing}. The logarithmic derivative of a function is $0$ at its zeros, so
$\logDerivMean{d}{M}{s}$ is defined for every~$s$; the proofs use it only as a meromorphic function and at points
where no $\entireL{\xi}$, $\entireL{\quadTwist{\xi}{d}}$ vanishes. The remainder $\primeError{d}{M}{x}$ is defined for
every real~$x$, with Lean's logarithm of a number $\le0$; only $x>0$ is used. The integral of a function that is not
integrable is~$0$, which is why Proposition~\ref{p:moments} and Input~\ref{i:laplace} assert integrability
separately. The weight $\lambdaClass{d}{M}{\nu}$ and the sums $\lambdaSum{d}{M}{i}{x}$ are defined for every natural
number~$M$, also $M=0$, where $\Z/0=\Z$ and the weight vanishes identically.

\subhead{Definition~\ref{d:polar}.} The formal definition asks that $\Phi$ be meromorphic at every point of~$\C$ as
a function on $\C$; its values at its poles are arbitrary. Condition (P3) is stated for the product
$(s-1)\Phi(s)$ as a function, so its value at $s=1$ is~$0$; the clause on that value in (P3) and
Remark~\ref{r:polar} are our reading of it, and the comment above the formal definition says the same: no pole
at~$1$ is allowed.

\subhead{Theorem~1.} The hypothesis $d\equiv7\pmod 8$ is used only through $d\ne2$ (Step~3). The formal proof
orders the steps differently: it proves the Laplace identity (our Steps 3 to~5) after it has used it, as a pending
goal, to show that the constant is real (Step~6); the order of the goals was inferred by the reader of the
dossier from the structure of the term, and the writer agrees with that reading. The formal statement speaks of
the function $v\mapsto\primeError{d}{M}{\ee^{v}}$ and of a residue function given by a formula; the names
$\Phi_{d,M}$ and $\operatorname{res}_{d,M}$ are ours, after the overview of the write-up, and the transform (5.1)
and the constant $|\kap{d}{M}|+5$ are the witnesses of the formal proof.

\subhead{Theorem~2.} The conclusion is about $\Ltv{\rho}{\chi}$; for principal $\chi$ this is
$(\rho-1)L(\rho,\chi)$. The hypothesis on the primes $d\equiv c$ is one conjunction in the formal text, and it is
used three times: for Input~\ref{i:density}, for Theorem~1, and for the one-sided bound. In the formal statement
$d$ runs over the type of prime numbers and $d\equiv c\pmod M$ is written as an equality in $\Z/M$.

\subhead{Proposition~\ref{p:moments}: the boundary of the statement.} The formal statement contains a local
definition of the function $v\mapsto v^{i}\ee^{-3v}\cE(v)$, written with a sign of definition. The census of the
workspace, which reads the text line by line, takes that sign for the end of the statement; the dossier corrects
this, and the writer read the complete statement in the source: it ends after the bound $i!\,v_0/\rr^{i}$. The
quantifiers of the formal statement are filters, \emph{for all $\rr>3$ close to~$3$} and \emph{for all large
$v_0$}; the proposition writes them out, in this order. Its proof uses neither the meromorphy of $\Phi$ nor (P4);
the residues of a polar function enter the proof only through Landau's lemma.

\subhead{Arguments that the census lists and Lean does not take.} The census reads the arguments of a declaration
from its text and counts every variable of the surrounding section whose name occurs in the statement. In three
places a bound variable of the statement has the name of such a section variable, and Lean, which the writer follows,
does not take it as an argument. \emph{(1)} Corollary~\ref{c:AB}: its statement binds its own $n$, $y$ and $\delta$
under existential quantifiers, and its only arguments are the parameters and the data. The reader of the dossier
could not confirm this with Lean; the writer confirms it from the use of the corollary in the proof of the main
inequality, where it is applied to the data alone and nothing could determine further arguments. \emph{(2)}
Input~\ref{i:density}: the character $\xi$ is bound by a universal quantifier inside the statement and is not an
argument. \emph{(3)} Definition~\ref{d:proj}: the formal declaration takes the data and $\gamma$, and its value is a
linear map; the census reads the arrow in the type of that map as one more argument. The notation of this note has
a third place for that reason, and it is left empty when the map itself is meant.

\subhead{Theorem~A.} The formal statement names the family of maps $\gamma\mapsto q\,\projectedPoincare{D}{\gamma}{}$:
it says that there are $q$ and a family equal to this one, with the two properties. We have substituted the
family. In~(b) the finite sets $J$ of Definition~\ref{d:index} are arbitrary finite sets of natural numbers, not
only sets of primes, as in the formal text. The formal proof of (b) obtains the bound for the determinant from
divisibility in the valuation ring $O_j$, as we do.

\subhead{Theorem~B.} Its first two assertions, that there are columns and rows, are part of the formal statement.
The filling theorem is stated for columns and the operator $\operatorname{ram}_D$, and Theorem~B for ramified
columns; the proof passes from one to the other.

\subhead{Theorem~C.} In the formal text the four hypotheses of Setting~\ref{st:family} are variables of the section,
shared by Lemma~\ref{l:volrows} and Theorem~C, in the order (i), (ii), (iii), (iv); Lemma~\ref{l:ratio} takes none of
them. The formal proof of Lemma~\ref{l:volrows} gets the injectivity of $\Th{D}{1}$ from a lemma that is applied
automatically and that we have not located; we derive it from $\zz{D}{1}\ne0$, which is the same fact.

\subhead{Two names for the chain.} The formal text has two maps for the chain of lifts on $\PF{D}$: one with values
in forms that are functions with vector values, for which pointwise sizes and integrality of values are defined,
and one with values in the space that carries the Petersson product. They have the same values, and the proofs of
Theorems A and~C pass from one to the other without comment; we write $\Thf{D}{\vartheta}$ for both.

\subhead{What the reader of the dossier was unsure of.} Besides the three points above (the boundary of the
statement of Proposition~\ref{p:moments}, the arguments of Corollary~\ref{c:AB} and of Definition~\ref{d:proj}): the
facts of Mathlib that Lemma~\ref{l:constant} needs, namely that $(s-1)\zeta(s)$ has the value $1$ and the derivative
$\Eu$ at $s=1$, were inferred by the first reader from the proof going through, because the pinned version of
Mathlib was not at hand. The second reader of this write-up has since found both at the pin, and also the value
$\prod_{p\mid M}(1-1/p)$ of Definition~\ref{d:Lt} and the two statements of Mathlib that Step~3 and Step~4 of the
proof of Theorem~1 quote (orthogonality of characters against von Mangoldt's function, and partial summation). For
this note the writer did not open Mathlib: every fact of Mathlib that the note uses rests on the second reader's
check at the pin. This is said of this note only; the three notes of Chapter~I list in their appendices the
declarations of Mathlib that their writer read at the pinned commit. The readings of the Inputs \ref{i:sizes} to~\ref{i:hadamard} come from their formal statements and
comments, not from their proofs; the chapters on the units that own them are the reference.

\subhead{The dossier, and its second reading.} The second reader of this write-up compared the dossier with the
source and asked for six corrections; this note agrees with all of them. Two concern words of the dossier that
were wrong as mathematics while its quotations were exact. In the structured proof of Theorem~4, one step says
that $(-\nu)\cdot\tau$ runs over the places of class $-\nu$ as $\tau$ runs over the tops; this holds for $\nu$ prime to
$N_1(D)$ only, and for the other $\nu$ the formal text makes $(-\nu)\cdot\tau=\tau$ (Step~1 above treats both cases).
In the structured proof of Lemma~\ref{l:crude}, an intermediate bound had absolute values around $z/\nu-Y$ where the
positive part stands; with the positive part the bound $|z|/\nu$ is correct. The other four are points of
exposition: the quadratic character is that of $K$ only for $d\equiv3\pmod 4$ (Definition~\ref{d:chi}); the value of
the entire $L$-function of a principal character at $s=1$ (Definition~\ref{d:Lt}); the vanishing at $s=1$ that makes
the singularity removable (Remark~\ref{r:polar}, Input~\ref{i:taylor}); and the difference between the rank $k(D)$
and the dimension of $\Col{D}$ (Setting~\ref{st:forms}).

\subhead{What of the unit the final theorems use.} The last unit of the proof uses Theorems 1, 2, 3, 4 and~C,
Corollary~\ref{c:AB} and Lemma~\ref{l:crude} directly; every other declaration of the unit is used by one of these.

\subhead{How this note was written.} From the dossier of unit~151: its statements, its \UnitProofs{} structured
proofs with \UnitSteps{} leaf steps, which a program compared with the source, and its readings of the declarations
of other units. The writer read in the Lean text itself: the four files of the unit completely, statements and
proofs; the file of the background library with the definition of polar functions, Landau's lemma and the
elementary inputs, completely; the definitions of the characters, of the entire $L$-function and of the
Chowla--Selberg constant; the statement and the last step of the proof of Input~\ref{i:density}; the record of the
data and the definitions attached to it that Setting~\ref{st:data} uses; the statements, as quoted in the dossier, of
every Input of Part~II, and in the source the definitions of the height, of fit families, of twists, of the index
relation of unit~149 and of the mechanism that compares sizes; and the last unit of the proof, to see how the
results are used. Every statement of the note was compared with its formal statement, argument by argument and in
the order of the quantifiers. The relation of Inputs \ref{i:nonvanishing}, \ref{i:landau} and~\ref{i:density} to the
literature is the second reader's comparison, as their boxes say. Not opened for this note: the proofs of the
Inputs of Part~II, and Mathlib, whose facts were checked at the pin by the second reader, as said above.

The arguments are those of the formal proofs, with these differences. In Theorem~3 the final numerical
comparison is done by hand here and by an arithmetic tactic there. In Theorem~4 the counting of Step~1 is written
out here; the formal proof reindexes the sum by a translation of the group of places and uses a lemma of unit~152,
which says that the place named by an integer $\nu$ has class $-1$ modulo $N_0$ exactly when $\nu\equiv-1$, also for
the integers that are not prime to $N_1(D)$. Remarks \ref{r:heuristic},
\ref{r:landau}, \ref{r:anatomy} and~\ref{r:window}, and the estimate of orders of magnitude after
Lemma~\ref{l:crude}, are ours and have no formal counterpart.

\head{Appendix A. Concordance with the Lean artifact}
{\fontsize{8.6}{11.6}\selectfont\renewcommand{\leanfont}{\ttfamily\fontsize{7.7}{10}\selectfont\color{leancol}}
\setlength{\LTleft}{0pt}\setlength{\LTright}{\fill}\setlength{\tabcolsep}{0pt}
\begin{longtable}{@{}p{30mm}p{56mm}p{43mm}r@{}}
\textbf{In this note} & \textbf{Lean declaration} & \textbf{File} & \textbf{Lines}\\\midrule\endhead
\ConcordanceRows
\end{longtable}}

\subhead{Legend.} The note changes some letters of the formal text, never the number of arguments. The modulus is
$M$ here, and $M$ or $N$ there. The unit class $c$ of Theorem~2 and Input~\ref{i:density} is \texttt{a}; the height
$z$ of Input~\ref{i:density} is \texttt{y}. The exponent $i$ of the weighted sums is \texttt{r}, and their real
argument $x$ is \texttt{y}; the natural number $a$ and the point $x$ of Lemma~\ref{l:Rsum} are \texttt{A} and
\texttt{z}. A polar function $\cE$ with residue function $\omega$ is \texttt{E} with \texttt{r}; in Theorem~3,
$\cE_1,\omega_1,\cE_2,\omega_2$ are \texttt{E}, \texttt{r}, \texttt{E'}, \texttt{r'}, the parameter $v_0$ is
\texttt{t}, and the point $v_1$ is \texttt{u}. In Proposition~\ref{p:moments} the radius $\rr$ is \texttt{R}, the
number $v_0$ is \texttt{t} and the order $i$ is \texttt{k}. In Input~\ref{i:taylor}, $s_0,s_1,s_2$ and $\Phi$ are
\texttt{z}, \texttt{a}, \texttt{σ} and \texttt{Ψ}, and the radius $\rr_0$ is \texttt{ε}. In Lemma~\ref{l:rational} the set $J$ is \texttt{T}, the exponents
$e_\ell$ are \texttt{g l}, and primes are \texttt{l} throughout. In Part~II the scales $\scx{D}$, $\scxo{D}$, $\scb{D}$ are the exponentials
of the components \texttt{logX}, \texttt{logX₁}, \texttt{logB} of the data; a place $\tau$ is \texttt{σ} or
\texttt{τ}; an element $\gamma$ of $\GD{D}$ is \texttt{s}; a finite subset $\mathcal{T}$ is \texttt{t}; a ramified
column $\psi$ is \texttt{φ}; a ring homomorphism $j$ is \texttt{ι}, of the class \texttt{PadicPlace}; but in Inputs~\ref{i:volume} and~\ref{i:lift} and in Lemma~\ref{l:ratio} the letter $j$ is the index of a lift, \texttt{j} in the formal text; the numbers
$\eta_\ell$ are \texttt{e l}, and a finite set $J$ of natural numbers is \texttt{U}. The objects of other units
have these formal names: $\tops{D}$, $\thick{D}$, $\budget{D}$, $\cost{D}$, $\hrow{D}$ are \texttt{Data.tops},
\texttt{Data.thick}, \texttt{Data.budget}, \texttt{Data.dirCost}, \texttt{Data.heightPerRow}; $\wD{D}{\tau}$,
$\Ask{D}{z}{\tau}$ and $\jet{D}{i}$ are \texttt{Data.nominalWeight}, \texttt{Data.primeSum} and
\texttt{Data.jetSize}; $\Col{D}$, $\Colram{D}$, $\Row{D}$, $\GD{D}$ are \texttt{Data.Column},
\texttt{Data.ramifiedColumns}, \texttt{Data.rows}, \texttt{Data.G}; $\Poin{D}$ and $\gpt{D}{\gamma}$ are
\texttt{Data.poincare} and \texttt{Data.goodProjTranslate}; $\Th{D}{\vartheta}$ and $\zz{D}{\vartheta}$ are
\texttt{Data.Θ} and \texttt{Data.𝔷}; the height is \texttt{HeightLe}, and fitness is \texttt{Data.IsFitFamily}.

\subhead{Six statements, as the artifact has them.} The remainder; Theorem~1; Theorem~2; polar functions;
Landau's lemma; Theorem~4.
\quoted{primeError}
\quoted{isPolarprimeError}
\quoted{entireLnezeroofprimeErrorle}
\quoted{IsPolar}
\quoted{landau}
\quoted{remaindereq}

\head{Notation}
\begingroup\fontsize{8.6}{10.9}\selectfont\raggedright
\setlength{\columnsep}{9mm}\setlength{\parskip}{0pt}
\begin{multicols}{2}
\nt{$\NN,\Z,\Q,\RR,\C$}{natural numbers with $0$, integers, rational, real, complex numbers}
\nt{$\C_\ell,\overline{\Q}_\ell$}{the completion of an algebraic closure of the $\ell$-adic numbers; an algebraic
closure; $|\cdot|_\ell$ their absolute value}
\nt{$\re,\im$}{real and imaginary part}
\nt{$\ee$, $\log$, $(z)_{+}$}{base of the natural logarithm; natural logarithm; $\max(z,0)$}
\nt{$O(\cdot)$, $o(1)$}{Landau's symbols}
\nt{$d$}{a prime number; in Part~II the second parameter}
\nt{$M$}{a natural number, the modulus of the characters}
\nt{$\nu$}{In Part~I also a natural number, the variable of the prime sums, and in Part~II an integer (Convention~0.6)}
\nt{$p$, $e$}{a prime; an integer exponent; $e_\ell$ the exponents of Lemma~\ref{l:rational}}
\nt{$\Lambda$, $\varphi$}{von Mangoldt's function; Euler's function}
\nt{$\chiK{d}$, $\chi_K$}{the quadratic character modulo~$d$ (Definition~\ref{d:chi}); it is the character of
$K=\Q(\sqrt{-d})$ if $d\equiv3\pmod 4$}
\nt{$\xi$, $\chi$, $\one$}{Dirichlet characters modulo~$M$; the character modulo~$1$}
\nt{$\quadTwist{\xi}{d}$}{the twist of $\xi$, a character modulo $Md$ (Definition~\ref{d:chi})}
\nt{$\iv{\chi}$}{the number $1$ if $\chi$ is principal, $0$ otherwise}
\nt{$L(s,\chi)$, $\Ltv{s}{\chi}$, $\entireL{\chi}$}{the Dirichlet $L$-function; the entire $L$-function, its values and
the function (Definition~\ref{d:Lt}); $\tilde L'$ its derivative}
\nt{$\zeta$}{the Riemann zeta function}
\nt{$\ord_s f$}{order of vanishing at $s$ of an entire function $f\ne0$}
\nt{$\Eu$, $\chowla{d}$}{Euler's constant; the Chowla--Selberg constant (Definition~\ref{d:cs})}
\nt{$\lambdaClass{d}{M}{\nu}$}{the weight (Definition~\ref{d:lambda})}
\nt{$\lambdaSum{d}{M}{i}{x}$}{the weighted sums (Definition~\ref{d:A})}
\nt{$\logDerivMean{d}{M}{s}$, $\kap{d}{M}$}{the mean of the logarithmic derivatives and its real part at $s=1$
(Definition~\ref{d:G})}
\nt{$\primeError{d}{M}{x}$}{the remainder (Definition~\ref{d:R})}
\nt{$\Ph{d}{M}{s}$, $\res{d}{M}{s}$}{the transform of the remainder, (5.1); its residue function (Theorem~1)}
\nt{$s$, $s'$, $s_0,s_1,s_2$}{complex numbers and not the exponent of the data (Convention~0.6), the variable of the transforms}
\nt{$\rho$}{a complex number with $\re\rho>0$; in Theorem~2 the point where a zero is excluded}
\nt{$\beta$}{a real number, the exponent of the one-sided bound}
\nt{$c$}{a unit of $\Z/M$, the class of the primes $d$}
\nt{$x$}{a bare $x$ is a free real variable (Convention~0.3): the argument of the remainder and of the weighted sums; $x_0$ (an abscissa of convergence) and $x_\ast$ are particular values}
\nt{$z$, $v$}{real numbers; $v$ the logarithm of the argument of the remainder, $v_\ast$ a particular value}
\nt{$i$}{a natural number: an exponent, the order of a moment or of a derivative; in Part~II also the index of a
row, with $i'$, $i_0$ (the index of a lift is $j$)}
\nt{$a$}{a natural number (Lemma~\ref{l:Rsum}); in the proofs of §\ref{s:rem} it abbreviates $a(D)$}
\nt{$\cE,\cE_0,\cE_1,\cE_2$; $\omega,\omega_1,\omega_2$}{functions $\RR\to\RR$; functions $\C\to\C$, residue functions}
\nt{$\isPolar{\cE}{\omega}$}{$\cE$ is polar with residue function $\omega$ (Definition~\ref{d:polar})}
\nt{$C$, $C_1$, $C_{\mathfrak{r}}$; $\Phi$}{real constants; a meromorphic function on $\C$, the transform; $\Phi^{(i)}$ its
$i$-th derivative}
\nt{$v_0$, $v_1$, $X$}{the parameter of a window, a point of it, a real factor (Theorem~3)}
\nt{$\mathfrak{r}$, $\mathfrak{r}_0$, $\Delta$}{radii of discs: $\mathfrak{r}_0$ the radius in Input~\ref{i:taylor} and the bound of Proposition~\ref{p:moments}, $3<\mathfrak{r}<\mathfrak{r}_0$; a positive real number, in the box of Input~\ref{i:density} and in the proof of that proposition}
\nt{$\varepsilon$}{a real number with $0<\varepsilon<\tfrac1{100}$, in the exponent of the one-sided bound of Chapter~I}
\nt{$I(v)$}{the auxiliary function in the proof of Theorem~3}
\nt{$P=(N_0,d)$, $D$}{parameters and admissible data (Setting~\ref{st:data})}
\nt{$\scx{D}$, $\scxo{D}$, $\scb{D}$}{the scales of $D$}
\nt{$\Dm{D}$, $\Dr{D}$}{its two natural numbers}
\nt{$K$, $E(D)$, $N_1(D)$, $f(D)$}{$\Q(\sqrt{-d})$; the field of the data, the order of its roots of unity, the
number of its infinite places}
\nt{$\tau$, $\nu\cdot\tau$}{an infinite place of $E(D)$; its translate by the integer~$\nu$}
\nt{$\tops{D}$, $\thick{D}$}{the places of class $1$ modulo $N_0$; the thick places}
\nt{$k(D)$, $t(D)$, $a(D)$}{$\Dr{D}+1-\Dm{D}$; $5\Dr{D}$; the least multiple of $N_1(D)$ that is $\ge\scx{D}\,\scb{D}$}
\nt{$\budget{D}$}{$f(D)\,a(D)\,(\log\scb{D})^{3}$}
\nt{$\wD{D}{\tau}$, $\wbar{D}$, $Y$}{the nominal weight, the mean of its logarithm; a real weight}
\nt{$\Ask{D}{z}{\tau}$, $\mainAsk{D}$}{the prime sum of order $z$ at $\tau$; its expected value
(Definition~\ref{d:askmain})}
\nt{$g(D)$, $\mathrm{rd}_{E(D)}$, $\GamD{D}$, $\Omega$, $\pi$}{product of two primes of the data; root discriminant;
the period term $\log(g(D)\,\mathrm{rd}_{E(D)}/2\pi)$; the period; the number $3.14\ldots$}
\nt{$\jet{D}{i}$, $\cost{D}$, $N(D)$, $\hrow{D}$}{jet size, cost of one direction, level, height per row
(Setting~\ref{st:data})}
\nt{$\mainHeight{D}$, $\remainder{D}$}{the main amount and the remainder of the data (Definitions \ref{d:hmain},
\ref{d:rem})}
\nt{$\Col{D}$, $\Colram{D}$, $\integralRamified{D}$}{columns, ramified columns, integral ramified columns}
\nt{$\psi$, $\psi_i$, $\psi'_i$, $\psi_0$}{columns}
\nt{$\GD{D}$, $\gamma$, $\mathcal{T}$}{the finite group of the data, an element, a finite subset}
\nt{$\PF{D}$, $\Row{D}$, $y=(y_i)$, $n$}{forms of principal level, rows, a family of $n$ rows; $n'$, $n''$ further
numbers of rows}
\nt{$\gpt{D}{\gamma}$, $\Poin{D}$, $\operatorname{ram}_D$}{translation followed by the projection to rows; Poincaré
series; the ramifying operator}
\nt{$\scaledInclusion{D}$, $\projectedPoincare{D}{\gamma}{\psi}$}{the scaled inclusion and the Poincaré maps
(Definitions \ref{d:inc}, \ref{d:proj})}
\nt{$\vartheta$, $Q(\vartheta)$, $\theta=(\theta_{\ell,\iota})$}{a twist, its set of ramified primes; a family of
$\Dm{D}$-fold twists}
\nt{$\Thf{D}{\vartheta}$, $\Th{D}{\vartheta}$, $\AAf{D}{\vartheta}$}{the chain of lifts on $\PF{D}$ and on the rows;
its target}
\nt{$\zz{D}{\vartheta}$, $\vol{D}{\vartheta}{j}$}{the volume of the chain; of its $j$-th lift}
\nt{$F$, $F'$, $F_i$}{forms, in $\AAf{D}{\vartheta}$ or in $\PF{D}$}
\nt{$\sigma$, $\AutC$}{an automorphism of the field $\C$; their group}
\nt{$\nsq{F}{\sigma}$, $\covol{F_1,\dots,F_n}{\sigma}$}{Petersson norm and squared covolume of $\sigma$-conjugates}
\nt{$\ell$, $j$, $O_j$}{a prime number (in Lemma~\ref{l:rational} and in sums, a natural number); a ring homomorphism
$\C\to\C_\ell$; its valuation ring in $\C$}
\nt{$\iota$}{a field isomorphism $\C\to\overline{\Q}_\ell$}
\nt{$\eta_\ell$, $J$}{real numbers indexed by $\ell\in\NN$; a finite set of natural numbers}
\nt{$u$, $u_0$}{square matrices}
\nt{$\IndexLe{D}{y}{\delta}$, $\delta$}{$y$ has $D$-index at most $\delta$ (Definition~\ref{d:index}); a real number}
\nt{$V$, $B$, $h$, $h'$}{real numbers: bounds for logarithms of covolumes, of sizes, of heights}
\nt{$q$, $S$}{the rational number and the spanning set of Theorem~A}
\nt{$\Psi$, $\Ch{D}{\vartheta,\Psi}$}{a good datum; the chain of lifts between definite spaces}
\nt{$\gain{D}{\ell}$, $\err{D}{\ell}$, $\nonsplit{D}$}{gain and allowed loss at $\ell$; the primes $\ell<\scxo{D}$ with
$\chiK{d}(\ell)\ne1$}
\nt{$\mathcal{S}$, $\mathcal{M}$, $\mathcal{N}$}{$O_j$-modules of columns and of rows}
\nt{$W$}{a subspace of $\Row{D}$}
\nt{$H$, $H_{\ell,\iota}$}{complex numbers used as factors}
\nt{$Z$, $Z'$, $\alpha$}{families of complex numbers indexed by the pairs $(\ell,\iota)$; a real number in the
definition of the height}
\nt{$\cst{see}{D}$, $\cst{vol}{D}$, $\cst{cr}{D}$}{the constants of Inputs \ref{i:seesaw}, \ref{i:volume},
\ref{i:crude}}
\end{multicols}
\endgroup

\head{References}
\begingroup\fontsize{9.2}{12.2}\selectfont\raggedright
The places are those found by the second reader of this write-up (Inputs \ref{i:landau} and~\ref{i:density}); the
writer has not consulted the texts.

\nt{[Ju]}{M.~Jutila, On mean values of Dirichlet polynomials with real characters, Acta Arithmetica 27 (1975),
191--198. Theorem~2, equation~(6), p.~193.}
\nt{[MV]}{H.~L.~Montgomery and R.~C.~Vaughan, \emph{Multiplicative Number Theory I. Classical Theory}, Cambridge,
2007. §15.1, Lemma~15.1, p.~463.}
\nt{[Wi]}{D.~V.~Widder, \emph{The Laplace Transform}, 1941; the second printing, of 1946, was the one compared.
Chapter~II, §5, Theorem~5b, p.~58.}
\endgroup

\end{document}
